\documentclass{article}
\usepackage{bookmark}
\usepackage{amsmath}
\usepackage{amsthm}
\usepackage{amssymb}
\usepackage{pgfplots}
\usepackage[utf8]{inputenc}
\usepackage{amsfonts}
\usepackage[margin=2.5cm]{geometry}
\usepackage{graphicx}
\usepackage[export]{adjustbox}
\usepackage{fancyhdr}
\usepackage[portuguese]{babel}
\usepackage{hyperref}
\usepackage{multirow}
\usepackage{lastpage}
\usepackage{mathtools}
\usepackage{fouriernc}

\pagestyle{fancy}
\fancyhf{}

\pgfplotsset{compat = 1.18}

\hypersetup{
  colorlinks,
  citecolor=black,
  filecolor=black,
  linkcolor=black,
  urlcolor=black
}
\newtheorem*{def*}{\underline{Defini\c c\~ao}}
\newtheorem*{theorem*}{\underline{Teorema}}
\newtheorem*{lemma*}{\underline{Lema}}
\newtheorem*{prop*}{\underline{Proposi\c c\~ao}}
\newtheorem{example}{\underline{Exemplo}}
\newtheorem*{proof*}{\underline{Prova}}
\newtheorem*{crl*}{\underline{Corolário}}
\renewcommand\qedsymbol{$\blacksquare$}

\rfoot{P\'agina \thepage \hspace{1pt} de \pageref{LastPage}}

\begin{document}
\begin{figure}[ht]
  \minipage{0.76\textwidth}
  \includegraphics[width=4cm]{icmc.png}
  \hspace{7cm}
  \includegraphics[height=4.9cm,width=4cm]{brasao_usp_cor.jpg}
  \endminipage  
\end{figure}

\begin{center}
  \vspace{1cm}
  \LARGE
  UNIVERSIDADE DE S\~AO PAULO

  \vspace{1.3cm}
  \LARGE
  INSTITUTO DE CI\^ENCIAS MATEM\'ATICAS E COMPUTACIONAIS - ICMC

  \vspace{1.7cm}
  \Large
  \textbf{Notas de Álgebra II}

  \vspace{1.3cm}
  \large
  \textbf{Renan Wenzel - 11169472}

  \vspace{1.3cm}
  \large
  \textbf{Professor(a): Behrooz Miraii}

  \textbf{E-mail: behrooz@icmc.usp.br}

  \vspace{1.3cm}
  \today
\end{center}

\newpage
\textbf{{\Huge Disclaimer}}
\vspace{5cm}

{\huge Essas notas não possuem relação com professor algum. 

  Qualquer erro é responsabilidade solene do autor.

Caso julgue necessário, contatar: renan.wenzel.rw@gmail.com}
\tableofcontents

\newpage

\section{Aula 01 - 09/08/2023}
\subsection{Motivações}
\begin{itemize}
  \item Anéis e Anéis Comutativos;
  \item Exemplos de Anéis;
  \item Grupo Multiplicativo e Elemento Inversível.
\end{itemize}
\subsection{Anéis - Nosso Objeto de Estudo no Curso}
\begin{def*}
  Um \textbf{anel} é um conjunto A munido de duas operações, \(+, \cdot \), tais que:
  \begin{itemize}
    \item[1)] \((A, +)\) é um grupo abeliano;
    \item[2)] \((A, \cdot )\) é associativo, com unidade \(1_{A}\);
    \item[3)] Para todos a, b, c em A, temos 
      \[
        (a+b)\cdot c = a \cdot c + b \cdot c \quad\&\quad a \cdot (b+c) = a \cdot b + a \cdot c.\quad\square
      \]
  \end{itemize}
\end{def*}
\begin{def*}
  Dado um anel \((A, +, \cdot )\), diremos que A é um \textbf{anel comutativo} se o grupo \((A, \cdot )\) é abeliano. \(\square\)
\end{def*}
\begin{example}
  \begin{itemize}
    \item[1)] Com as operações usuais, \(\mathbb{Z}, \mathbb{R}, \mathbb{C}\) são todos anéis (comutativos).
    \item[2)] Os anéis de polinômios na variável X com coeficientes inteiros ou complexos são anéis, i.e., 
      \[
        \mathbb{Z}[x]\coloneqq \biggl\{\sum\limits_{i=0}^{n}a_{i}x^{i}: n\in \mathbb{N}, a_{i}\in \mathbb{Z}\biggr\}\quad\&\quad \mathbb{C}[x]\coloneqq \biggl\{\sum\limits_{i=0}^{n}a_{i}x^{i}:n\in \mathbb{N}, a_{i}\in \mathbb{C}\biggr\};
      \]

  \end{itemize}
\end{example}
\textbf{\underline{Observação 1}:} A unidade é única - Se \(1_{A}, e\) são duas unidades de \((A, \cdot ),\) então 
\[
  1_{A}\cdot e = e = e \cdot 1_{A} = 1_{A} \Rightarrow e = 1_{A}.
\]

\textbf{\underline{Observação 2}:} Se A é um anel e \(0\in A\) é o elemento neutro da soma, temos, para todo \(a\in A\), 
\[
  a0 = a(0+0) = a0 + a0 \Rightarrow a0 = 0.
\]
\begin{def*}
  Seja \((A, +, \cdot )\) um anel. Um subconjunto \(B\subseteq{A}\) é um \textbf{subanel} de A se \((B, +, \cdot )\) é um anel e \(1_{A} = 1_{B}.\quad\square\)
\end{def*}
A segunda condição é importante para evitar casos como o seguinte: Considere \(\mathbb{Z}\times \{0\}\subseteq{\mathbb{Z}\times \mathbb{Z}}.\) Note que a restrição
das operações de \(\mathbb{Z}\times \mathbb{Z}\) faria de \(\mathbb{Z}\times \{0\}\) um subanel, mas \(1_{\mathbb{Z}\times \mathbb{Z}} = (1, 1)\) e \(1_{\mathbb{Z}\times\{0\}}=(1, 0).\)
\begin{example}
  \begin{itemize}
    \item[i)] Temos a seguinte cadeia de subanéis: \(\mathbb{Z}\subseteq{\mathbb{Q}}\subseteq{\mathbb{R}}\subseteq{\mathbb{C}}.\)
    \item[ii)] Analogamente, temos a seguinte cadeia de subanéis de polinômios: \(\mathbb{Z}[x]\subseteq{\mathbb{Q}[x]}\subseteq{\mathbb{R}[x]}\subseteq{\mathbb{C}[x]}\).
    \item[iii)]\texttt{(Anel zero)} Seja S um conjunto unitário \(S = \{a\}.\) Podemos definir as operações \(a + a\coloneqq a\) e \(a \cdot a\coloneqq a\). Com estas operações,
      S torna-se um anel, o \textbf{Anel Zero}, no qual \(1_{S} = 0_{S}\). Este anel é especial não só porque podemos brincar que ``1 = 0'', como também porque ele é \textit{único}.

      Com efeito, seja A um anel qualquer no qual \(1_{A} = 0_{A}\). Para \(a\in A\), vale \(a = a \cdot 1_{A} = a \cdot 0_{A} = 0.\) Assim, \(a = 0\) e \(A = \{0\}\) é 
      um anel zero. 

      Portanto, denotaremos o único anel zero por \(\{0_{A}\} = (0).\)
    \item[iv)] \(\mathbb{Z}[\sqrt[]{d}]\) é subanel de \(\mathbb{Q}[\sqrt[]{d}]\);
    \item[v)] Se A é um anel, então 
      \[
        A\subseteq A[x] \subseteq A[x, y] \subseteq A[x, y, z].
      \]
  \end{itemize}
\end{example}
\begin{example}[Exercícios]
  Mostre que os seguintes conjuntos, com suas operações usuais, são anéis:
  \begin{itemize}
    \item[a)] \(\mathbb{Z}[i] \subseteq{\mathbb{C}}, \mathbb{Z}[\sqrt[]{p}]\subseteq{\mathbb{R}},\) em que p é primo;
    \item[b)] \(\mathbb{Z}\biggl[\frac{1}{p}\biggr]\coloneqq \biggl\{\frac{a}{p^{r}}: a\in \mathbb{Z}, r\in \mathbb{N}\biggr\}\subseteq{\mathbb{Q}}\);
    \item[c)] \(\mathbb{Z}_{(p)}\coloneqq \biggl\{\frac{a}{b}: a, b\in \mathbb{Z}\quad \text{e}\quad \text{p não divide b} \biggr\}\subseteq{\mathbb{Q}}\).
  \end{itemize}
\end{example}
\begin{def*}
  Seja A um anel. Um elemento \(a\in A\) é dito \textbf{invertível} (ou \textbf{unitário}) se existe \(b\in B\) tal que \(ab = 1_{A}.\) O \textbf{grupo multiplicativo} de A
  é o conjunto 
  \[
    A^{*}\coloneqq \{a\in A: \text{a é invertível}\}.\quad\square
  \]
\end{def*}
Vale observar que, caso exista o inverso, ele é único.
\begin{prop*}
  O grupo multiplicativo de um anel \((A, +, \cdot )\) é, de fato, um grupo.
\end{prop*}
\begin{proof*}
  Primeiro, note que \(1_{A}\in A^{*},\) pois \(1_{A} \cdot 1_{A} = 1_{A}.\) 

  Além disso, todo elemento \(a\in A^{*}\) possui inverso \(a^{-1}\in A^{*}.\) A operação herda a associatividade do anel. 

  Finalmente, basta provarmos que A é fechado por multiplicação. Com efeito, se \(a, b\in A^{*},\) considere \(x = b^{-1}a^{-1}\) e note que \(abx = abb^{-1}a^{-1} =
  a 1_{A} a^{-1} = aa^{-1} = 1_{A}\) e, assim, \(ab\in A^{*}.\)

  Portanto, \(A^{*}\) é um grupo. \qedsymbol
\end{proof*}
\begin{def*}
  Se A é um anel, dizemos que A é um \textbf{corpo} se todo elemento não nulo de A possui um elemento inverso também em A. Em outras palavras, se \(A = A^{*}.\quad\square\)
\end{def*}
\begin{example}
  Os anéis \(\mathbb{Q}, \mathbb{R}, \mathbb{C}\) são todos corpos. Para \(\mathbb{Q}\), basta notar que, dado \(0\neq \frac{p}{q}\in \mathbb{Q}\), seu inverso
  \(\frac{q}{p} = \biggl(\frac{p}{q}\biggr)^{-1}\) é um elemento de \(\mathbb{Q}\) também. Para \(\mathbb{R},\) dado \(a\in \mathbb{R},\) segue que \(a^{-1} = \frac{1}{a}\in \mathbb{R}\).
  Finalmente, para \(\mathbb{C},\) considere \(z = a + ib\in \mathbb{C}.\) Como 
  \[
    z\overline{z} = a^{2} + b^{2} = \Vert Z \Vert^{2},
  \]
  segue que \(z^{-1}=\frac{\overline{z}}{\Vert z \Vert^{2}}.\)

  Os inteiros de Gauss, \(\mathbb{Z}[i]\), não formam um corpo, mas \(\mathbb{Q}[i]\) forma um corpo. 
\end{example}
\begin{example}[Exercício]
  Seja \(d\in \mathbb{Z}\) livre de quadrados. Mostre que \(\mathbb{Q}[\sqrt[]{d}]\) é um corpo.
\end{example}
\begin{theorem*}
  O conjunto \(\mathbb{Z}_{n} = \{m\mod n: m\in \mathbb{Z}\} = \{\overline{0}, \overline{1}, \dotsc , \overline{n}\}\) é um corpo se, e somente se, n é primo.
\end{theorem*}
Caso p seja primo, denotamos 
\[
  \mathbb{Z}_{p} = \mathbb{F}_{p},
\]
representando o corpo finito de característica (``tamanho'') p.

Seja \(\mathbb{C}_{p}\coloneqq \mathbb{F}_{p}\times \mathbb{F}_{p}\) e defina, neste conjunto, as operações 
\[
  (\overline{a}, \overline{b}) + (\overline{c}, \overline{d}) = (\overline{a}+\overline{c}, \overline{b}+\overline{d})
\]
e 
\[
  (\overline{a}, \overline{b})\cdot (\overline{c}, \overline{d}) = (\overline{a}\overline{c} - \overline{b}\overline{d}, \overline{a}\overline{d}+\overline{b}\overline{c}).
\]
Fica de exercício mostrar que \((\mathbb{C}_{p}, +, \cdot )\) é um anel com unidade multiplicativa \(1_{\mathbb{C}_{p}} = (1, 0).\) e aditiva \(0_{\mathbb{C}_{p}} = (0, 0).\)
Além disso, dado \(a\neq 0\) em \(\mathbb{C}_{p},\) note que \(y = -\frac{b}{a}x\) implica que 
\begin{align*}
  &ax - b(\frac{-b}{a}x) = 1\\
  \Rightarrow & ax + \frac{b^{2}}{a}x = 1.\\
  \Rightarrow & (a^{2} + b^{2})x = a\\
  \Rightarrow & x = \frac{a}{a^{2}+b^{2}} \quad\&\quad y = \frac{-b}{a^{2}+b^{2}}.
\end{align*}
Disto segue que \(a^{2} + b^{2}\neq 0\), ou, equivalentemente, \(1 + \biggl(\frac{b}{a}\biggr)^{2}\neq 0\). Conclui-se que \(x^{2}+1\in \mathbb{F}_{p}[x]\) não tem que ter raíz.
\begin{prop*}
  O conjunto \(\mathbb{C}_{p}\) é um corpo com \(p^{2}\) elementos se, e somente se, o polinômio \(x^{2}-1 \) em \(\mathbb{F}_{p}[x]\) não tem raíz em \(\mathbb{F}_{p}.\)
\end{prop*}
\begin{example}[Exercício]
  Prove que os grupos multiplicativos dos anéis a seguir são estes:
  \begin{itemize}
    \item[1)] \((\mathbb{Z}/n \mathbb{Z})^{*} = \{\overline{r}: (n, r) = 1\}\), em que \((n, r)\) denota o \textit{maior divisor comum} entre n e r;
    \item[2)] \((\mathbb{Z}[i])^{*}=\{\pm1, \pm i\}\);
    \item[3)] \((\mathbb{Z}[\sqrt[]{p}])^{*} = \{\pm1\}\);
    \item[4)] \((\mathbb{Z}[1/p])^{*} = \{\pm p^{r}: r\in \mathbb{Z}\}\);
    \item[5)] \((\mathbb{Z}_{(p)})^{*} = \biggl\{\frac{a}{b}: a, b\in \mathbb{Z}, \text{ em que p não divide a e nem b}\biggr\}\);
    \item[6)] Se A é um anel comutativo, \(M_{n}(A)^{*} = \{A\in M_{n}(A): \det{(A)}\in A^{*}\}\).
  \end{itemize}
\end{example}
\newpage

\section{Aula 02 - 14/08/2023}
\subsection{Motivações}
\begin{itemize}
  \item Domínios;
  \item Ideais, Ideais Principais e Ideais Finitamente Gerados;
  \item Corpos;
  \item Operações com Ideais e Ideais em \(\mathbb{Z}\).
\end{itemize}
\subsection{Domínios e Ideais}
\begin{def*}
  Seja A um anel. Um elemento não nulo \(a\in A\) é dito \textbf{divisor de zero} se existe \(b\in A\setminus{\{0\}}\) tal que
  \(ab = 0.\) Caso A não possua divisores de zero, A é chamado \textbf{domínio.} \(\square\)
\end{def*}
\begin{example}
  O anel \(\mathbb{Z}/4\) não é um domínio, pois \(\overline{2}\cdot \overline{2} = \overline{4} = \overline{0}.\)
\end{example}
\begin{example}
  \begin{itemize}
    \item[1)] Se A é um anel, \(A\times A\) não é um domínio. Basta notar que \((1, 0)\cdot (0, 1) = (0, 0).\)
    \item[2)] O conjunto \(A = \mathbb{Q} + \mathbb{Q}x \subseteq{\mathbb{Q}[x]}\) (apenas como um subconjunto) com operação de soma usual e um
      produto diferente definido por 
      \[
        (a+bx)(c+dx) = ac + (bc+ad)x.
      \]
      Neste caso, \(x \cdot x = 0.\) Essa construção será futuramente conhecida como o anel quociente \(\mathbb{Q}/(x^{2}).\)
    \item[3)] Como \(\mathbb{C}\) é um domínio, qualquer subanel de \(\mathbb{C}\) é um domínio, ou seja, \(\mathbb{Z}, \mathbb{Q}, \mathbb{R}, \mathbb{Z}[i]\) e \(\mathbb{Q}[i]\)
      são todos domínio.
    \item[4)] Se \(d\in \mathbb{Z}\setminus{\{0\}}\) é livre de quadrados, então \(\mathbb{Z}[\sqrt[]{\alpha }]\leq \mathbb{Q}[\sqrt[]{\alpha }] \) são domínio.
    \item[5)] O anel nulo é um domínio.
  \end{itemize}
\end{example}
\begin{example}[Exercício]
  \begin{itemize}
    \item[1)] Temos \(\mathbb{Z}/n \mathbb{Z}\) é um domínio se, e somente se, n é primo;
    \item[2)] A é domínio se, e somente se, \(A[x]\) é domínio. Em particular, \(A[x_{1}, \dotsc , x_{n}]\) é um domínio.
  \end{itemize}
\end{example}
\begin{def*}
  Seja A um anel. Um subconjunto \(\mathfrak{i}\subseteq{A}\) é um \textbf{ideal} se:
  \begin{itemize}
    \item[1)] Para todos \(a, b\in \mathfrak{i}, a + b\in \mathfrak{i}\), ou seja, \textbf{um ideal é fechado pela soma};
    \item[2)] Para \(a\in \mathfrak{i}\) e \(r\in A, r \cdot a\in \mathfrak{i}\), ou seja, \textbf{um ideal absorve produtos}.
  \end{itemize}
  Ideais serão denotados por letras em \textit{Old English Script} - letras frescas, como \(\mathfrak{i}, \mathfrak{j}, \mathfrak{o}, \mathfrak{a}, \dotsc\).
  Além disso, para dizer que um subconjunto é um ideal, escrevemos \(\mathfrak{i} \trianglelefteq{A}.\quad\square\)
\end{def*}
\begin{prop*}
  Todo ideal \(\mathfrak{i}\trianglelefteq{A}\) é um subgrupo pela soma de A, ou seja, \((\mathfrak{i}, +)\leq (A, +).\)
\end{prop*}
\begin{proof*}
  Já vimos que, por definição, os ideais são fechados pela soma. Além disso, como \(0\in A,\) 
  tome \(a\in \mathfrak{i}\) qualquer. Teremos \(0 = 0a\in \mathfrak{i},\) fornecendo o elemento neutro do grupo.
  Finalmente, \(-a = (-1)a\in \mathfrak{i}\). Portanto, \((\mathfrak{i}, +)\) é um subgrupo de \((A, +).\) \qedsymbol
\end{proof*}
\begin{def*}
  Se \((M, +)\) é um grupo abeliano munido de uma aplicação \(A\times \mathfrak{i}\longrightarrow \mathfrak{i}\) com \((r, a)\mapsto r \cdot a\)
  e satisfazendo
  \begin{itemize}
    \item \((r + r')a = ra + r'a\);
    \item \(r(a + a') = ra + ra'\);
    \item \(r(r'a) = (rr')a\);
    \item \(1 \cdot a = a,\)
  \end{itemize}
  dizemos que M é um \textbf{A-módulo.} \(\square\)
\end{def*}
Em particular, se \(\mathfrak{i} \trianglelefteq{A},\) então \(\mathfrak{i}\) é um A-módulo.
\begin{example}
  \begin{itemize}
    \item[i)] Seja A um anel. Então, \(\{0\}\) e A são ideais de A. 
    \item[ii)] Um conjunto \(\mathfrak{i}\) é um ideal de \(\mathbb{Z}\) se, e somente se, existe n inteiro tal que \(\mathfrak{i} = n \mathbb{Z}\coloneqq \{nr:r\in \mathbb{Z}\}\).
      Com efeito, suponha que \(\mathfrak{i}\) é um ideal de \(\mathbb{Z}.\) Caso \(\mathfrak{i} = \{0\},\) então \(\mathfrak{i} = 0 \mathbb{Z},\) ou seja, podemos supor que \(\mathfrak{i}\neq \{0\}.\) Seja 
      \[
        n = \min\{x\in \mathfrak{i}: x > 0\}.
      \]
      Temos \(n \mathbb{Z} = \{nr:r\in \mathbb{Z}\}\subseteq \mathfrak{i}\). Por outro lado, seja \(x\in \mathfrak{i}.\) Pelo algoritmo de divisão, existem
      \(q, rzin \mathbb{Z}\) tais que \(x = nq + r,\) em que \( 0 \leq r \leq n\). Assim, \(r = x-nq = x + (-q)n\in \mathfrak{i}\).
      Porém, \(r\in \mathfrak{i}\) implica que \(0\leq r < n,\) tal que \( r = 0\). Logo,
      \(x = nq \), que pertence a \(n \mathbb{Z}.\) Conclui-se que \(\mathfrak{i}\subseteq n \mathbb{Z}\) e, portanto, \(\mathfrak{i} = n \mathbb{Z}\). O outro lado da 
      afirmação é automático.
    \item[iii)] Se \(a\in A,\) então \(\langle a \rangle\coloneqq aA = \{ra: r\in A\}\trianglelefteq{A}.\) De fato,
      verifiquemos os axiomas de ideal para este caso.

      Caso \(ra, r'a\in \langle a \rangle\), então \(ra + r'a = (r+r')a\in \langle a \rangle\), já que \(r + r'\in A.\)
      Além disso, se \(r'\in A\) e \(ra\in \langle a \rangle\), temos \(r'(ra) = (r'r)a\in \langle a \rangle\), pois \(rr'\in A.\)

      Portanto, \(\langle a \rangle \trianglelefteq{A}\). Em particular, \((0) = \langle 0 \rangle\) e \(\langle 1 \rangle = A.\)
  \end{itemize}
\end{example}
Vale observar que subaneis próprios de um anel, i.e., \(B\subsetneq A\), \textbf{não} são ideais, pois, se fossem então para todo \(a\in A,\) 
\[
  a = a \cdot 1_{A} = a \cdot 1_{B}\in B,
\]
o que resultaria em \(A\subseteq B\), uma contradição.

Em particular, o último tipo do exemplo compõe uma classe importante de ideais, definida a seguir
\begin{def*}
  Seja A um anel. Um ideal \(\mathfrak{i}\trianglelefteq{A}\) é \textbf{principal} se existe \(a\in A\) tal que \(\langle a \rangle = \mathfrak{i}.\quad\square\)
\end{def*}
\begin{lemma*}
  Seja \(\{\mathfrak{i}_{k}\}_{k\in K}\) uma família qualquer de ideais de um anel A. Então, 
  \[
    \bigcap_{k\in K}^{}{\mathfrak{i}_{k}}\trianglelefteq{A}.\quad \text{(Em outras palavras, a interseção de ideais é um ideal)}
  \]
\end{lemma*}
\begin{proof*}
  Para não carregar a notação, coloque \(\mathfrak{i} = \bigcap_{k\in K}^{}{\mathfrak{i_{k}}}.\) Seja \(a, a'\in \mathfrak{i}\) e \(r\in A\).
  Para \(k\in K, a, a'\in \mathfrak{i}_{k}\). Como cada \(\mathfrak{i}_{k}\) é um ideal, \(a + a'\in \mathfrak{i}_{k}\) e, assim,
  \(a + a'\in \mathfrak{i}.\) Além disso, como \(ra\in \mathfrak{i}_{k}\) para todo \(k\in K\), temos também \(ra \in \mathfrak{i}.\)

  Portanto, \(\mathfrak{i}\) é um ideal. \qedsymbol
\end{proof*}
\begin{prop*}[Exercício]
  Seja A um anel. Se \(\mathfrak{i}, \mathfrak{j}\trianglelefteq{A},\) mostre que \(\mathfrak{i}\cup \mathfrak{j}\trianglelefteq{A}\) se, e somente se,
  \(\mathfrak{i}\subseteq{\mathfrak{j}}\) ou \(\mathfrak{j}\subseteq{\mathfrak{i}}.\)
\end{prop*}
Como contra exemplo de que é preciso a hipótese de estarem pelo menos um contido no outro, tome os ideais \(2 \mathbb{Z}\) e \(3 \mathbb{Z}\) dos inteiros.
Então, \(2\mathbb{Z}\cup 3 \mathbb{Z} = \{2r, 3s: r, s \in \mathbb{Z}\},\) tal que é possível ver que esse conjunto não é um ideal de \(\mathbb{Z}\) usando o exemplo
fornecido previamente que dita a forma dos ideais dos inteiros.
\begin{def*}
  Se \(S\subseteq{A}\) é um subconjunto, o ideal gerado por S é o ideal 
  \[
    \langle S \rangle\coloneqq \bigcap_{S\subseteq{\mathfrak{i}}\trianglelefteq{A}}^{}{\mathfrak{i}} \trianglelefteq{A}.\quad\square
  \]
\end{def*}
Esta definição passa a noção de \textit{menor ideal que contém o conjunto S.}
\begin{lemma*}
  Vale a igualdade \(\langle S \rangle = \{\sum\limits_{i=1}^{n}a_{i}s_{i}: a_{i}\in A, s_{i}\in S, n\in \mathbb{N}\}.\)
\end{lemma*}
\begin{proof*}
  Novamente, a fim de não carregar a notação, seja \(\mathfrak{i} = \{\sum\limits_{i=1}^{n}a_{i}s_{i}: a_{i}\in A, s_{i}\in S, n\in \mathbb{N}\}\). Note que
  \(\mathfrak{i}\) é um ideal, pois a soma de dois elementos de \(\mathfrak{i}\) está em \(\mathfrak{i}\) e o produto por um elemento do anel é absorvido.
  A soma segue por definição, e a absorção segue de, se \(r\in A,\) 
  \[
    r \sum\limits_{i=1}^{n}a_{i}s_{i} = \sum\limits_{i=1}^{n}(ra_{i})s_{i}\in \mathfrak{i}.
  \]
  Para ver que \(\langle S \rangle \subseteq{\mathfrak{i}}.\) Para isso, note que \(S\subseteq{\mathfrak{i}},\) bastando tomar \(a_{1} = 1\) para todo
  \(s\in S\). Assim, \(\langle S \rangle\) é a interseção de todos os ideais que contém S e, como \(\mathfrak{i}\) é um deles, a inclusão segue.

  Por outro lado, mostremos que \(\mathfrak{i}\subseteq{\langle S \rangle}.\) Se \(\mathfrak{j}\) é um ideal que contém S, pela definição de ideal, 
  temos \(s_{1}, \dotsc, s_{n}\in S\) implica que \(s_{1} + \dotsc + s_{n}\in \mathfrak{j}\) e, se \(s\in S, r\in A\), então \(rs\in \mathfrak{j}.\)
  Temos, assim, que se \(r_{1}, \dotsc, r_{n}\in A\) e \(s_{1}, \dotsc, s_{n}\in S\), então \(r_{i}s_{i}\in \mathfrak{j}\) e \(\sum\limits_{}^{}r_{i}s_{i}\in \mathfrak{j}.\)
  Com isso, \(\mathfrak{i}\subseteq{\mathfrak{j}}.\) Repetindo isso para todo \(\mathfrak{j}\) que contém S, temos \(\mathfrak{i}\subseteq{\langle S \rangle}.\)

  Portanto, \(\langle S \rangle = \mathfrak{i}.\) \qedsymbol
\end{proof*}
\begin{def*}
  Seja A um anel. Dizemos que \(\mathfrak{i} \trianglelefteq{A}\) é \textbf{finitamente gerado (f.g.)} se existe um conjunto finito
  \(S = \{s_{1}, \dotsc, s_{n}\}\) tal que \(\mathfrak{i} = \langle S \rangle.\quad\square\)
\end{def*}
Em particular, todo ideal principal é finitamente gerado.
\begin{def*}
  Dizemos que um domínio A é um \textbf{domínio de ideais principais (D.I.P.)} se todo ideal é principal. \(\square\)  
\end{def*}
\begin{def*}
  Dizemos que um anel A é um \textbf{corpo} se \((0)\) e A são os únicos ideais principais. \(\square\)
\end{def*}
\begin{example}[Exercício]
  \begin{itemize}
    \item[1)] Se \(\mathfrak{i}\trianglelefteq{A}\) e \(\mathfrak{i}\cap A^{*}\neq\emptyset\), então \(\mathfrak{i} = A\);
    \item[2)] Todo corpo é um domínio;
    \item[3)] Todo domínio finito é um corpo;
    \item[4)] A é um corpo se, e somente se, \(A^{*} = A \setminus{\{0\}}\);
    \item[5)] \(\mathbb{Q}[i]\) é um corpo.
  \end{itemize} 
\end{example}
\begin{lemma*}
  Vale que \(\mathbb{Z}\) e \(\mathbb{Z}/n \mathbb{Z}\) são domínio de ideais principais.
\end{lemma*}
\begin{proof*}
  Seja \(\mathfrak{i} \trianglelefteq{Z}.\) Então, \((\mathfrak{i}, +)\leq (\mathbb{Z}, +)\). Usando o resultado
  de classificação de subgrupos de \((\mathbb{Z}, +)\), que afirma que todo subgrupo é da forma \(n \mathbb{Z}\) para
  algum \(n\in \mathbb{N},\) segue que \(\mathfrak{i} = n \mathbb{Z}\) para algum n e, assim, \(\mathfrak{i} = \langle n \rangle\) 
  é principal.

  O caso \(\mathbb{Z}/n \mathbb{Z}\) fica de exercício. \qedsymbol
\end{proof*}
\begin{def*}
  Seja A um anel. Se \(\mathfrak{i}, \mathfrak{j} \trianglelefteq{A},\) podemos definir os ideais \(\mathfrak{i} + \mathfrak{j}\coloneqq \{x+y:x\in \mathfrak{i}\text{ e } y\in \mathfrak{j}\}\)
  e \(\mathfrak{i} \cdot \mathfrak{j}\coloneqq \{x_{1}y_{1}+\dotsc+x_{n}y_{n}: x_{i}\in \mathfrak{i}, y_{i}\in \mathfrak{j}, \text{ e } n\in \mathbb{N}\}.\quad\square\)
\end{def*}
Vamos mostrar que esses conjuntos são, de fato, ideais. No primeiro caso, sejam \(a, a'\in \mathfrak{i} + \mathfrak{j}.\)
Assim, eles podem ser escritos como \(a = x + y\) e \(a'= x'+ y'\) para \(x, x'\in \mathfrak{i}\) e \(y, y'\in \mathfrak{j}.\)
Então, \(a + a' = (x + x') + (y + y')\in \mathfrak{i} + \mathfrak{j}.\) Considere agora \(r\in A\), tal que 
\[
  ra = r(x'+y') = rx' + ry'.
\]
Como \(\mathfrak{i}, \mathfrak{j}\) são ideais, cada parcela destas está no seu ideal respectivo.
Logo, \(ra \in \mathfrak{i} + \mathfrak{j}.\)

No segundo caso, se \(a, b\in \mathfrak{i}\cdot \mathfrak{j},\) temos 
\[
  a = \sum\limits_{}^{}x_{i}y_{i}\quad\&\quad b = \sum\limits_{}^{}u_{j}v_{j},
\]
com \(x_{i}, u_{j}\in \mathfrak{i}\) e \(y_{i}, v_{j}\in \mathfrak{j}.\) Reorganizando os termos, podemos ver
\(a+b\) como a soma finita de termos da forma \(x_{k}y_{k}, x_{k}\in \mathfrak{i}, y_{k}\in \mathfrak{j}\), ou seja, 
\(a+b\in \mathfrak{i}\cdot \mathfrak{j}\). Além disso, se \(r\in A,\) temos 
\[
  ra = r \sum\limits_{}^{}x_{i}y_{i} = \sum\limits_{}^{}(rx_{i})y_{i}\in \mathfrak{i}\cdot \mathfrak{j},
\]
pois \(rx_{i}\in \mathfrak{i},\) já que \(\mathfrak{i}\) é ideal, provando, assim, que ambos os conjuntos definidos são ideais.
\begin{lemma*}
  Se A é um anel e \(\mathfrak{i}, \mathfrak{j}\trianglelefteq{A},\) então \(\mathfrak{i}\cdot \mathfrak{j}=\langle \{xy: x\in \mathfrak{i}, y\in \mathfrak{j}\} \rangle\) e
  \(\mathfrak{i}+\mathfrak{j} = \langle \mathfrak{i}\cup \mathfrak{j} \rangle\).
\end{lemma*}
\begin{proof*}
  Denotemos por \(\mathfrak{b}\) o conjunto 
  \[
    \mathfrak{b}\coloneqq \langle \{xy:x\in \mathfrak{i}, y\in \mathfrak{j}\} \rangle.
  \]
  Primeiramente, note que \(\mathfrak{i}\cdot \mathfrak{j}\subseteq{\mathfrak{b}}.\) De fato, seja \(\mathfrak{k}\) um ideal tal que \(\{xy:x\in \mathfrak{i}, y\in \mathfrak{j}\}\subseteq{\mathfrak{k}}.\)
  Como \(\mathfrak{k}\) é fechado pela soma, se \(x_{1}, \dotsc, x_{n}\in \mathfrak{i}\) e \(y_{1},\dotsc,y_{n}\in \mathfrak{j},\) temos  
  \(\sum\limits_{}^{}x_{i}y_{i}\in \mathfrak{k}.\) Logo, \(\mathfrak{i}\cdot \mathfrak{j}\subseteq{\mathfrak{k}}\) e \(\mathfrak{i}\cdot \mathfrak{j}\) está
  na intersecção \(\langle \{xy: x\in \mathfrak{i}, y\in \mathfrak{j}\} \rangle.\)

  Por outro lado, por definição, elementos da forma \(xy\) tal que \(x\in \mathfrak{i}, y\in \mathfrak{j}\) são elementos
  de \(\mathfrak{i}\cdot \mathfrak{j}.\) Assim, \(\{xy:x\in \mathfrak{i}, y\in \mathfrak{j}\}\subseteq{\mathfrak{i}\cdot \mathfrak{j}}.\) Como
  \(\mathfrak{i}\cdot \mathfrak{j}\) é ideal, a intersecção de ideais com ele está contido nele, ou seja, \(\mathfrak{b}\subseteq{\mathfrak{i}\cdot \mathfrak{j}}\).
  Logo, \(\mathfrak{b} = \mathfrak{i}\cdot \mathfrak{j}.\)

  Com relação à segunda parte, começamos por mostrar que \(\mathfrak{i}+\mathfrak{j}\subseteq{\langle \mathfrak{i}\cup \mathfrak{j} \rangle}.\) De fato,
  seja \(\mathfrak{k}\) um ideal tal que \(\mathfrak{i}\cup \mathfrak{j}\subseteq{\mathfrak{k}}.\) Como \(\mathfrak{k}\) é fechado pela soma,
  se \(x\in \mathfrak{i}\) e \(y\in \mathfrak{j},\) temos \(x+y\in \mathfrak{k}\) tal que, assim, \(\mathfrak{i} + \mathfrak{j}\subseteq{\mathfrak{k}}\).
  Finalmente, como \(\mathfrak{k}\) é qualquer, \(\mathfrak{i}+\mathfrak{j}\subseteq{\langle \mathfrak{i}\cup \mathfrak{j} \rangle}.\)

  Em contrapartida, pela definição de \(\mathfrak{i} + \mathfrak{j},\) conseguimos 
  que \(\mathfrak{i}\cup \mathfrak{j}\subseteq{\mathfrak{i}+\mathfrak{j}}\) (basta tomar 0 + y ou x + 0).
  Como \(\mathfrak{i} + \mathfrak{j}\) é ideal e contém \(\mathfrak{i}\cup \mathfrak{j},\) segue que \(\langle \mathfrak{i}\cup \mathfrak{j} \rangle\subseteq{\mathfrak{i}+\mathfrak{j}}\).
  Portanto, \(\mathfrak{i}+\mathfrak{j} = \langle \mathfrak{i}\cup \mathfrak{j} \rangle.\) \qedsymbol
\end{proof*}
\begin{prop*}
  Seja \(\{\mathfrak{i}_{k}\}_{k\in K}\) uma família de ideais de A. Então,
  \begin{itemize}
    \item[1)] O conjunto 
      \[
        \sum\limits_{k\in K}^{}\mathfrak{i}_{k}\coloneqq \biggl\{\sum\limits_{\text{finita}}^{}a_{i}: a_{i}\in \mathfrak{i}\biggr\},
      \]
      é um ideal e, além disso, 
      \[
        \sum\limits_{k\in K}^{}\mathfrak{i}_{k}\coloneqq \biggl\langle \bigcup_{k\in K}^{}{\mathfrak{i}_{k}} \biggr\rangle
      \]
    \item[2)] Se \((K,\leq )\) é um conjunto totalmente ordenado e a família satisfaz a propriedade que
      \textbf{para todos k, k' em K tais que} \(k\leq k'\), \textbf{temos} \(\mathfrak{i}_{k}\subseteq{\mathfrak{i}_{k'}}\). Então, 
      \[
        \bigcup_{k\in K}^{}{\mathfrak{i}_{k}}\trianglelefteq{A}.
      \]
  \end{itemize}
  \begin{proof*}
    Provar que \(\sum\limits_{}^{}\mathfrak{i}_{k}\) é um ideal e que é igual ao gerado pela união é análogo ao caso finito.

    Para o item dois, sejam \(a, a'\in \cup \mathfrak{i}_{k}.\) Assim, existem \(k, k'\) tais que \(a\in \mathfrak{i}_{k}\) e \(a'\in \mathfrak{i}_{k'}\).
    Como K é totalmente ordenado, podemos supor, sem perda de generalidade, que \(k\leq k'.\) Pela propriedade que K possui,
    \(\mathfrak{i}_{k}\subseteq{\mathfrak{i}_{k'}}\) e \(a, a'\in \mathfrak{i}_{k'}.\) Como \(\mathfrak{i}_{k'}\) é ideal, a soma
    \(a + a'\in \mathfrak{i}_{k}\) e, assim, \(a + a'\in\cup \mathfrak{i}_{k}.\) Agora, considere \(r\in A\)
    e \(a\in \cup \mathfrak{i}_{k}.\) Então, existe \(k\in K\) tal que \(a\in \mathfrak{i}_{k}\). Já que \(\mathfrak{i}_{k}\)
    é ideal, \(ra\in \mathfrak{i}_{k}\), tal que \(ra\in \cup \mathfrak{i}_{k}.\) Portanto, \(\bigcup_{k\in K}^{}{\mathfrak{i}_{k}}\) é um
    ideal de A. \qedsymbol
  \end{proof*}
\end{prop*}
\begin{example}[Exercício]
  Se \(n \mathbb{Z}\) e \(m \mathbb{Z}\) são ideais de \(\mathbb{Z},\) prove que:
  \begin{itemize}
    \item[1)] \(n \mathbb{Z}\cdot m \mathbb{Z} = nm \mathbb{Z};\)
    \item[2)] \(n \mathbb{Z} + m \mathbb{Z} = \mathrm{mdc}(n, m) \mathbb{Z};\)
    \item[3)] \(n \mathbb{Z}\cap m \mathbb{Z} = \mathrm{mmc}(n, m) \mathbb{Z}.\)
  \end{itemize}
\end{example}
\begin{def*}
  Seja A um anel e \(\mathfrak{i}, \mathfrak{j}\trianglelefteq{A}.\) Dizemos que \(\mathfrak{i}, \mathfrak{j}\) são \textbf{coprimos}
  se \(\mathfrak{i} + \mathfrak{j} = A.\quad\square\)
\end{def*}
\begin{prop*}[Exercício]
  Sejam \(\mathfrak{i}, \mathfrak{j}\trianglelefteq{A}\). Se \(\mathfrak{i}, \mathfrak{j}\) são coprimos, então \(\mathfrak{i}\cdot \mathfrak{j} = \mathfrak{i}\cap \mathfrak{j}.\)
\end{prop*}
\newpage

\section{Aula 03 - 16/08/2023}
\subsection{Motivações}
\begin{itemize}
  \item Ideais Primos e Ideais Maximais;
  \item Espectro de um Anel;
  \item Anéis Locais e Elementos Nilpotentes.
\end{itemize}
\subsection{Ideais Primos e  Maximais}
\begin{def*}
  Uma ideal \(\mathfrak{i} \trianglelefteq{A}\) é dito \textbf{primo} se \(\mathfrak{i}\neq A\) e \(ab\in \mathfrak{i}\)
  implica \(a\in \mathfrak{i}\) ou \(b\in \mathfrak{i}.\) 

  Um ideal \(\mathfrak{i}\trianglelefteq{A}\) é dito maximal se \(\mathfrak{i}\neq A\) e se, dado \(\mathfrak{j}\trianglelefteq{A}\)
  com \(\mathfrak{i}\subseteq{\mathfrak{j}}\), então 
  \[
    \mathfrak{i} = \mathfrak{j}\text{ ou } \mathfrak{j} = A.\quad\square
  \]
\end{def*}
Diremos que uma \textbf{cadeia} \(K\subseteq{S}\) de um conjunto parcialmente ordenado é um subconjunto totalmente
ordenado do mesmo.
\begin{prop*}
  \begin{itemize}
    \item[1)] Os ideais maximais são primos;
    \item[2)] Todo anel não nulo tem um ideal maximal e logo tem um ideal primo;
    \item[3)] Todo ideal próprio de A está contido em algum ideal maximal de A.
  \end{itemize}
\end{prop*}
\begin{proof*}
  1. Seja M maximal e \(ab\in M.\) Mostremos \(a\in M\) ou \(b\in M\). Suponha
  \(a\not\in M\). Considere o anel:
  \[
    M + \langle a \rangle = \{x+ra: x\in M\text{ e }r\in A\} \trianglelefteq{A}.
  \] 
  Note que \(M\subsetneq{M + \langle a \rangle},\) pois \(a\not\in M.\) Como M é maximal, segue que
  \(M + \langle a \rangle = A.\) Assim, \(1\in M + \langle a \rangle\) e existe \(x\in M\), \(r\in A\)
  tais que \(x + ra = 1\) e, então, 
  \[
    b = b1 = b(x+ra) = bx + rab.
  \]
  Como M é ideal, \(ab\in M\), o que implica em \(rab\in M\), e \(x\in M\) implica em \(bx\in M\).
  Assim, \(b = bx + rab\in M\). 

  2. Para obtermos o ideal maximal, usaremos o Lema de Zorn. Construiremos o conjunto parcialmente ordenado 
  dado por \(S\coloneqq \{\mathfrak{i}\trianglelefteq{A}: \mathfrak{i}\neq A\}\), sendo a ordem dada por \(\subseteq{}.\)
  Sabemos que \(S \neq\emptyset\), pois \((0)\in S.\) Considere, agora, uma cadeia \(\{\mathfrak{i}_{l}\}_{l\in L}\) de elementos
  de S. Por indução, a união destes ideais é um ideal e, além disso, é um ideal próprio
  \[
    \mathfrak{i}\coloneqq \bigcup_{l\in L}^{}{\mathfrak{i}_{l}}\vartriangleleft{A},
  \]
  pois para cada \(l\in L, 1\not\in \mathfrak{i}_{l}\), já que, se \(1\in \mathfrak{i}_{l}\) para algum \(l\in L\), teríamos
  \(\mathfrak{i} = A\). Então, a união é um elemento de S e é uma cota superior da cadeia, no sentido de \(\mathfrak{i}_{l}\subseteq{\mathfrak{i}}\)
  para cada \(l\in L.\) Pelo Lema de Zorn, segue que existe um elemento maximal no conjunto S, que será o ideal maximal por definição.

  3. Considere \(\mathfrak{i}\) um ideal próprio de A e defina o conjunto \(S_{\mathfrak{i}}\coloneqq \{\mathfrak{j}\vartriangleleft{A}:\mathfrak{i}\subseteq{\mathfrak{j}}\}\).
  Como \(\mathfrak{i}\in S_{\mathfrak{i}}, S_{\mathfrak{i}}\neq\emptyset.\) Seja \(\{\mathfrak{j}_{l}\}_{l\in L}\) uma cadeia de \(S_{\mathfrak{i}}.\)
  Assim como no item 2, a união de elementos da cadeia é cota superior para ela e, aplicando o Lema de Zorn, temos um elemento
  maximal dessa família - um anel maximal que contém \(\mathfrak{i}.\) 
  \qedsymbol
\end{proof*}
\begin{def*}
  Seja A um anel. Definimos o \textbf{espectro primo} de A por \(\mathrm{Spec}(A) = \{\mathfrak{p}\trianglelefteq{A}: \mathfrak{p}\text{ é ideal primo de } A\}\)
  e o \textbf{espectro maximal} de A como \(\mathrm{Specm}(A) = \{\mathfrak{m}\trianglelefteq{A}:\mathfrak{m} \text{ é ideal maximal de } A\}.\quad\square\)
\end{def*}
\begin{prop*}[Exercício]
  \begin{itemize}
    \item[1)] A é domínio se, e somente se, \((0)\in \mathrm{Spec}(A)\)
    \item[2)] Dado um ideal \(\mathfrak{i} = n \mathbb{Z},\) então \(\mathfrak{i}\in \mathrm{Spec}(\mathbb{Z})\) se, e somente se,
      n é primo ou n = 0. Assim, \(\mathrm{Spec}(\mathbb{Z}) = \{p \mathbb{Z}: p \text{ é primo}\}\cup \{(0)\}\).
    \item[3)] Dado um ideal \(\mathfrak{i} = n \mathbb{Z}\), então \(\mathfrak{i}\in \mathrm{Specm}(\mathbb{Z})\) se, e somente se,
      n é primo. Logo, \(\mathrm{Specm}(\mathbb{Z}) = \{p \mathbb{Z}:p \text{ é primo}\}\).
    \item[4)] Mostre que \(\mathfrak{i}\trianglelefteq A\) é primo se, e somente se, para \(\mathfrak{j}, \mathfrak{k}\trianglelefteq A\) são tais que \(\mathfrak{j}\mathfrak{k}\subseteq \mathfrak{i}\), então 
      \(\mathfrak{j}\subseteq \mathfrak{i}\) ou \(\mathfrak{k}\subseteq \mathfrak{i}\).
  \end{itemize}
\end{prop*}
\begin{def*}
  Um anel é dito \textbf{local} se tem só um ideal maximal. \(\square\)
\end{def*}
\begin{lemma*}
  Seja \(\mathfrak{m}\in \mathrm{Specm}(A)\). Então, A é local se, e somente se, \(A^{*} = A\setminus{\mathfrak{m}}\).
\end{lemma*}
\begin{proof*}
  \(\Rightarrow )\) Suponha que \(\mathfrak{m}\) seja o único ideal maximal de A. Como \(\mathfrak{m} \neq\emptyset, A^{*}\cap \mathfrak{m} = \emptyset\).
  Assim, \(A^{*}\subseteq{A\setminus{\mathfrak{m}}}.\) Tome \(a\in A\setminus{\mathfrak{m}}\) e considere o ideal principal \(\langle a \rangle \trianglelefteq{A}.\)
  Se \(\langle a \rangle\neq A,\) então \(\langle a \rangle \subseteq{\mathfrak{m}},\) pois todo ideal próprio está contido num maximal, que supomos ser unicamente \(\mathfrak{m}\).
  Logo, \(a\in \mathfrak{m}\), o que é um absurdo. Então, \(\langle a \rangle = A\) e existe \(r\in A\) tal que \(ra = 1,\) i.e., \(a\in A^{*}.\)

  \(\Leftarrow )\) Suponha que existe outro ideal maximal \(\mathfrak{m}'\neq \mathfrak{m}.\) Então, existe \(a\in \mathfrak{m}'\)
  tal que \(a\in A\setminus{\mathfrak{m}} = A^{*}.\) Portanto, \(\mathfrak{m}'\cap A^{*} \neq\emptyset\) e, assim, \(\mathfrak{m}' = A,\)
  contrariando a maximalidade de \(\mathfrak{m}.\) \qedsymbol
\end{proof*}
\begin{theorem*}
  Seja A um anel não-nulo. Então, todo ideal próprio de A está contido num ideal maximal. Em particular, todo anel
  não nulo tem pelo menos um ideal maximal.
\end{theorem*}
\begin{proof*}
  Seja \(\mathfrak{i}\trianglelefteq A\) e defina a coleção 
  \[
    \mathcal{A} \coloneqq \{\mathfrak{j}\trianglelefteq A: \mathfrak{j} \text{ é um ideal próprio de A que contém }\mathfrak{i}\}.
  \]
  Por definição, \(\mathfrak{i}\in \mathcal{A},\) tal que \(\mathcal{A}\neq\emptyset.\) Seja, agora, \(\{\mathfrak{j}_{l}\}_{l\in L}\) uma cadeia
  dos ideais de A que estão em \(\mathcal{A}\) de modo tal que, para todo \(l, l'\in L\),
  \(\mathfrak{j}_{l}\subseteq \mathfrak{j}_{l'}\) ou \(\mathfrak{j}_{l'}\subseteq \mathfrak{j}\). Temos 
  \[
    \mathfrak{m}\coloneqq \bigcup_{l\in L}^{}\mathfrak{j}_{l}\trianglelefteq A
  \]
  e 
  \[
    \mathfrak{i}\subseteq \bigcup_{l\in L}^{}\mathfrak{j}_{l}\neq A.
  \]
  De fato, dados \(a, b\in \mathfrak{m}\), existem \(l, l'\in L\) tais que \(a\in \mathfrak{j}_{l}, b\in \mathfrak{j}_{l'}\).
  Caso \(\mathfrak{j}_{l}\subseteq \mathfrak{j}_{l'},\) então ambos \(a, b\) pertencem a \(\mathfrak{j}_{l'},\) ou seja, \(a + b, ra\in \mathfrak{j}_{l'}.\) Caso
  seja o oposto e \(\mathfrak{j}_{l'}\subseteq \mathfrak{j}_{l}\), o raciocínio aplica-se para \(a + b, ra\in \mathfrak{j}_{l}\). Logo, \(\mathfrak{m}\trianglelefteq A.\)

  Com relação à segunda propriedade, é claro que \(\mathfrak{i}\subseteq \bigcup_{l\in L}^{}\mathfrak{j}_{l}.\) Assim, se
  \(\mathfrak{m} = A,\) então \(1\in \mathfrak{m}.\) Logo, existe \(l\in L\) tal que \(1\in \mathfrak{j}_{l'},\) donde segue que 
  \(\mathfrak{j}_{l} = A.\) Contradição.

  Em outras palavras, temos um candidato para ideal maximal. Agora, utilizando o \textbf{Lema de Zorn}\footnotemark{1} \footnotetext{Relembre o enunciado do \textbf{Lema de Kuratowski-Zorn: Suponha que um conjunto parcialmente ordenado P tem a propriedade de que toda
  cadeia em P tem um limitante superior. Então, o conjunto P contém ao menos um elemento maximal.} Ele é equivalente ao axioma da escolha e ao princípio da boa-ordenação.}, \(\mathcal{A}\) tem
  ao menos um elemento maximal - \(\mathfrak{m}\). Para ver que este elemento maximal é, de fato, um ideal maximal, suponha que ele não é. Então, existe \(\mathfrak{j}_{i} \trianglelefteq A\) tal que 
  \(\mathfrak{m}\subsetneq \mathfrak{j}_{i} \subsetneq A\). Como \(\mathfrak{i}\subseteq \mathfrak{j}_{i}\neq A,\) segue que \(\mathfrak{j}_{i}\in \mathcal{A}.\) Mas isto é uma contradição,
  pois \(\mathfrak{m}\) é o elemento maximal de \(\mathcal{A}.\) Portanto, \(\mathfrak{m}\) é um ideal maximal de A. \qedsymbol
\end{proof*}
\begin{example}
  Segue que \(n \mathbb{Z}\) é primo se, e somente se, n é primo ou n = 0 (isso está em um dos exercício).
  De fato, suponha que p é primo e tome \(ab\in p \mathbb{Z}.\) Então, p divide ab e, como p é primo, p divide a ou p divide b.
  Isto significa exatamente que \(a\in p\mathbb{Z}\) ou \(b\in p \mathbb{Z}.\) Caso p = 0, se \(ab\in (0),\) então
  \(ab=0\). Sendo \(\mathbb{Z}\) um domínio, isto significa que ou a = 0, ou b = 0, tal que \(a\in (0)\) ou \(b\in (0)\).

  Por outro lado, seja \(n \mathbb{Z}\) primo. Caso  n = 0, não há nada para provar, então suponha que \(n\neq 0\).
  Neste caso, suponha que n não é primo, ou seja, existem \(l, k\) diferentes de 1 e de n tais que \(n = lk\). Note que 
  \(n = lk\in n\mathbb{Z}\), mas \(l\not\in n \mathbb{Z}\) e \(k\not\in n \mathbb{Z}.\) Contradição. Portanto, n tem que ser primo.
\end{example}
\begin{example}
  Vamos olhar o caso de \(\mathbb{Z}_{(n)} = \mathbb{Z}/n \mathbb{Z}.\) Se n for primo, então \(\mathbb{Z}_{(n)}\) é um corpo, ou seja, tem apenas \((\overline{0})\) como 
  ideal maximal. Por outro lado, se n não é primo, existem l, k diferentes de 1 e de n tais que \(n = lk\). Em particular, isto significa que 
  \[
    \overline{0} = \overline{n} = \overline{lk} = \overline{l}\overline{k}\in (\overline{0}),
  \]
  mas tanto \(\overline{l}\) quanto \(\overline{k}\) são diferentes de \(\overline{0}\) por serem diferentes de n. Logo, conclui-se que \((\overline{0})\) não é
  um ideal primo neste caso. Qual seria, então, a cara dos ideais primos de \(\mathbb{Z}_{n}\) para n não sendo primo?

  Para responder isso, seja p um primo tal que \(p\mid n\). Como \(p\leq n, \langle \overline{p} \rangle\subseteq n \mathbb{Z}\) e p é primo, segue que 
  \(\langle \overline{p} \rangle\) é um ideal primo. Além disso, suponha que n é tal que \(n = pl\). Isso nos permite explicitar a cara do ideal gerado pela classe
  de equivalência de p:
  \[
    \langle \overline{p} \rangle = \{\overline{0}, \overline{p}, 2\overline{p}, \dotsc , (l-1)\overline{p}\}.
  \]
\end{example} 
\begin{example}[Exercício]
  \begin{itemize}
    \item[1)] Mostre que todo corpo é um anel local.
    \item[2)] Mostre que \(\mathbb{Z}_{(p)},\) para p primo, é um anel local.
    \item[3)] Seja \(\mathbb{K}\) um corpo e seja \(f(x)\in \mathbb{K}[x]\) um polinômio não nulo irredutível.
      Defina o seguinte conjunto:
      \[
        \mathbb{K}[x]_{(f(x))} = \biggl\{\frac{h(x)}{g(x)}: h(x), g(x)\in \mathbb{K}[x] \text{ e f(x) não divide g(x)}\biggr\}.
      \]
      Mostre que este é um anel local de \(\mathbb{K}[x].\)
  \end{itemize}
\end{example}
\begin{example}
  Sejam F um corpo e seja \(A = F[x]\). Vamos mostrar que todo ideal de A é formado por um polinômio em x.
  Começando pelo caso trivial, se \(\mathfrak{i} = (0),\) então \(\mathfrak{i} = \langle 0 \rangle.\) Assim, suponha que \(\mathfrak{i}\neq (0)\)
  e seja \(f(x)\in \mathfrak{i}\) um polinômio com menor grau em \(\mathfrak{i}.\) Por um lado, é claro que \(\langle f((x) \rangle\subseteq \mathfrak{i},\)
  afinal \(f(x)\) é um elemento de \(\mathfrak{i}.\) Por outro lado, seja \(h(x)\in \mathfrak{i},\) de modo que, pelo algoritmo da divisão, existem \(q(x), r(x)\in F[x]\) tais que 
  \[
    h(x) = q(x)f(x) + r(x),\quad 0\leq \deg{(r)}<\deg{(f)}.
  \]
  Note que, em particular, \(r(x) = h(x)-q(x)f(x)\in \mathfrak{i},\) tal que \(r(x)\in \mathfrak{i}\) e \(\deg{(r)} <\deg{(f)}\) 
  imlicam que \(r(x) = 0\). Com isso, \(h(x) = q(x)f(x),\) ou seja, \(h(x)\in \langle f(x) \rangle,\) mostrando que \(\mathfrak{i}\subseteq \langle f(x) \rangle.\)
  Portanto, \(\mathfrak{i} = \langle f(x) \rangle.\)
\end{example}
\begin{lemma*}[Exercício]
  Seja \(\mathfrak{i}\in F[x]\), em que F é um corpo. Então, 
  \begin{itemize}
    \item[1)] \(\mathfrak{i}\) é primo se \(\mathfrak{i} = 0\) ou \(\mathfrak{i} = \langle f(x) \rangle,\) com f(x) irredutível;
    \item[2)] \(\mathfrak{i}\) é maximal se \(\mathfrak{i} = \langle f(x) \rangle,\) f(x) irredutível
  \end{itemize}
\end{lemma*}
\begin{def*}
  Para um anel A, os ideais a seguir:
  \[
    \mathrm{nil}(A) = \bigcap_{\mathfrak{p}\in \mathrm{Spec}(A)}^{}{\mathfrak{p}}\quad \& \quad J(A)\coloneqq \bigcap_{\mathfrak{m}\in \mathrm{Specm}(A)}^{}{\mathfrak{m}}
  \]
  são chamados nil-radical e radical Jacobson de A. \(\square\)
\end{def*}
\begin{prop*}
  Temos \(\mathrm{nil}(A) \subseteq{J(A)}.\)
\end{prop*}
\begin{proof*}
  De fato, como sabemos que os ideais maximais são primos, vale que \(\mathrm{Specm}(A)\subseteq{\mathrm{Spec}(A)}.\) Assim, se
  \(x\in \mathrm{nil}(A),x\in \mathfrak{p}\) para todo \(\mathfrak{p}\in \mathrm{Spec}(A).\) Em particular, \(x\in \mathfrak{p}\) para todo \(\mathfrak{p}\in \mathrm{Specm}(A)\)
  e então \(x\in J(A)\). Portanto, \(\mathrm{nil}(A) \subseteq{J(A)}.\) \qedsymbol
\end{proof*}
\begin{def*}
  Um elemento \(a\in A\) é dito \textbf{nilpotente} se existir um inteiro \(n\geq 1\) tal que \(a^{n} = 0.\) Um ideal
  \(\mathfrak{i}\) é dito \textbf{nilpotente} se existir um inteiro \(n\geq 1\) tal que \(\mathfrak{i}^{n} = (0).\)
\end{def*}
\begin{prop*}[Exercício]
  \begin{itemize}
    \item[1)] Mostre que \(\mathrm{nil}(A)\) é o conjunto de todos os elementos nilpotentes de A;
    \item[2)] Mostre que se \(\mathrm{nil}(A)\) é finitamente gerado, ele é um ideal nilpotente;
    \item[3)] Se A é domínio, \((0)\in \mathrm{Spec}(A)\) e, logo, \(\mathrm{nil}(A) = (0);\)
    \item[4)] Se A é local com ideal maximal \(\mathfrak{m}\), então \(J(A) = \mathfrak{m}.\)
  \end{itemize}  
\end{prop*}
\newpage

\section{Aula 04 - 21/08/2023}
\subsection{Motivações}
\begin{itemize}
  \item Morfismos entre Anéis;
  \item Núcleo e Imagem;
  \item Morfismos que Preservam Estruturas.
\end{itemize}
\subsection{Morfismos de Anéis}
\begin{def*}
  Sejam \((A, +_{A}, \cdot_{A})\) e \(B, +_{B}, \cdot_{B})\) anéis. Um mapa \(f:A\rightarrow B\) é chamado um \textbf{homomorfismo de anéis}
  (ou, neste texto, apenas morfismo), se:
  \begin{itemize}
    \item[i)] \(f(a +_{A} a') = f(a) +_{B} f(a')\) para todo \(a, a'\in A;\)
    \item[ii)] \(f(a \cdot_{A} a') = f(a)\cdot_{B}f(a')\) para todo \(a, a'\in A;\)
    \item[iii)] \(f(1_{A}) = 1_{B}.\quad\square\)
  \end{itemize}
\end{def*}
\begin{example}
  A função \(f:\mathbb{Z}\rightarrow \mathbb{Z}_{n}\) dada por \(f(r) = \overline{r}\) é um homomorfismo de anéis. No entanto, não existe homomorfismo
  \(g:\mathbb{Z}_{n}\rightarrow \mathbb{Z},\) pois, caso existisse, então 
  \[
    g(\overline{1}) = 1 \Rightarrow g(n\overline{1}) = g(\overline{0}) = 0,
  \]
  mas \(g(n\overline{1}) = \underbrace{g(\overline{1}) + \dotsc + g(\overline{1})}_{\text{n-vezes}} = 1 + \dotsc + 1 = n\). Contradição, pois isso
  significaria que 0 = n se g fosse um homomorfismo.
\end{example}
\begin{lemma*}
  Um homomorfismo possui as seguintes propriedades:
  \begin{itemize}
    \item[1)] \(f(0_{A}) = 0_{B}\)
    \item[2)] Se \(a\in A\) é inversível, então \(f(a)\) é inversível e \(f(a)^{-1} = f(a^{-1})\)
    \item[3)] \(f(-1_{A}) = -1_{B}\)
    \item[4)] \(f(n1_{A}) = n1_{B}\)
    \item[5)] \(f(a_{1} + \dotsc + a_{n}) = f(a_{1})+\dotsc +f(a_{n})\) e \(f(a_{1}\cdot \dotsc \cdot a_{n}) = f(a_{1})\dotsc f(a_{n}).\)
  \end{itemize}
\end{lemma*}
\begin{example}
  Para todo anel A, \(i_{A}:A\rightarrow A,\) dada por \(i_{A}(a) = a\), é um homomorfismo.
\end{example}
\begin{def*}
  Dizemos que um homomorfismo de anéis \(f:A\rightarrow B\) é:
  \begin{itemize}
    \item[1)] \textbf{monomorfismo} se f é um morfismo injetor;
    \item[2)] \textbf{epimorfismo} se f é um morfismo sobrejetor;
    \item[3)] \textbf{isomorfismo} se f é um morfismo bijetor.
  \end{itemize}
  Caso exista um isomorfismo entre dois anéis, dizemos que eles são \textbf{isomorfos}. \(\square\)
\end{def*}
\begin{example}
  A função de antes, \(f:\mathbb{Z}\rightarrow \mathbb{Z}_{n}, r \mapsto \overline{r},\) é um epimorfismo, mas não é mono (note que 0, n e qualquer múltiplo de n
  são mapeados para o mesmo elemento).
\end{example}
\begin{prop*}[Exercício]
  \begin{itemize}
    \item[1)] A identidade \(id_{A}:A\rightarrow A\) é um isomorfismo;
    \item[2)] Se A é subanel de B, então a inclusão \(i:A\rightarrow B\), \(i(a) = a\), é
      um monomorfismo.
  \end{itemize}
\end{prop*}
\begin{example}[Exercício]
  \begin{itemize}
    \item[1)] Se A é um anel qualquer, então:
      \begin{align*}
        \varphi:&\mathbb{Z}\rightarrow A\\
                &n\mapsto n1_{A},
      \end{align*}
      é um morfismo.
    \item[2)] Se \(B=(0), f:A\rightarrow (0)\) é um morfismo.
    \item[3)] Se A é um anel, então:
      \begin{align*}
        \varphi_{a}:&A[x]\rightarrow A\\
                    &f(x)\mapsto f(a).
      \end{align*}
  \end{itemize}
\end{example}
\begin{def*}
  Seja \(f:A\rightarrow B\) um homomorfismo. O conjunto \(\ker{(f)}\coloneqq \{a\in A: f(a) = 0\}\)
  é chamado \textbf{núcleo} ou \textbf{kernel} de f. O conjunto \(\mathrm{Im}(f)\coloneqq f(A)\) é chamado 
  \textbf{imagem} de f. \(\square\)
\end{def*}
\begin{lemma*}
  O kernel de um homomorfismo é um ideal do Anel domínio e a imagem de um homomorfismo é subanel do codomínio. Em outras palavras,
  \[
    \ker{(f)}\trianglelefteq{A}\quad\&\quad \mathrm{Im}(f)\text{ é subanel de }B.
  \]
\end{lemma*}
\begin{proof*}
  1.) Da Álgebra I, sabemos que \((\ker{(f)}, +)\leq (A, +).\) Seja \(a\in A\) e 
  considere \(x\in\ker{(f)}.\) Note que \(f(ax) = f(a)f(x) = f(a)0 = 0\) e, então, \(ax\in\ker{(f)}.\)

  2.) Para provar este, basta usarmos as propriedades de morfismo de anéis. Segue que
  \(f(a), f(b)\in \mathrm{Im}(f),\quad f(a) + f(b) = f(a+b)\in \mathrm{Im}(f).\) Além disso,
  \(f(a)f(b) = f(ab)\in \mathrm{Im}(f)\) e \(1_{A}f(a) = f(1)f(a) = f(1a) = f(a).\) \qedsymbol
\end{proof*}
\begin{example}
  Seja A um anel e considere \(A[x].\) Dado \(a\in A\), defina o mapa
  \begin{align*}
    \varphi_{a}:&A[x]\rightarrow A\\
                &f(x)\mapsto f(a).
  \end{align*}
  Note que, assim, \(\ker{(\varphi_{a})} = \{f(x)\in A[x]: f(a) = 0\} \trianglelefteq F[x].\) Ainda mais, esse ideal é principal, sendo gerado por 
  \[
    \ker{\varphi_{a}} = \langle (x-a) \rangle = \{g(x)\in F[x]: (x-a)\mid g(x)\}.
  \]
\end{example}
\begin{lemma*}
  As seguintes propriedades são equivalentes:
  \begin{itemize}
    \item[1)] f é monomorfismo;
    \item[2)] \(\ker{(f)} = (0);\)
    \item[3)] Se \(f(a) = 0\), então \(a=0.\)
  \end{itemize}
\end{lemma*}
A prova é análoga ao caso de grupos, então será omitida.
\begin{example}[Exercício]
  \begin{itemize}
    \item[1)] Os anéis com um elemento são todos isomorfos;
    \item[2)] Todo anel com p elementos, sendo p um primo, é isomorfo à \(\mathbb{Z}/p \mathbb{Z};\)
    \item[3)] Os anéis A e B são isomorfos se, e somente se, existirem homomorfismos
      \(f:A\rightarrow B\) e \(g:A\rightarrow B\) tais que \(f\circ{g}=id_{B}\) e \(g\circ{f} = id_{A}.\)
  \end{itemize} 
\end{example}
\begin{prop*}
  Seja \(f:A\rightarrow B\) um homomorfismo de anéis.
  \begin{itemize}
    \item[1)] Se \(\mathfrak{j}\trianglelefteq{B}\), então \(f^{-1}(\mathfrak{j})\trianglelefteq{A}.\) Em particular, \(\ker{(f)} = f^{-1}((0))\trianglelefteq{A};\)
    \item[2)] Se \(\mathfrak{q}\in \mathrm{Spec}(B)\), então \(f^{-1}(\mathfrak{q})\in \mathrm{Spec}(A).\)
    \item[3)] Se A é subanel de B e \(\mathfrak{j}\trianglelefteq{B},\) então \(\mathfrak{j}\cap A \trianglelefteq{A}.\) Se \(i:A\hookrightarrow B\),
      então \(i^{-1}(\mathfrak{j})=\mathfrak{j}\cap A.\)
  \end{itemize}
  \begin{proof*}
    1.) De fato, já sabemos que \(f^{-1}(\mathfrak{j})\leq A\) como grupo. Porém, sejam \(x\in f^{-1}(\mathfrak{j})\)
    e \(y\in A\) tais que \(f(x) = a\in \mathfrak{j}.\) Como \(\mathfrak{j}\) é ideal, \(af(y)\in \mathfrak{j},\)
    ou seja, \(f(x)f(y) = f(xy)\in \mathfrak{j}\). Logo, \(xy\in f^{-1}(\mathfrak{j})\).

    2.) Se \(xy\in f^{-1}(\mathfrak{q}),\) temos \(f(xy) = f(x)f(y)\in \mathfrak{q}.\) Como \(\mathfrak{q}\) é primo, 
    então \(f(x)\in \mathfrak{q}\) ou \(f(y)\in \mathfrak{q}.\) Assim, \(x\in f^{-1}(\mathfrak{q})\) ou \(y\in f^{-1}(\mathfrak{q}).\)

    3.) A primeira parte segue do fato de que subanel é um subconjunto fechado pelas operações restritas. Além disso,
    \(i^{-1}(\mathfrak{j}) = \mathfrak{j}\cap A.\) \qedsymbol
  \end{proof*}
\end{prop*}
\begin{example}[Exercício]
  Seja A um anel. Se \(\mathfrak{p}\trianglelefteq A\) é um ideal primo de A, então 
  \[
    \mathfrak{p}[x] = \{a_{0} + a_{1}x + \dotsc + a_{n}x^{n}:p\mid a_{i}\}
  \]
  é um ideal primo de A[x].
\end{example}
\begin{example}
  Defina \(\alpha :\mathbb{Z}[x]\rightarrow \mathbb{F}_{p}[x]\) por 
  \[
    \alpha \biggl(\sum\limits_{j=0}^{n}a_{j}x^{j}\biggr) = \overline{f(x)}\coloneqq \sum\limits_{j=0}^{n}\overline{a}_{j}x^{n}.
  \]
  Com isso, \(\ker{(\alpha )} = p \mathbb{Z}[x] = \{a_{0} + a_1x + \dotsc +a_{n}x^{n}: p\mid a_{i}\}\). Note que 
  \(\ker{(\alpha )} = f^{-1}((0))\) e, como \(\mathbb{F}_{p}[x]\) é um domínio, segue que \((0)\trianglelefteq \mathbb{F}_{p}[x]\) é um ideal primo.
  Pela proposição recém-provada, isso quer dizer que \(\ker{(\alpha )}\) é um ideal primo. No entanto, ele não é maximal, pois 
  \[
    (p \mathbb{Z})[x] \subsetneq \langle p \mathbb{Z}, x \rangle = p \mathbb{Z} + \langle x \rangle \subsetneq \mathbb{Z}[x].
  \]
\end{example}
Note que a imagem inversa de um ideal maximal não é necessariamente maximal.
Basta considerar a inclusão \(i:\mathbb{Z}\hookrightarrow \mathbb{Q}, (0)\in \mathrm{Specm}(\mathbb{Q}),\)
mas \(i^{-1}((0)) = (0)\not\in \mathrm{Specm}(\mathbb{Z}).\)

Além disso, a imagem de um ideal não é, em geral, um ideal. Para ver isso, basta considerar \(\varphi_{0}:A[x]\rightarrow A[x].\)
Assim, \(\mathrm{Im}(\varphi_{0}) = A\subseteq A[x],\) tal que 
\[
  f(\mathfrak{i}) = \mathrm{Im}(\varphi_{0}) = A
\]
não é um ideal de A[x]. No entanto, é possível pôr uma condição em f para que ocorra o desejado:
\begin{lemma*}
  Seja \(f:A\rightarrow B\) um epimorfismo de anéis.
  \begin{itemize}
    \item[1)] Se \(\mathfrak{i}\trianglelefteq{A},\) então \(f(\mathfrak{i})\trianglelefteq{B}.\) De fato, sabemos que a imagem
      de grupo é um grupo.
    \item[2)] Todo ideal \(\mathfrak{j}\) de B é da forma \(f(\mathfrak{i})\) para algum ideal \(\mathfrak{i}\trianglelefteq{A}.\)
  \end{itemize}
\end{lemma*}
\begin{proof*}
  1.) Já sabemos que a imagem é grupo. Seja \(y\in B.\) Como f é epimorfismo, existe
  \(a\in A\) tal que \(f(a) = y.\) Assim, se \(f(x)\in f(\mathfrak{i}), x\in \mathfrak{i}\) e 
  \(ax\in \mathfrak{i}.\) Então, \(f(ax) = f(a)f(x) = y f(x)\in f(\mathfrak{i}).\)

  2.) Seja \(\mathfrak{j}\trianglelefteq{B}.\) Com isso, \(f^{-1}(\mathfrak{j})\) é ideal de A, pelo lema anterior.
  Como f é epimorfismo, \(f(f^{-1}(\mathfrak{j})) = \mathfrak{j}.\) Assim, \(\mathfrak{j}\) é imagem de um ideal
  \(f^{-1}(\mathfrak{j})\) de A. \qedsymbol
\end{proof*}
\newpage

\section{Aula 05 - 23/08/2023}
\subsection{Motivações}
\begin{itemize}
  \item Anel Quociente.
\end{itemize}
\subsection{Anel Quociente}
\begin{def*}
  Seja A um anel e \(\mathfrak{i}\trianglelefteq{A}.\) Como \((\mathfrak{i}, +)\leq (A, +)\) como grupos e
  como é um subgrupo normal, podemos considerar o \textbf{grupo quociente} \((A/\mathfrak{i}, +)\), em que 
  \[
    A/\mathfrak{i} = \{a + \mathfrak{i}: a \in A\}.\quad\square
  \]
\end{def*}
A operação de soma no grupo quociente é definida pela regra 
\[
  (a+\mathfrak{i})+(b+\mathfrak{i})\coloneqq (a+b)+\mathfrak{i}.
\]
Definamos, também, uma nova operação neste grupo - a multiplicação, dada por 
\[
  (a+\mathfrak{i})\cdot (b+\mathfrak{i})\coloneqq ab + \mathfrak{i}. 
\]
Precisamos mostrar que esta operação está bem definida. Com efeito, suponha que
\(a+\mathfrak{i} = a'+\mathfrak{i}\) e \(b+\mathfrak{i} = b'+\mathfrak{i},\) ou seja,
\(a'-a\in \mathfrak{i}\) e \(b-b'\in \mathfrak{i}.\) Vamos provar que 
\[
  (a+\mathfrak{i})(b+\mathfrak{i}) = (a'+\mathfrak{i})(b'+\mathfrak{i}) \Longleftrightarrow ab - a'b'\in \mathfrak{i}.
\]
Note que podemos somar \(0=a'b-a'b\), a fim de obter 
\[
  ab - a'b + a'b - a'b' = b(a-a') + a'(b-b')\in \mathfrak{i}.
\] 
Como \(b(a-a') + a'(b-b') = ba - a'b',\) mostramos o que queríamos. Por fim,
defina a identidade do produto por \(1 + \mathfrak{i}.\) A partir dessa construção,
definimos o \textbf{anel quociente A por \(\mathfrak{i}\)} como \((A/\mathfrak{i}, +, \cdot )\)
\begin{example}
  Tome o ideal \(3 \mathbb{Z}\) de \(\mathbb{Z}.\) Neste caso, 
  \[
    \mathbb{Z}/3 \mathbb{Z} = \{\overline{r} = r + 3 \mathbb{Z}:r\in \mathbb{Z}\} = \{\overline{0}, \overline{1}, \overline{2}\}.
  \]
  Note que \(2\not\in \mathbb{Z}^{*},\) mas \(\overline{2}\cdot \overline{2} = \overline{1}\) e, portanto, \(\overline{2} = (\overline{2})^{-1},\) ou seja,
  \(\overline{2}\in (\mathbb{Z}/3 \mathbb{Z})^{*}.\)

  Vamos olhar para este exemplo com uma visão mais ampla. Considere, para \(n\in \mathbb{Z},\) o anel quociente 
  \[
    \mathbb{Z}/n \mathbb{Z} = \{\overline{0}, \overline{1}, \overline{2}, \dotsc , \overline{n-1}\}.
  \]
  Com isso, o mapa 
  \begin{align*}
    \varphi : &\mathbb{Z}/n \mathbb{Z}\rightarrow \mathbb{Z}_{1}\\
              &r + n\mathbb{Z}\mapsto \overline{r}
  \end{align*}
  é um isomorfismo de anéis.
\end{example}
\begin{example}
  Considere o anel de polinômios \(\mathbb{R}[x]\) e seu ideal \(\langle x^{2}+1 \rangle\trianglelefteq \mathbb{R}[x]\).
  Como \(x^{2}+1\) é irredutível (afinal, não tem solução em \(\mathbb{R}\)), ele é um ideal primo, cuja forma é 
  \[
    \langle x^{2}+1 \rangle = \{g(x)\cdot (x^{2}+1):g(x)\in \mathbb{R}[x]\}.
  \]
  Neste exemplo, estudaremos o anel quociente 
  \[
    \frac{\mathbb{R}[x]}{\langle x^{2}+1 \rangle} = \{\overline{a} + \overline{b}x: a, b\in \mathbb{R}\}.
  \]
  Dado \(f(x)\in \mathbb{R}[x],\) o algoritmo da divisão garante que existem \(q(x), r(x)\in \mathbb{R}[x]\) tais que 
  \(f(x) = q(x)(x^{2}+1) + r(x),\) em que \(0 \leq \deg{(r(x))} < \deg{(x^{2}+1)} = 2,\) tal que r tem forma \(r(x) = ax + b\).
  Com isso, 
  \begin{align*}
    \overline{f(x)} &= \overline{q(x)}\cdot \overline{(x^{2}+1)} + \overline{r(x)}\\
                    &= \overline{r(x)} = \overline{a}+\overline{b}x.
  \end{align*}
  Em particular, como consquência disso, colocando r(x) = 0 e q(x) = 1, segue que 
  \[
    \overline{x^{2}+1} = \overline{0} \Rightarrow \overline{x}^{2} + \overline{1} = \overline{0} \Rightarrow \overline{x}^{2} = -\overline{1}.
  \]
  Isso induz a existência de um isomorfismo entre o corpo dos complexos e o quociente \(\mathbb{R}[x]/\langle x^{2}+1 \rangle\), explicitado como (exercício)
  \begin{align*}
    \varphi :&\mathbb{C}\rightarrow \mathbb{R}[x]/\langle x^{2}+1 \rangle\\
             &a + bi\mapsto \overline{a} + \overline{b}x.
  \end{align*}
\end{example}
\begin{prop*}[Exercício]
  Seja F um corpo, \(A\neq 0\) um anel e \(\varphi : F\rightarrow A\) um homomorfismo. Então, \(\varphi \) é monomorfismo.
\end{prop*}
\begin{example}
  Vejamos uma construção parecida com a do exemplo anterior. Desta vez, considere o anel quociente 
  \[
    \frac{\mathbb{R}[x]}{\langle x^{2}-1 \rangle} = \{\overline{a}+\overline{b}x:a, b\in \mathbb{R}\}.
  \]
  Mostraremos que \(\mathbb{R}[x]/\langle x^{2}-1 \rangle\cong{\mathbb{R}\times \mathbb{R}}\), com as operações definidas por
  \[
    \left\{\begin{array}{ll}
        (a, b) + (c, d) = (a + c, b + d)\\
        (a, b) \cdot (c, d) = (ac, bd).
    \end{array}\right.
  \]
  Para isso, definiremos a função \(\varphi :\mathbb{R}[x]/\langle x^{2}-1 \rangle\rightarrow \mathbb{R}\times \mathbb{R}\) dada por \(\overline{f(x)}\mapsto (f(1), f(-1)).\)

  \underline{ \(\varphi \)\textbf{ está bem-definido}:} Tome dois elementos com classes de equivalência iguais, ou seja,
  \(\overline{f}(x) = \overline{g}(x)\). Então, \(f(1) = g(1) \) e \(f(-1) = g(-1)\). Com isso,
  \(f(x) - g(x)\in \langle x^{2} - 1 \rangle,\) tal que \(f(x) - g(x) = q(x)(x^{2}-1)\). Logo, 
  \[
    f(1) - g(1) = q(1)(1-1) = 0
  \]
  e 
  \[
    f(-1)-g(-1) = 0.
  \]

  \underline{\(\varphi\) \textbf{ é homomorfismo}:} Para isso, basta perceber que 
  \begin{align*}
    \varphi (\overline{f(x)} + \overline{g(x)}) &= \varphi (\overline{f(x)+g(x)}) \\
                                                &= (f(1) + g(1), f(-1)+g(-1))\\
                                                &= (f(1), f(-1)) + (g(1), g(-1))\\
                                                &= \varphi (\overline{f(x)}) + \varphi (\overline{g(x)}).
  \end{align*}

  \underline{\(\varphi \)\textbf{ é epimorfismo}:} Dado \((a, b)\in \mathbb{R}\times \mathbb{R},\) temos 
  \[
    f(x) = \frac{(a-b)}{2}x + \frac{a+b}{2}. 
  \]

  \underline{\(\varphi \)\textbf{ é monomorfismo}:} Observe que 
  \[
    \varphi (\overline{f}(x)) = 0 = (0, 0) \Rightarrow (f(1), f(-1)) = (0, 0) \Rightarrow f(1) = f(-1) = 0.
  \]
  Disto, segue que 
  \[
    \left.\begin{array}{ll}
        f(1) = 0 \Rightarrow x-1\mid f(x)\\
        f(-1) = 0 \Rightarrow x+1\mid f(x)
      \end{array}\right\{ \Rightarrow x^{2}-1 = (x-1)(x+1)\mid f(x).
      \]
      Portanto, \(f(x)\in \langle x^{2}-1 \rangle,\) ou seja, \(\overline{f(x)} = 0.\)
    \end{example}
    \begin{example}[Exercícios]
      Mostre que \(\langle x^{2} + 1 \rangle \trianglelefteq \mathbb{R}[x]\) é maximal e que \(\langle x^{2}-1 \rangle \trianglelefteq \mathbb{R}[x]\) não é primo. 
    \end{example}

    Vale observe que, junto dessa construção, obtemos acesso canônico à definição do \textbf{mapa quociente} \(\pi :A\rightarrow A/\mathfrak{i},\)
    dado por \(\pi(a) = a + \mathfrak{i}.\) Este mapa é um morfismo de anéis sobrejetor - um epimorfismo.
    Com efeito, a sobrejetividade ocorre pois, se \(a+\mathfrak{i}\in A/\mathfrak{i}\), sabemos que \(\pi(a)= a + \mathfrak{i}.\)
    Sobre a parte do morfismo, se \(a, b\in A, \pi (a+b) = (a+b)+\mathfrak{i} = (a+\mathfrak{i})+(b+\mathfrak{i}) = \pi(a) + \pi(b)\).
    Além disso, \(\pi (ab) = ab + \mathfrak{i} = (a+\mathfrak{i})(b+\mathfrak{i}) = \pi (a)\pi (b).\)
    \begin{prop*}
      Seja \(\mathfrak{i}\trianglelefteq{A}\) um ideal de A.
      \begin{itemize}
        \item[1)] Se \(\mathfrak{j}\trianglelefteq{A}\) e \(\mathfrak{i}\subseteq{\mathfrak{j}}, \) então \(\mathfrak{j}/\mathfrak{i}\trianglelefteq{A/\mathfrak{i}},\) em que 
          \(\mathfrak{j}/\mathfrak{i}=\{a+\mathfrak{i}\}\);
        \item[2)] Se \(\mathfrak{k}\trianglelefteq{A/\mathfrak{i}},\) então \(\mathfrak{k}=\mathfrak{j}/\mathfrak{i}\) para algum \(\mathfrak{j}\trianglelefteq{A};\)
        \item[3)] Exite uma correspondência biunívoca entre os conjuntos:
          \[
            \biggl\{\mathfrak{j}\trianglelefteq{A}: \mathfrak{i}\subseteq{\mathfrak{j}}\biggr\}\longleftrightarrow \biggl\{\mathfrak{k}\trianglelefteq{A/\mathfrak{i}}\biggr\};
          \]
        \item[4)] Se \(\mathfrak{k}\in \mathrm{Spec}(A/\mathfrak{i}),\) então \(\mathfrak{k} = \mathfrak{p}/\mathfrak{i}\) para algum \(\mathfrak{p}\in \mathrm{Spec}(A).\)
          Assim, \(\mathrm{Spec}(A/\mathfrak{i}) = \{\mathfrak{p}/\mathfrak{i}:\mathfrak{i}\subseteq{\mathfrak{p}}\in \mathrm{Spec}(A)\}\);
        \item[5)] Se \(\mathfrak{k}\in \mathrm{Specm}(A/\mathfrak{i}),\) então \(\mathfrak{k} = \mathfrak{m}/\mathfrak{i}\) para algum \(\mathfrak{m}\in \mathrm{Specm}(A).\)
          Com isso, \(\mathrm{Specm}(A/\mathfrak{i}) = \{\mathfrak{m}/\mathfrak{i}: \mathfrak{i}\subseteq{\mathfrak{m}}\in \mathrm{Specm}(A)\}.\)
      \end{itemize}
    \end{prop*}
    \begin{proof*}
      1.) Se \(a+\mathfrak{i}, b+\mathfrak{i}\in \mathfrak{j}/\mathfrak{i}\) e \(r+\mathfrak{i}\in A/\mathfrak{i},\) temos
      \(a, b\in \mathfrak{j}.\) Como \(\mathfrak{j}\) é ideal, \(a+b\in \mathfrak{j}\) e assim \((a+\mathfrak{i}) + (b+\mathfrak{i}) = (a+b)
      +\mathfrak{i}\in \mathfrak{j}/\mathfrak{i}.\) Além disso, \(ar\in \mathfrak{j}\), tal que \((a+\mathfrak{i})(r+\mathfrak{i})=ar + \mathfrak{i}\in
      \mathfrak{j}/\mathfrak{i}.\)

      2.) Tome \(\mathfrak{j} = \pi^{-1}(\mathfrak{k})\trianglelefteq{A}.\) Se \(\mathfrak{k}\trianglelefteq{A/\mathfrak{i}}, 0_{A/\mathfrak{i}}\in \mathfrak{k}\)
      implica que \(\mathfrak{i}\subseteq{\mathfrak{j}}.\) Como \(\pi \) é sobrejetora, podemos escrever
      \(\mathfrak{k} = \pi(\mathfrak{j}) = \{a+\mathfrak{i}: a\in \mathfrak{j}\} = \mathfrak{j}/\mathfrak{i}.\)

      3.) Segue usando 1 e 2.

      4.) Dado \(\mathfrak{k}\in \mathrm{Spec}(A/\mathfrak{i})\), temos \(\mathfrak{k} = \mathfrak{p}/\mathfrak{i},\)
      para \(\mathfrak{i}\subseteq{\mathfrak{p}}\trianglelefteq{A}.\) Vamos, então, provar que \(\mathfrak{p}\) é primo. 
      De fato, se \(ab\in \mathfrak{p},\) então \(\pi (ab)=\pi(a)\pi(b)\in \mathfrak{k}\). Como \(\mathfrak{k}\) é primo,
      isto significa que \(\pi(a)\in \mathfrak{k}\) ou \(\pi(b)\in \mathfrak{k}\). Logo, \(a\in \mathfrak{p}\) ou \(b\in \mathfrak{p},\) do que
      segue que \(\mathfrak{p}\) é primo.

      5.) Suponha que \(\mathfrak{k}\in \mathrm{Specm}(A/\mathfrak{i}).\) Com isso, \(\mathfrak{k} = \mathfrak{m}/\mathfrak{i}\)
      para algum \(\mathfrak{m}\subseteq{\mathfrak{i}}.\) Assuma que \(\mathfrak{j}/\mathfrak{i}\) é outro ideal tal que \(\mathfrak{m}/\mathfrak{i}\subseteq{\mathfrak{j}/\mathfrak{i}}.\)
      Então, \(\mathfrak{j}/\mathfrak{i} = \mathfrak{m}/\mathfrak{i}\) ou \(\mathfrak{j}/\mathfrak{i} = A/\mathfrak{i}\), isto é,
      \(\mathfrak{j} = \mathfrak{m}\) ou \(\mathfrak{j} = A\). Portanto, \(\mathfrak{m} \) é maximal. \qedsymbol
    \end{proof*}
    \begin{example}
      Considere o ideal \(n \mathbb{Z}\subseteq \mathbb{Z}\) e o anel quociente \(\mathbb{Z}/n \mathbb{Z}.\)
      Um ideal \(\mathfrak{j}\trianglelefteq \mathbb{Z}/n \mathbb{Z}\) do anel quociente terá a forma 
      \[
        \mathfrak{j} = \mathfrak{i}/n \mathbb{Z},\quad n \mathbb{Z}\subseteq \mathfrak{i}\subseteq \mathbb{Z}
      \]
      Seja \(\mathfrak{i} = l \mathbb{Z}\) satisfazendo essa condição. Consequentemente, \(l\mid n\). Assim, os ideais de
      \(\mathbb{Z}/n \mathbb{Z}\) são da forma \(\frac{l \mathbb{Z}}{n \mathbb{Z}},\) em que \(l\mid n\). Por exemplo, os ideais
      de \(\mathbb{Z}_{20} = \mathbb{Z}/20 \mathbb{Z}\) serão 
      \[
        (0) = \frac{20 \mathbb{Z}}{20 \mathbb{Z}},\quad \frac{10 \mathbb{Z}}{20 \mathbb{Z}}, \quad \frac{5 \mathbb{Z}}{20 \mathbb{Z}},\quad \frac{4 \mathbb{Z}}{20 \mathbb{Z}}, \quad \frac{2 \mathbb{Z}}{20 \mathbb{Z}}, \quad\&\quad \frac{\mathbb{Z}}{20 \mathbb{Z}}.
      \]
    \end{example}
    \newpage

    \section{Aula 06 - 28/08/2023}
    \subsection{Motivações}
    \begin{itemize}
      \item Relação Entre Anel e Quociente;
      \item Primeiro Teorema do Isomorfismo;
      \item Segundo Teorema do Isomorfismo.
    \end{itemize}
    \subsection{Relação entre Anéis e seus Quocientes}
    \begin{prop*}
      Seja \(\mathfrak{i}\trianglelefteq{A}.\)
      \begin{itemize}
        \item[1)] \(\mathfrak{i}\) é primo se, e somente se, \(A/\mathfrak{i}\) é domínio;
        \item[2)] \(\mathfrak{i}\) é maximal se, e somente se, \(A/\mathfrak{i}\) é corpo.
      \end{itemize}
    \end{prop*}
    \begin{proof*}
      1. \(\Rightarrow ):\) Sejam \(a+\mathfrak{i}, b+\mathfrak{i}\in A/\mathfrak{i}\) tais que \((a+\mathfrak{i})(b+\mathfrak{i})=\mathfrak{i}=0_{A/\mathfrak{i}}.\) Assim,
      \(ab+\mathfrak{i} = \mathfrak{i}\) e \(ab\in \mathfrak{i}.\) como \(\mathfrak{i}\) é primo, \(a\in \mathfrak{i}\) ou \(b\in \mathfrak{i}\), o que implica que
      \(a+\mathfrak{i} = \mathfrak{i} = 0_{A/\mathfrak{i}}\) ou \(b+\mathfrak{i} = \mathfrak{i} = 0_{A/\mathfrak{i}}.\) Logo, \(A/\mathfrak{i}\) é domínio.

      \(\Leftarrow ):\) Suponha agora que \(A/\mathfrak{i}\) é domínio. Tome \(ab\in \mathfrak{i}\). Assim, \(ab + \mathfrak{i} = (a+\mathfrak{i})(b+\mathfrak{i}) = \mathfrak{i}.\)
      Como \(A/\mathfrak{i}\) é domínio, \(a+\mathfrak{i} = \mathfrak{i}\) ou \(b+\mathfrak{i} = \mathfrak{i}\), i.e., \(a\in \mathfrak{i}\) ou \(b\in \mathfrak{i}.\) Destarte, \(\mathfrak{i}\) é primo.

      2. \(\Rightarrow ):\) Suponha que \(\mathfrak{i}\) é maximal e seja \(\mathfrak{k}\trianglelefteq{A/\mathfrak{i}}.\) Pela proposição anterior,
      \(\mathfrak{k} = \mathfrak{j}/\mathfrak{i}\) para algum \(\mathfrak{i}\subseteq \mathfrak{j}\trianglelefteq{A}.\) Como \(\mathfrak{i}\) é maximal,
      \(\mathfrak{i}=\mathfrak{j}\) ou \(\mathfrak{j} = A.\) Com isso, \(\mathfrak{k} = \mathfrak{i}/\mathfrak{i} = (\mathfrak{i})\) ou \(\mathfrak{k} = A/\mathfrak{i}.\) 
      Visto que \(\mathfrak{k}\) é genérico, \(A/\mathfrak{i}\) é um corpo.

      \(\Leftarrow ):\) Assuma que \(A/\mathfrak{i}\) seja um corpo. Seja \(\mathfrak{j}\trianglelefteq{A}\) com \(\mathfrak{i}\mathfrak{j}\).
      Considere o ideal no quociente \(\mathfrak{k} = \mathfrak{j}/\mathfrak{i}\trianglelefteq{A/\mathfrak{i}}\). Como \(A/\mathfrak{i}\) é um corpo,
      segue que \(\mathfrak{k} = (\mathfrak{i})\) ou \(\mathfrak{k} = A/\mathfrak{i}.\) Assim, \(\mathfrak{j} = \mathfrak{i}\) ou \(\mathfrak{j} = A.\)
      Portanto, \(\mathfrak{i}\) é maximal. \qedsymbol
    \end{proof*}
    \begin{example}
      \begin{itemize}
        \item[1)] Se p é um primo, então \(\mathbb{Z}/p \mathbb{Z}\) é um corpo, ou seja, \(p \mathbb{Z}\) é um ideal maximal de \(\mathbb{Z}.\)
        \item[2)] Seja \(\mathfrak{i} = \langle 2 \rangle - 2 \mathbb{Z}\subseteq \mathbb{Z}\) e considere o anel 
          \[
            \mathfrak{i}[x] \trianglelefteq \mathbb{Z}[x],\quad  \mathfrak{i}[x]\coloneqq \{a_{0}+a_{1}x+\dotsc +a_{n}x^{n}\in \mathbb{Z}[x]: a_{i}\in \mathfrak{i}, n\in \mathbb{N}.\}
          \]
          Vamos verificar que o mapa 
          \[
            \varphi : \frac{\mathbb{Z}[x]}{\mathfrak{i}[x]}\rightarrow \biggl(\frac{\mathbb{Z}}{\mathfrak{i}}\biggr)[x] = \mathbb{F}_{2}[x]
          \]
          dado por \(\varphi \biggl(\sum\limits_{i=0}^{n}a_{i}x^{i}\biggr) = \sum\limits_{i=0}^{n}\overline{a_{i}}x^{i}\) é um isomorfismo. Como \(\mathbb{F}_{2}[x]\) é um domínio,
          pelo teorema anterior, \(\mathfrak{i}[x]\) é um ideal primo de \(\mathbb{Z}[x]\).
      \end{itemize}
    \end{example}
    \begin{prop*}[Exercício]
      Seja \(\mathfrak{p}\) um ideal primo de A. Mostre que 
      \[
        \frac{A[x]}{\mathfrak{p}[x]}\cong{\biggl(\frac{A}{\mathfrak{p}}\biggr)}[x]
      \]  
      e que \(\mathfrak{p}[x]\) é um ideal primo de A[x].
    \end{prop*}
    \begin{example}[Exercício]
      Seja p um número primo. Mostre que:
      \begin{itemize}
        \item[1)] \(\mathbb{Z}[1/p]\cong \mathbb{Z}[x]/(\langle px-1 \rangle)\);
        \item[2)] \(\mathbb{Z}[\sqrt[]{p}]\cong \mathbb{Z}[x]/(\langle x^{2}-p \rangle)\);
        \item[3)] Se \(\mathfrak{m}\) é o ideal maximal do anel local \(\mathbb{Z}_{(p)},\) então \(\mathbb{Z}_{(p)}/\mathfrak{m}\cong{\mathbb{Z}/p \mathbb{Z}}.\)
      \end{itemize}
    \end{example}
    \begin{prop*}[Exercício]
      Seja \(\mathfrak{i}\trianglelefteq{A}\) e \(\mathfrak{i}[x]\subseteq A[x]\) o conjunto de polinômios 
      com coeficientes em \(\mathfrak{i}.\) Mostre que:
      \begin{itemize}
        \item[1)] \(\mathfrak{i}[x] \trianglelefteq{A[x]}\);
        \item[2)] \(A[x]/\mathfrak{i}[x]\cong{(A/\mathfrak{i})}[x];\)
        \item[3)] \(\mathfrak{i}\) é primo se, e somente se, \(\mathfrak{i}[x]\) é primo.
      \end{itemize}
    \end{prop*}
    \begin{prop*}[Exercício]
      Se \(\mathfrak{m}\in \mathrm{Specm}(A),\) mostre que \(A/\mathfrak{m}^{n}\) é um anel local para todo
      \(n \geq 1\), com ideal maximal \(\mathfrak{m}/\mathfrak{m}^{n}.\)
    \end{prop*}
    \begin{example}[Exercício]
      Mostre que \(\mathbb{R}[x]/(\langle x^{2}+1 \rangle)\cong{\mathbb{C}}, \mathbb{Q}[x]/(\langle x^{2}+1 \rangle)\cong{\mathbb{Q}[i]}\)
      e \(\mathbb{Z}[x]/(\langle x^{2}+1 \rangle)\cong{\mathbb{Z}[i]}.\)
    \end{example}
    \begin{prop*}[Exercício]
      Seja k um corpo e seja \(\mathfrak{m} = \{x_{1}-a_{1}, \dotsc, x_{n}-a_{n}\}\trianglelefteq{k[x_{1}, \dotsc, x_{n}]}\) para
      \(a_{i}\in k.\) Mostre que \(k[x_{1}, \dotsc, x_{n}]/\mathfrak{m}\cong{k}\), tal que \(\mathfrak{m}\) é maximal. 
    \end{prop*}
    \subsection{Teoremas do Isomorfismo}
    Enunciaremos a seguir dois resultados de extrema importância para a álgebra como um todo.
    \hypertarget{first_isomorphism}{
      \begin{theorem*}[Primeiro Teorema do Isomorfismo]
        Seja \(f:A\rightarrow B\) um homomorfismo de anéis. O mapa \(\overline{f}:A/\ker{(f)}\rightarrow B\)
        é um morfismo injetor de anéis e \(A/\ker{(f)}\cong{\mathrm{Im}(f)}\)
    \end{theorem*}}
    \hypertarget{second_isomorphism}{
      \begin{theorem*}[Segundo Teorema do Isomorfismo]
        Seja \(\mathfrak{i}\trianglelefteq{A}.\) Se \(\mathfrak{j}\subseteq \mathfrak{i}\) e \(\mathfrak{j}/\mathfrak{i}\trianglelefteq{A/\mathfrak{i}},\) então:
        \[
          \frac{A/\mathfrak{i}}{\mathfrak{j}/\mathfrak{i}}\cong{\frac{A}{\mathfrak{j}}}
        \]
    \end{theorem*}}
    \begin{proof*}
      \textbf{\underline{Primeiro Teorema}:} Denotaremos por \(\mathfrak{i}\) o kernel \(\ker{(f)}.\)
      A priori, precisamos checar se \(\overline{f}\) está bem definida. De fato, se \(a+\mathfrak{i}=b+\mathfrak{i}, a-b\in \mathfrak{i}\) e, assim
      \(0 = f(a-b) = f(a)-f(b)\), tal que \(f(a) = f(b).\) Assim, \(\overline{f}(a+\mathfrak{i}) = f(a) = f(b) = \overline{f}(b+\mathfrak{i}).\)

      Agora, conferiremos que \(\overline{f}\) é um morfismo. Com efeito, sejam \(a+\mathfrak{i}, b+\mathfrak{i}\),
      \[
        \overline{f}(a+b+\mathfrak{i}) = f(a+b) = f(a)+f(b) = f(a+\mathfrak{i})+f(b+\mathfrak{i}).
      \]
      A mesma conta prova isto para o produto:
      \[
        \overline{f}(ab+\mathfrak{i}) = f(ab) = f(a)f(b) = f(a+\mathfrak{i})f(b+\mathfrak{i}).
      \]

      A posteriori, provemos a injetividade e o isomorfismo. Para o primeiro, observe que, se \(\overline{f}(a+\mathfrak{i}) = 0\), então
      \(f(a) = 0\) e \(a\in\ker{(f)}\). Assim, \(a+\ker{(f)} = \ker{(f)} = 0_{A/\mathfrak{i}}\). Por fim,
      para ver que \(\overline{f}\) é sobrejetora com a imagem (e, portanto, que há o isomorfismo com a imagem),
      seja \(b\in \mathrm{Im}(f).\) Assim, existe \(a\in A\) tal que \(f(a) = b\). Tome \(\overline{f}(a+\mathfrak{i}) = f(a)=b.\)

      Portanto, segue o teorema.

      \textbf{\underline{Segundo Teorema}:} Defina o mapa \(\theta :A/\mathfrak{i}\rightarrow A/\mathfrak{j},\) dado por \(\theta(a+\mathfrak{i}) = a + \mathfrak{j}.\)
      Deste modo, note que \(\theta \) está bem-definido, pois, se \(a+\mathfrak{i} = b+\mathfrak{i}\), então \(a-b\in \mathfrak{i}\). Como
      \(\mathfrak{i}\subseteq \mathfrak{j}, a-b\in \mathfrak{j}.\) Destarte, \(a+\mathfrak{j} = b+\mathfrak{j}\) e \(\theta(a+\mathfrak{i}) = \theta (b+\mathfrak{i}).\)

      A partir disto, faz sentido checarmos se \(\theta \) é epimorfismo injetor da forma desejada no enunciado. A parte do morfismo segue da conta 
      \[
        \theta (a+b+\mathfrak{i}) = a + b + \mathfrak{j} = (a+\mathfrak{j}) + (b+\mathfrak{j}).
      \]
      A sobrejetividade vem de qualquer elemento \(a+\mathfrak{j}\in A/\mathfrak{j}\) ser escrito por um correspondente \(a+\mathfrak{i}\in A/\mathfrak{i}\).
      Finalmente, a injetividade para o isomorfismo desejado segue ao provarmos que \(\ker{(\theta )}=\mathfrak{j}/\mathfrak{i}.\) Com efeito,
      \(a+\mathfrak{i}\in\ker{(\theta )}\) equivale a dizer que \(a\in \mathfrak{j}\) e \(a+\mathfrak{i}\in \mathfrak{j}/\mathfrak{i}\).

      Portanto, podemos usar o primeiro teorema do isomorfismo e obtemos o resultado desejado. \qedsymbol
    \end{proof*}
    \begin{example}
      Seja \(\mathfrak{p}\) um primo e \(\mathbb{Z}_{(p)} = \biggl\{\frac{a}{b}\in \mathbb{Q}:a, b\in \mathbb{Z}, \text{ p não divide b.}\biggr\} \subseteq \mathbb{Q}\).
      Note que \(\mathbb{Z}_{(p)}\) é subanel de \(\mathbb{Q}.\) Além disso, o conjunto 
      \[
        \mathfrak{i}_{(p)}=\biggl\{\frac{a}{b}\in \mathbb{Q}: p\mid a, p \text{ não divide b}\biggr\}
      \]
      é um ideal de \(\mathbb{Z}_{(p)}\). Vamos estudar este ideal por meio do Teorema do Isomorfismo. Defina a função 
      \begin{align*}
        \varphi : &\mathbb{Z}_{(p)}\rightarrow \mathbb{F}_{p} = \mathbb{Z}/p \mathbb{Z}\\
                  &\frac{a}{b}\mapsto \frac{\overline{a}}{\overline{b}} = \overline{a}\overline{b}^{-1}.
      \end{align*}
      Em particular, como p não divide b, \(\mathrm{mdc}(p, b) = 1,\) tal que existem r, s em \(\mathbb{Z}\) que satisfazem 
      \[
        rp + sb = 1 \Rightarrow \overline{rp + sb} = \overline{1}
      \]
      Como o anel quociente é \(\mathbb{Z}/p \mathbb{Z}\), isto quer dizer que 
      \[
        \overline{r}\overline{0} + \overline{s}\overline{b} = \overline{1} \Rightarrow \overline{s}\overline{b} = \overline{1},
      \]
      tal que \(\overline{b}\) tem inverso multiplicativo no anel quociente. Além disso, \(\varphi \) é epimorfismo. Vamos mostrar que 
      \(\ker{(\varphi )} = \mathfrak{i}_{(p)}.\) 

      Com efeito, para ver que \(\mathfrak{i}_{(p)}\subseteq \ker{(\varphi )},\) dado \(\frac{a}{b}\in \mathfrak{i}_{(p)},\) segue que \(p\mid a\) e p não divide b, tal que 
      \(\overline{a} = \overline{0}\) e \(\overline{b}\neq \overline{0}.\) Com isso, 
      \[
        \varphi \biggl(\frac{a}{b}\biggr) = \frac{\overline{a}}{\overline{b}} = \frac{\overline{0}}{\overline{b}} = \overline{0}.
      \]
      Consequentemente, \(\frac{a}{b}\in \ker{(\varphi )}\) e \(\mathfrak{i}_{(p)}\subseteq \ker{(\varphi )}\).

      Por outro lado, suponha que \(\varphi \biggl(\frac{a}{b}\biggr)=\overline{0}\), tal que \(\frac{\overline{a}}{\overline{b}} = \overline{0}\). Em particular, isto significa que 
      \(\overline{a} = \overline{0}\) e que, então, \(p\mid a\), ou seja, \(\frac{a}{b}\in \mathfrak{i}_{(p)},\) mostrando que \(\ker{(\varphi )}\subseteq \mathfrak{i}_(p)\).

      Assim, mostramos que \(\ker{(\varphi )} = \mathfrak{i}_{(p)}.\) Agora, pelo Teorema do Isomorfismo,
      \[
        \mathbb{Z}_{(p)}/\mathfrak{i}_{(p)} \cong{\mathbb{Z}_{(p)}/\ker{(\varphi )}} \cong{\mathbb{F}_{p}}.
      \]
      Logo, \(\mathfrak{i}_{(p)}\) é um ideal maximal de \(\mathbb{Z}_{(p)}\).

      Podemos levar isto um pouco mais a fundo e concluiremos que \(\mathbb{Z}_{(p)}\) tem apenas dois ideais primos, sendo eles \((0)\) e \(\mathfrak{i}_{(p)}=\langle \mathfrak{p}/\mathfrak{i} \rangle.\)
      Para a unicidade de \(\mathfrak{i}_{(p)}\) como ideal maximal, seja \(\mathfrak{m}\trianglelefteq \mathbb{Z}_{(p)}\) maximal tal que \(\mathfrak{m}\neq \mathfrak{i}_{(p)}\). Seja 
      \(\frac{a}{b}\in \mathfrak{m}\setminus{\mathfrak{i}_{(p)}},\) de forma que p não divide a e nem b, sendo \(\frac{b}{a}\in \mathbb{Z}_{(p)}\). Temos:
      \[
        1 = \frac{b}{a}\cdot \frac{a}{b}\in \mathfrak{m}\Rightarrow \mathfrak{m}=\mathbb{Z}_{(p)}.
      \]
      Contradição. Logo, \(\mathbb{Z}_{(p)}\) não tem outro ideal maximal. 

      Fica para o leitor checar os passos finais de que \(\mathfrak{i}_{(p)} = \langle \mathfrak{p}/\mathfrak{i} \rangle.\) Portanto, \(\mathbb{Z}_{(p)}\) tem apenas
      dois ideais primos, sendo eles \((0)\) e \(\mathfrak{i}_{(p)}.\)
    \end{example}
    \begin{prop*}[Exercício]
      Sejam \(f:A\rightarrow B\) um morfismo de anéis e \(\mathfrak{i}\trianglelefteq{A}.\)
      \begin{itemize}
        \item[1)] Se \(\mathfrak{i}\subseteq \ker{(f)}\), mostre que \(\overline{f}:A/\mathfrak{i}\rightarrow B,\) definido por \(\overline{f}(a+\mathfrak{i}) = f(a)\) é 
          um morfismo bem-definido;
        \item[2)] Se \(\mathfrak{j}\trianglelefteq{B}\), mostre que \(\overline{f}:A/f^{-1}(\mathfrak{j})\rightarrow B/\mathfrak{j},\) definido por
          \(\overline{f}(a+f^{-1}(\mathfrak{j})) = f(a) + \mathfrak{j}\) é um morfismo bem definido e injetor;
        \item[3)] Se \(\mathfrak{j}\trianglelefteq{B}\) e \(\mathfrak{i}\trianglelefteq{A}\) tal que \(\mathfrak{i}\subseteq f^{-1}(\mathfrak{j}),\) então
          \(\overline{f}:A/\mathfrak{i}\rightarrow B/\mathfrak{j}, \) definido por \(\overline{f}(a+\mathfrak{i})=f(a)+\mathfrak{j}\),
          é um morfismo bem definido e \(\ker{(\overline{f})} = f^{-1}(\mathfrak{j})/\mathfrak{i}.\)
      \end{itemize}
    \end{prop*}
    \subsubsection{EXTRA: O Terceiro e o Quarto Teoremas do Isomorfismo}
    Um pouco menos conhecidos, existem dois outros Teoremas do Isomorfismo em Teoria de Anéis. 

    \begin{theorem*}[Terceiro Teorema do Isomorfismo]
      Seja A um anel e \(\mathfrak{i} \trianglelefteq{A}\) um ideal de A. Então, 
      \begin{itemize}
        \item[i)] Se B é um subanel de A, tal que \(\mathfrak{i}\subseteq B\subseteq A,\) então \(B/\mathfrak{i}\) é um
          subanel de \(A/\mathfrak{i};\)
        \item[ii)] Todo subanel de \(A/\mathfrak{i}\) é da forma \(B/\mathfrak{i}\) para algum subanel B de A tal que
          \(\mathfrak{i} \subseteq B \subseteq A\);
        \item[iii)] Se \(\mathfrak{j}\) é um ideal de A tal que \(\mathfrak{i}\subseteq \mathfrak{j}\subseteq A,\) então
          \(\mathfrak{j}/\mathfrak{i}\) é um ideal de \(A/\mathfrak{i}.\)
        \item[iv)] Todo ideal de \(A/\mathfrak{i}\) é da forma \(\mathfrak{j}/\mathfrak{i}\) para algum ideal \(\mathfrak{j}\) tal que
          \(\mathfrak{i}\subseteq \mathfrak{j}\subseteq R.\) 
      \end{itemize}
    \end{theorem*}
    \begin{theorem*}[Quarto Teorema do Isomorfismo]
      Seja \(\mathfrak{i}\) um ideal de A. A correspondência \(A\longleftrightarrow A/\mathfrak{i}\) é uma bijeção, entre os subconjuntos B de A que contêm \(\mathfrak{i}\) e o 
      conjunto de subanéis de \(A/\mathfrak{i}\), que preserva inclusões. Além disso, o subanel B com \(\mathfrak{i} \subseteq B\) é um ideal de A se, e somente se,
      \(B/\mathfrak{i}\) é um ideal de \(A/\mathfrak{i}\).
    \end{theorem*}
    Em algumas fontes é possível encontrar o segundo teorema do isomorfismo como outro resultado (O apresentado aqui seria parte do terceiro), relacionando o quociente da soma de ideais
    com a interseção:
    \begin{theorem*}[Segundo Teorema do Isomorfismo Alternativo]
      Seja A um anel, S um subanel de A e \(\mathfrak{i}\) um ideal de A. Então, 
      \[
        \frac{(S+\mathfrak{i})}{\mathfrak{i}}\cong{\frac{S}{S\cap \mathfrak{i}}}.
      \]
    \end{theorem*}
    \newpage

    \section{Aula 07 - 30/08/2023}
    \subsection{Motivações}
    \begin{itemize}
      \item O Teorema Chinês dos Restos;
      \item Anel Quociente de Polinômios Complexos.
    \end{itemize}
    \subsection{Teorema Chinês dos Restos}
    Antes de mais nada, relembramos a noção de ideais coprimos:
    \begin{quote}
      ``Dois ideais \(\mathfrak{i}, \mathfrak{j}\trianglelefteq{A}\) são ditos \textbf{coprimos} se \(\mathfrak{i}+\mathfrak{j} = A,\) ou seja,
      se \(1\in \mathfrak{i}+\mathfrak{j}.\)''
    \end{quote}
    \begin{example}
      \begin{itemize}
        \item[1)] Se m, n são números inteiros tais que \(\mathrm{mdc}(m, n) = 1,\) então \(m \mathbb{Z} + n \mathbb{Z} = \mathbb{Z}\).
        \item[2)] Seja k um corpo e \(f(x), g(x)\in k[X].\) Se \(\mathrm{mdc}(f(x), g(x)) = 1\), i.e.,
          existe \(h(x)\in k[X]\) com \(\deg(h(x))\geq 1\) tal que \(h(x)\mid f(x)\) e \(h(x)\mid g(x)\), então os ideais
          \(\langle f(x) \rangle, \langle g(x) \rangle\) são coprimos.
        \item[3)] Seja A anel tal que \(\mathfrak{m}_{1}, \mathfrak{m}_{2}\in \mathrm{Specm}(A), \mathfrak{m}_{1}\neq \mathfrak{m}_{2}.\) Então,
          \(\mathfrak{m}_{1} + \mathfrak{m}_{2} \trianglelefteq{A},\) mas \(\mathfrak{m}_1\subseteq \mathfrak{m}_{1}+\mathfrak{m}_{2}.\) 
          Como \(\mathfrak{m}_{1}\) é maximal e \(\mathfrak{m}_{1}\neq \mathfrak{m}_{2},\) resta a opção \(\mathfrak{m}_{1} + \mathfrak{m}_{2} = A.\)
      \end{itemize}
    \end{example}
    \begin{prop*}
      Sejam \(\mathfrak{i}, \mathfrak{j}\trianglelefteq{A}\) com \(\mathfrak{i}+\mathfrak{j} = A.\) Então, para todo \(n\geq 1, m\geq 1,\) os ideais
      \(\mathfrak{i}^{m}\) e \(\mathfrak{j}^{n}\) são coprimos.
    \end{prop*}
    \begin{proof*}
      Se \(\mathfrak{i}, \mathfrak{j}\) são coprimos, existem \(x\in \mathfrak{i}, y\in \mathfrak{j}\) com \(x+y = 1.\)
      Assim, \(1 = 1^{m+n} = (x+y)^{m+n}.\) Expandindo o binômio, obtemos
      \begin{align*}
  &1 = \sum\limits_{k=0}^{m+n}\binom{m+n}{k}x^{k}y^{m+n-k}\\
  &1 = \sum\limits_{k=0}^{m}\binom{m+n}{k}x^{k}y^{m+n-k} + \sum\limits_{k=m+1}^{m+n}\binom{m+n}{k}x^{k}y^{m+n-k}\\
  &1 = \biggl(\sum\limits_{k=0}^{m}\binom{m+n}{k}x^{k}y^{m-k}\biggr)y^{n}+\biggl(\sum\limits_{k=m+1}^{m+n}\binom{m+n}{k}x^{k-m}y^{m+n-k}\biggr)x^{m.}
      \end{align*}
      Assim, escrevemos \(1 = ay^{n}+bx^{m}\), isto é, como a soma de elementos de \(\mathfrak{i}^{m}\) e \(\mathfrak{j}^{n}.\) Portanto,
      \(\mathfrak{i}^{m}\) e \(\mathfrak{j}^{n}\) são coprimos. \qedsymbol
    \end{proof*}
    \hypertarget{chinese_remainder}{
      \begin{theorem*}[Teorema Chinês dos Restos]
        Seja A um anel e \(\mathfrak{i}_{1},\dotsc,\mathfrak{i}_{n}\trianglelefteq{A}\) ideais de A dois a dois coprimos, ou seja,
        \(\mathfrak{i}_{i}+\mathfrak{i}_{j}=A, i\neq j.\) Então:
        \begin{itemize}
          \item[1)] \(\mathfrak{i}_{1}\cdot \dotsc \cdot \mathfrak{i}_{n} = \mathfrak{i}_{1}\cap \dotsc\cap \mathfrak{i}_{n};\)
          \item[2)] Existe um isomorfismo 
            \begin{align*}
  &\varphi :\frac{A}{\bigcap_{i=1}^{n}\mathfrak{i}_{i}}\rightarrow \frac{A}{\mathfrak{i}_{1}}\times \dotsc\times \frac{A}{\mathfrak{i}_{n}}\\
  &a+\bigcap_{i=1}^{n}\mathfrak{i}_{i}\mapsto(a+\mathfrak{i}_{1}, \dotsc, a+\mathfrak{i}_{n}).
            \end{align*}
        \end{itemize} 
    \end{theorem*}}
    \begin{proof*}
      Provaremos os dois fatos juntos utilizando a indução a partir do primeiro caso n=2. 

      Já sabemos que 1 vale para este caso, então iremos encontrar o morfismo. Para tanto, considere \(\varphi :A/(\mathfrak{i}_{1}\cap \mathfrak{i}_{2})\rightarrow A/\mathfrak{i}_{1}\times A/\mathfrak{i}_{2}.\)
      Observe que \(\varphi \) está bem-definido, dado que, se \(a+\mathfrak{i}_{1}\cap \mathfrak{i}_{2} = b + \mathfrak{i}_{1}\cap \mathfrak{i}_{2}, a-b\in \mathfrak{i}_{1}\cap \mathfrak{i}_2\), 
      então \(a+\mathfrak{i}_{1} = b + \mathfrak{i}_{1}\) e \(a + \mathfrak{i}_{2} = b + \mathfrak{i}_{2}.\) Então, \(\varphi(a+\mathfrak{i}_{1}\cap \mathfrak{i}_{2}) = \varphi(b+\mathfrak{i}_{1}\cap \mathfrak{i}_{2}).\)
      Resta provar a injetividade, que é homomorfismo e a sobrejetividade. Faremos nesta ordem

      Suponha \(a+\mathfrak{i}_{1}\cap \mathfrak{i}_{2}\in\ker{(\varphi )}\). Então, \(\varphi (a + \mathfrak{i}_{1}\cap \mathfrak{i}_{2}) = (\mathfrak{i}_{1}, \mathfrak{i}_{2})\),
      isto é, \(a\in \mathfrak{i}_{1}\) e \(a\in \mathfrak{i}_{2}.\) Assim, \(a\in \mathfrak{i}_{1}\cap \mathfrak{i}_{2}\), tal que \(\ker{(\varphi )} = \{0_{A/(\mathfrak{i}_{1}\cap \mathfrak{i}_{2})\} = \mathfrak{i}_{1}\cap \mathfrak{i}_{2}}\)
      e \(\varphi \) é injetora.

      Além disso, \(\varphi (a + b + \mathfrak{i}_{1}\cap \mathfrak{i}_{2}) = (a+b+\mathfrak{i}_{1}, a + b + \mathfrak{i}_{2}) =
      (a + \mathfrak{i}_{1}, a + \mathfrak{i}_{2}) + (b+\mathfrak{i}_{1}, b + \mathfrak{i}_{2}) = \varphi(a + \mathfrak{i}_{1}\cap \mathfrak{i}_{2}) + \varphi(b + \mathfrak{i}_{1}\cap \mathfrak{i}_{2}).\)
      Além disso, \(\varphi (ab + \mathfrak{i}_{1}\cap \mathfrak{i}_{2}) = (ab + \mathfrak{i}_{1}, ab+ \mathfrak{i}_{2}) = (a+\mathfrak{i}_{1}, a+\mathfrak{i}_{2}).\)

      Finalmente, para ver que \(\varphi \) é sobrejetora, tome um elemento genérico \((a+\mathfrak{i}_{1}, b+\mathfrak{i}_{2})\in A/\mathfrak{i}_{1}\times A/\mathfrak{i}_{2}.\)
      Note que essa \(\varphi \) só atinge elementos da forma \((z + \mathfrak{i}_{1}, z + \mathfrak{i}_{2})\) e que 
      \begin{align*}
        (a+\mathfrak{i}_{1}, b+\mathfrak{i}_{2}) &= (a+\mathfrak{i}_{1}, \mathfrak{i}_{2}) + (\mathfrak{i}_{1}, b + \mathfrak{i}_{2})\\
                                                 &= (a+\mathfrak{i}_{1}, a + \mathfrak{i}_{2})(1+\mathfrak{i}_{1}, 0 + \mathfrak{i}_{2}) + (b+\mathfrak{i}_{1}, b+\mathfrak{i}_{2})(0+\mathfrak{i}_{1}, 1 + \mathfrak{i}_{2})\\
                                                 &= \varphi(a)(1+\mathfrak{i}_{1}, 0 + \mathfrak{i}_{2}) + \varphi(b)(0+\mathfrak{i}_{1}, 1 + \mathfrak{i}_{2}).
      \end{align*}
      Com isso, resolvemos nosso problema se encontrarmos \(x_{1}, x_{2}\) tais que \(\varphi (x_{1}+\mathfrak{i}_{1}\cap \mathfrak{i}_{2}) = (0 +\mathfrak{i}_{1}, 1 +\mathfrak{i}_{2})\) e 
      \(\varphi (x_{2} + \mathfrak{i}_{1}\cap \mathfrak{i}_{2}) = (1+\mathfrak{i}_{1}, 0+\mathfrak{i}_{2}).\) Considere
      \(x_{1}\in \mathfrak{i}_{1}, x_{2}\in \mathfrak{i}_{2}\) tais que \(x_{1} + x_{2} = 1.\) Temos:
      \begin{align*}
        \varphi (x_{1}+\mathfrak{i}_{1}\cap \mathfrak{i}_{2}) &= (x_{1}+\mathfrak{i}_{1}, x_{1}+\mathfrak{i}_{2})\\
                                                              &= (\mathfrak{i}_{1}, 1 - x_{2} + \mathfrak{i}_{2})\\
                                                              &= (\mathfrak{i}_{1}, 1+\mathfrak{i}_{2}).
      \end{align*}
      Usando o mesmo raciocínio, 
      \begin{align*}
        \varphi (x_{2} + \mathfrak{i}_{1}\cap \mathfrak{i}_{2}) &= (x_{2}+\mathfrak{i}_{1}, x_{2}+\mathfrak{i}_{2})\\
                                                                &= (1-x_{1}+\mathfrak{i}_{1}, \mathfrak{i}_{2})\\
                                                                &= (\mathfrak{i}_{1}, 1 + \mathfrak{i}_{2}).
      \end{align*}
      Temos todos os elementos que precisávamos. Basta tomarmos \(\varphi (bx_{1} + ax_{2} + \mathfrak{i}_{1}\cap \mathfrak{i}_{2}) = (a+\mathfrak{i}_{1}, b+\mathfrak{i}_{2})\) e, pelos resultados anteriores,
      \(\varphi \) é um isomorfismo de anéis.

      Agora, utilizaremos este caso base para conseguir a hipótese indutiva. Suponha que o resultado vale para \(n-1\) e provemos para n.
      Sejam \(\mathfrak{i}_{1}, \dotsc, \mathfrak{i}_{n}\) ideais de A dois a dois coprimos. Primeiro, note que,
      se \(\mathfrak{j}\coloneqq \mathfrak{i}_{2}\cdot \dotsc \cdot \mathfrak{i}_{n}\), \(\mathfrak{i}_{1}\) e \(\mathfrak{j}\) são coprimos.

      De fato, por hipótese, vale \(\mathfrak{i}_{1}+\mathfrak{i}_{i} = A\) para \(i=2, \dotsc, n.\) Então, existem
      \(x_{1, i}\in \mathfrak{i}_{1}\) e \(y_{i}\in \mathfrak{i}_{i}\) tais que \(x_{1, i}+y_{i} = 1.\) Assim, o produto:
      \[
        1 = (x_{1, 2} + y_{2})(x_{1, 3}+y_{3})\dotsc(x_{1, n}+y_{n}) = y_{2}\cdot \dotsc \cdot y_{n}+z,
      \]
      em que z é o resto da expansão do produto e está em \(\mathfrak{i}_{1},\) porque será uma soma de produtos de termos \(x_{1, i}\in \mathfrak{i}_{1}\) com alguns 
      termos \(y_{},\) que estão no anel, e \(\mathfrak{i}_{1}\) é ideal. Como \(y_{2}\cdot \dotsc \cdot y_{n}\in \mathfrak{j},\) segue que  \(\mathfrak{i}_{1} + \mathfrak{j} = A.\)
      Supondo, por indução, que \(\mathfrak{i}_{2}\cdot \dotsc \cdot \mathfrak{i}_{n} = \mathfrak{i}_{2}\cap \dotsc \cap \mathfrak{i}_{n},\) donde tiramos
      \(\mathfrak{i}_{1}\cdot (\mathfrak{i}_{2}\cdot \dotsc \cdot \mathfrak{i}_{n}) = \mathfrak{i}_{1}\cap (\mathfrak{i}_{2}\cap \dotsc\cap \mathfrak{i}_{n})\) pelo caso \(n=2\).
      Ademais, para a segunda parte do teorema, note que 
      \begin{align*}
        \frac{A}{\mathfrak{i}_{1}\cap(\bigcap_{i=2}^{n}\mathfrak{i}_{i})} = \frac{A}{\mathfrak{i}_{1}\cap \prod\limits_{i=2}^{n}\mathfrak{i}_{i}}\cong \frac{A}{\mathfrak{i}_{1}}\times \frac{A}{\prod\limits_{i=2}^{n}\mathfrak{i}_{i}}&=\frac{A}{\mathfrak{i}_{1}}\times \frac{A}{\bigcap_{i=2}^{n}\mathfrak{i}_{i}}\\  
                                                                                                                                                                                                                                        &\cong \frac{A}{\mathfrak{i}_{1}}\times \dotsc\times \frac{A}{\mathfrak{i}_{i}}.
      \end{align*}
      Portanto, o isomorfismo definido existe. \qedsymbol
    \end{proof*}
    \begin{crl*}
      Sejam \(x_{1}, \dotsc, x_{r}\in \mathbb{Z}\) e \(n_{1}, \dotsc, n_{r}\in \mathbb{N}\) dois a dois coprimos. Então,
      o sistema 
      \begin{equation*}
        \left\{\begin{array}{ll}
            x\equiv x_{1}\mod n_{1}\\
            \vdots\\
            x\equiv x_{r}\mod n_{r}
        \end{array}\right.
      \end{equation*}
      possui uma única solução mod \(n_{1}\cdot \dotsc n_{r}.\)
    \end{crl*}
    \begin{example}
      Seja \(f(x)\in \mathbb{C}[x].\) Pelo teorema fundamental da álgebra, podemos escrever \(f(x) = a(x-a_{1})^{r_{1}}\cdot \dotsc \cdot (x-a_{r})^{r_{n}},\)
      em que \(a_{i}\) são raízes distintas de ordem \(r_{1}, \dotsc, r_{n}\) respectivamente, com \(\deg{(g)} = \sum\limits_{}^{}r_{i}.\) Temos
      \(\langle x-a_{i} \rangle \trianglelefteq{\mathbb{C}[x]}\) é ideal maximal para \(i=1, \dotsc, n.\) Além disso, se \(i\neq j\), os ideais
      \(\langle x-a_{i} \rangle^{r_{i}}\) e \(\langle x-a_{j} \rangle^{r_{j}}\) são coprimos. Aplicando o Teorema Chinês dos Restos, obtemos o isomorfismo:
      \[
        \frac{\mathbb{C}[x]}{\langle f(x) \rangle}\cong \frac{\mathbb{C}[x]}{\langle x-a_{1} \rangle^{r_{1}}}\times \dotsc\times \frac{\mathbb{C}[x]}{\langle x-a_{n} \rangle^{r_{n}}}.    
      \]
      Se temos o caso especial em que f(x) é irredutível, isto é, \(r_{1} = \dotsc = r_{n} = 1\), então:
      \[
        \frac{\mathbb{C}[x]}{\langle x-a_{i} \rangle}\cong{\mathbb{C}},
      \]
      e, assim, 
      \[
        \frac{\mathbb{C}[x]}{\langle f(x) \rangle}\cong{\mathbb{C}^{n}}.
      \]
      Em particular, note que isto significa que 
      \[
        \frac{\mathbb{C}[x]}{\langle x^{2}+1 \rangle}\cong \mathbb{C}.
      \]
    \end{example}
    \begin{example}[Exercícios]
      Sejam A, B anéis. Mostre que:
      \begin{itemize}
        \item[1)] \(\mathrm{Specm}(A\times B) = \{A\times \mathfrak{j}:\mathfrak{j}\in \mathrm{Specm}(B)\}\cup\{\mathfrak{i}\times B:\mathfrak{i}\in \mathrm{Specm}(A)\}.\)
        \item[2)] Encontre todos os ideais maximais de \(A_{1}\times \dotsc\times A_{n}\), em que cada
          \(A_{i}\) é um anel.
        \item[3)] Prove que \(\mathrm{Specm}(\mathbb{C}^{n}) = \mathrm{Spec}(\mathbb{C}^{n}).\) Anéis que satisfazem essa propriedade
          são chamados \textbf{anéis artinianos.}
      \end{itemize}
    \end{example}
    \begin{def*}
      Um anel A é dito \textbf{semilocal} se \(\mathrm{Specm}(A)\) é finito. \(\square\)
    \end{def*}
    \begin{prop*}[Exercício]
      Mostre que para todo polinômio \(f(x)\in \mathbb{C}[x]\) tal que f(x) é não-constante e irredutível, o anel
      \(\mathbb{C}[x]/\langle f(x) \rangle\) é semilocal.
    \end{prop*}
    \newpage

    \section{Aula 08 - 11/09/2023}
    \subsection{Motivações}
    \begin{itemize}
      \item Anéis Multiplicativos;
      \item Anel e Mapa de Localização.
    \end{itemize}
    \subsection{Anéis Multiplicativos}
    Começamos esta seção com exemplo de base. Considere o conjunto \(\mathbb{Q}=\biggl\{\frac{a}{b}: a, b\in \mathbb{Z}, b\neq 0\biggr\}.\)
    Segue que \(\mathbb{Q}\) é um corpo com apenas dois ideais - \((0), \mathbb{Q}.\) Além disso,
    \(\mathbb{Z}\subseteq \mathbb{Q}\) é subanel e \(\mathbb{Z}\) não tem subanel próprio. Como consequência, se
    A é um subanel de \(\mathbb{Q},\) então obrigatoriamente \(\mathbb{Z}\subseteq A\subseteq \mathbb{Q}\). Vamos ver alguns subanéis de \(\mathbb{Q}.\)
    \begin{example}
      \begin{itemize}
        \item[1)] Um subanel é \(\mathbb{Z}\);
        \item[2)] Outro é, para p primo, o conjunto 
          \[
            \mathbb{Z}_{(p)} = \biggl\{\frac{a}{b}\in \mathbb{Q}:\text{ p não divide b}\biggr\},\quad \mathbb{Z}\subseteq \mathbb{Z}_{(p)}\subseteq \mathbb{Q}.
          \]
          Em particular, vimos que este anel é local e tem apenas dois ideais primos: (0) e \(\mathfrak{i}_{(p)}\);
        \item[3)] \(2\mathbb{Z}\) é subanel de \(\mathbb{Q}\);
        \item[4)] \(\mathbb{Z}[\sqrt[]{2}]=\{a + b\sqrt[]{2}: a, b\in \mathbb{Z}\}\subseteq \mathbb{R}\) não é ideal de \(\mathbb{Q}\) pois \(\sqrt[]{2}\) não é racional;
        \item[5)] O conjunto \(\mathbb{Z}\biggl[\frac{1}{2}\biggr]=\biggl\{a + \frac{b}{2}: a, b\in \mathbb{Z}\biggr\}\) \textbf{não} é ideal de \(\mathbb{Q}\). De fato,
          \begin{align*}
  &\biggl(a+\frac{b}{2}\biggr)+\biggl(a'+\frac{b'}{2}\biggr) = (a+a') + \frac{(b+b')}{2} \\
  &\biggl(a+\frac{b}{2}\biggr)\biggl(\frac{a'+b'}{2}\biggr) = aa' + \frac{ab'+a'b}{2} + \frac{bb'}{4}\not\in \mathbb{Z}\biggl[\frac{1}{2}\biggr].
          \end{align*}
      \end{itemize}
    \end{example}
    Este último exemplo possui outras formas alternativas de definí-lo:
    \[
      \mathbb{Z}\biggl[\frac{1}{2}\biggr] = \biggl\{\frac{a}{2^{r}}: r\in \mathbb{Z}, a\in \mathbb{Z}\biggr\} = \biggl\{f \biggl(\frac{1}{2}\biggr): f(x)\in \mathbb{Z}[x].\biggr\}
    \]
    Em particular, a forma de potência fonece uma propriedade interessante de que \(\biggl(\mathbb{Z}\biggl[\frac{1}{2}\biggr]\biggr)^{*} = \{2^{r}: r\in \mathbb{Z}\}.\) Isso pode ser observado
    em anéis parecidos, como em 
    \[
      \mathbb{Z}\biggl[\frac{1}{6}\biggr] = \biggl\{\frac{a}{6^{r}}: r\in \mathbb{Z}, a\in \mathbb{Z}\biggr\} \Rightarrow \biggl(\mathbb{Z}\biggl[\frac{1}{6}\biggr]\biggr)^{*}=\{2^{\alpha }3^{\beta }:\alpha , \beta \in \mathbb{Z}\}.
    \]
    Com base nessas coisas, definimos, para \(\frac{a}{b}\in \mathbb{Q}\) com \(\mathrm{mdc}(a, b) = 1\) e b livre de quadrados, o anel 
    \[
      \mathbb{Z}\biggl[\frac{a}{b}\biggr] = \biggl\{f \biggl(\frac{a}{b}\biggr): f(x)\in \mathbb{Z}[x]\biggr\} \eqqcolon \mathbb{Z}\biggl[\frac{1}{b}\biggr] = \biggl\{\frac{a}{b^{r}}: a\in \mathbb{Z}, r\in \mathbb{Z}\biggr\}.
    \]
    \begin{def*}
      Seja A um anel. Um subconjunto \(S\subseteq A\) é dito \textbf{multiplicativo} se:
      \begin{itemize}
        \item[1)] \(1\in S;\)
        \item[2)] S é fechado pela multiplicação, i.e., \(s, s'\in S\) implica que \(s \cdot s'\in S.\)
      \end{itemize}
    \end{def*}
    \begin{example}
      \begin{itemize}
        \item[1)] Se \(a\in A,\) com \(a\neq 0,\) o conjunto \(S = \{a^{i}: i\geq 0\}\) é multiplicativo;
        \item[2)] Se \(\mathfrak{p}\in \mathrm{Spec}(A),\) então \(S\coloneqq A\setminus{\mathfrak{p}}\) é multiplicativo. De fato, como
          \(\mathfrak{p}\vartriangleleft A, 1\not\in \mathfrak{p}\) e \(1\in A\setminus{\mathfrak{p}}.\)
          Além disso, S é fechado pelo produto. De fato, suponha, por hipótese de absurdo, que existam \(s, s'\in S\) tais que \(s, s'\not\in S.\) Então,
          \(s \cdot s'\in \mathfrak{p}\) e, como \(\mathfrak{p}\) é primo, segue que \(s\in \mathfrak{p}\) ou \(s'\in \mathfrak{p}.\) Mas, então, \(s\not\in S\) ou
          \(s'\not\in S,\) o que é um absurdo. Logo, S é multiplicativo. Em particular, se A for um domínio, temos \((0)\in \mathrm{Spec}(A)\) e, assim,
          \(S=A\setminus{\{0\}}\) é multiplicativo.
        \item[3)] Se \(\mathfrak{p}_{1}, \dotsc , \mathfrak{p}_{n}\in \mathrm{Spec}(A)\) então o conjunto:
          \[
            S = A\setminus{\bigcup_{i=1}^{n}\mathfrak{p}_{i}}
          \]
          é multiplicativo.
        \item[4)] Se \(a, b\in A \setminus{\{0\}},\) então \(S = \{a^{i}\cdot b^{j}: i, j\geq 0\}\) é multiplicativo (o mesmo vale para número finito de
          elementos do anel).
        \item[5)] Se \(f:A\rightarrow B\) é um morfismo de anéis e \(S\subseteq A\) e \(S\cap \ker{(f)} = \emptyset\) é um conjunto multiplicativo,
          então \(f(S)\) é multiplicativo.

          De fato, como \(1_{A}\in S\) e f é morfismo, \(f(1_{A}) = 1_{B}\in f(S).\) Se \(x, y\in f(S),\) existem
          \(a, b\in S\) tais que \(x=f(a)\) e \(y=f(b).\) Como S é multiplicativo, \(a \cdot b\in S\) e, então, \(f(a \cdot b) = f(a)f(b)\in S.\) 
        \item[6)] O conjunto das unidades \(S = A^{*}\subseteq A\) é multiplicativo.
      \end{itemize}
    \end{example}
    Seja A um domínio e \(S\subseteq A\) um conjunto multiplicativo com \(0\not\in S.\) Definiremos uma relação de equivalência no conjunto
    \(A\times S:\) Se \((a, s), (a', s')\in A\times S\), então defina 
    \[
      (a, s)\sim (a', s') \Longleftrightarrow as' = sa'.
    \]
    Provemos que esta relação é de equivalência. É reflexiva, já que, para todo
    \((a, s)\in A\times S, as = sa,\) ou seja, \((a, s)\sim (a, s).\) É simétrica, pois,
    se \((a, s)\sim (a'\Sigma s')\) vale \(as' = sa'\) e, então, \(a's = s'a, \) por comutatividade. 
    Logo, \((a', s')\sim (a, s).\) A transitividade segue do fato de que, se \((a, b), (c, d), (x, y)\in A\times S\) são
    tais que \((a, b)\sim (c, d)\) e \((c, d)\sim (x, y)\), então valem \(ad = bc\) e \(cy = dx.\) Queremos provar
    que \((a, b)\sim (x, y)\). No entanto, note que:
    \[
      bcy = bcy \Rightarrow (da)y = b(xd) \Rightarrow ay = bx.
    \]
    Na última linha, foi utilizada a propriedade dos cancelamentos dos domínios para podermos cancelar d. De fato, se \(d\in S, d\neq 0.\)
    Vale a seguinte propriedade: Se \(Ad = Bd,\) vale \(a = B.\) Caso contrário, teríamos \(A\neq B\) e \(A - B\neq 0,\) mas \(Ad = Bd\) implica que
    \(d(A-B) = 0\) e então d seria um divisor de zero, o que é um absurdo.

    \begin{def*}
      A classe de equivalência de um elemento \((a, s)\in A\times S\) será denotada por 
      \[
        \frac{a}{s} = \{(a', s')\in A\times S: (a, s)\sim (a', s')\}.\quad\square
      \] 
      Denotaremos o conjunto de todas as classes \(S^{-1}A\) o conjunto de todas essas classes de equivalência. 
    \end{def*}
    Podemos caracterizar as classes fazendo:
    \[
      \frac{a}{s} = \frac{a'}{s'} \Longleftrightarrow as'=a's.
    \]
    Em \(S^{-1}A\), podemos definir as operações soma e multiplicação como segue:
    \[
      \frac{a}{s}+\frac{b}{t}\coloneqq \frac{at+bs}{st}\quad\&\quad \frac{a}{s}\frac{b}{t}\coloneqq \frac{ab}{st}.
    \]
    Precisamos mostrar que essas operações estão bem definidas. Quanto à soma,
    \begin{align*}
      \frac{a}{s} + \frac{b}{t} = \frac{a'}{s'} + \frac{b'}{t'} &\Longleftrightarrow \frac{at+bs}{st} = \frac{a't' + b's'}{s't'}\\
                                                                &\Longleftrightarrow (at+bs)(s't') = (a't'+b's')(st)\\
                                                                &\Longleftrightarrow as't t'+bt'ss' = a'st t' + b'tss'.
    \end{align*}
    Note que \((as')t t' = (a's)t t'\) e \((bt')ss' = (b't)ss',\) o que faz com que valam as igualdades acima.

    Agora, para o produto, 
    \begin{align*}
      \frac{a}{s}\frac{b}{t} = \frac{a'}{s'}\frac{b'}{t'} &\Longleftrightarrow \frac{ab}{st}=\frac{a'b'}{s't'}\\
                                                          &\Longleftrightarrow as'bt' = a'sb't. 
    \end{align*}
    Portanto, as operações estão bem-definidas.
    \begin{prop*}
      O conjunto \(S^{-1}A\) com as operações definidas acima é um anel comutativo e com unidade.
    \end{prop*}
    \begin{proof*}
      Note que a classe \(0/s = 0/1\) para todo \(s\in S\). Além disso, se \(a/s\in S^{-1}A,\) temos:
      \[
        \frac{a}{s}+\frac{0}{1} = \frac{a1 + 0s}{s1} = \frac{a}{s}.
      \]  
      Então, \(0/1\) é o elemento neutro da soma neste conjunto. Dado um elemento genérico \(a/s\in S^{-1}A,\) podemos tomar seu inverso \((-a)/s:\)
      \[
        \frac{a}{s}+\frac{-a}{s} = \frac{as + (-a)s}{s^{2}} = \frac{0}{s^{2}} = \frac{0}{1}.
      \]
      A associatividade fica como exercício, assim como a distributiva. A unidade \(1_{S^{-1}A} = 1/1 = s/s \) para todo \(s\in S\). De fato,
      \(a/b \cdot s/s = as/bs = a/b\). A comutatividade fica como exercício também. \qedsymbol
    \end{proof*}
    \begin{lemma*}
      O mapa \(\rho :A\rightarrow S^{-1}A,\) definido por \(\rho (a) = a/1\) é um morfismo de anéis e é injetor, chamado \textbf{mapa de localização}.
    \end{lemma*}
    \begin{proof*}
      Vamos provar, primeiro, que é morfismo. Se \(a, b\in A,\) então \(\rho (ab) = ab/1 = a/1 \cdot b/1 = \rho (a)\rho (b)\) e \(\rho (a+b) = (a1+b1)/1 = a/1 + b/1 = \rho (a) + \rho (b).\)
      Além disso, \(\rho (1) = 1/1 = 1_{S^{-1}A}.\)

      Agora, se \(a\in\ker{(\rho )},\) segue que \(\rho (a) = a/1 = 0/1.\) Isso implica que \(a1 = 0\) e, então, \(a = 0.\) Portanto,
      \(\rho \) é um morfismo injetor. \qedsymbol
    \end{proof*}
    \begin{prop*}
      O anel \(S^{-1}A\) é um domínio. Além disso, \(S^{-1}A = \biggl\{\frac{a}{b}: a, b\in A, b\neq 0\biggr\}\).  
    \end{prop*}
    \begin{proof*}
      Se \(a/s, b/t\in S^{-1}A\) tais que \(a/s \cdot b/t = 0_{S^{-1}A}\), temos \(ab/st = 0/1\) e, assim,
      \(ab = 0.\) Como A é domínio, \(a = 0\) ou \(b = 0\). Portanto,
      \[
        a/s = 0/1\quad\&\quad b/t = 0/1.\quad \text{\qedsymbol}
      \]
    \end{proof*}
    \newpage

    \section{Aula 09 - 13/09/2023}
    \subsection{Motivações}
    \begin{itemize}
      \item Caracterizações da Localização;
      \item Corpo das Frações.
    \end{itemize}
    \subsection{Caracterizando a Localização}
    \begin{prop*}
      Sejam A um domínio, S multiplicativo e \((S^{-1}A, \rho )\) localização de A por S.
      \begin{itemize}
        \item[1)] Se \(\mathfrak{i}\trianglelefteq{A},\) então \(S^{-1}\mathfrak{i} = \{a/s: a\in \mathfrak{i}, s\in S\}\trianglelefteq{S^{-1}A};\) 
        \item[2)] Se \(\mathfrak{j}\trianglelefteq{S^{-1}A}\), existe \(\mathfrak{i}\trianglelefteq{A}\) tal que \(\mathfrak{j} = S^{-1}\mathfrak{i};\)
        \item[3)] Se \(\mathfrak{i}\trianglelefteq{A}\), então \(\mathfrak{i}\cap S \neq\emptyset\) se, e somente se, \(S^{-1}\mathfrak{i} = S^{-1}A;\)
        \item[4)] Vale que \(\mathrm{Spec}(S^{-1}A) = \{S^{-1}\mathfrak{p}:\mathfrak{p}\in \mathrm{Spec}(A) \text{ e } S\cap \mathfrak{p}=\emptyset\}\). Além disso, se considerarmos A como subanel de \(S^{-1}A\),
          então \(S^{-1}\mathfrak{p}\cap A = \mathfrak{p}\) para todo \(\mathfrak{p}\in \mathrm{Spec}(A).\)
      \end{itemize}
    \end{prop*}
    \begin{proof*}
      1.) Seja \(a/s\in S^{-1}A\) e \(b/t\in S^{-1}\mathfrak{i}.\) Assim, \(a\in A\) e \(b\in \mathfrak{i}\) implicam
      que \(ab\in \mathfrak{i},\) e S ser fechado por multiplicação implica que \(st\in S.\) Então, \(a/s \cdot b/t = ab/st\in S^{-1}\mathfrak{i}.\)
      Agora, considere \(b/t, b'/t'\in S^{-1}\mathfrak{i}.\) Como \(t, t'\in S\subseteq A\) e \(\mathfrak{i}\) é ideal, \(bt'+b't\in \mathfrak{i}\) e,
      com isso, \(b/t + b'/t' = (bt'+b't)/t t'\in S^{-1}\mathfrak{i}.\)

      2.) Seja \(\mathfrak{j}\trianglelefteq{S^{-1}A}.\) Tome \(\mathfrak{i} = \rho^{-1}(\mathfrak{j})\trianglelefteq{A}.\) Observe que o
      \(\mathfrak{i} = \{a\in A:\exists s\in S, a/s\in \mathfrak{j}\}\). Seja \(a/s\in S^{-1}\mathfrak{i},\) e então
      \(a\in \mathfrak{i}\) e \(\rho (a)\in \mathfrak{j}.\) Já que \(a/s = 1/s \cdot a/1, 1/s\in S^{-1}A\) e \(a/1\in \mathfrak{j}\),
      junto do fato de \(\mathfrak{j}\) ser ideal, conclui-se que \(a/s\in \mathfrak{j}.\)

      Por outro lado, considere \(x=a/s\in \mathfrak{j}.\) Como \(s/1\in S^{-1}A, a/1 = s/1 \cdot a/s\in \mathfrak{j}.\) Logo, 
      como \(\rho (a) = a/1,\) temos \(a\in \mathfrak{i}\) e \(a/s\in S^{-1}\mathfrak{i}.\)

      3.) Suponha, primeiro, que \(S\cap \mathfrak{i} \neq\emptyset\) e seja \(s\in S\cap \mathfrak{i}.\) Então,
      \(1 = s/s\in S^{-1}\mathfrak{i}.\) Já que \(S^{-1}\mathfrak{i}\) é ideal e contém 1, \(S^{-1}\mathfrak{i} = S^{-1}A.\)

      Agora, suponha que \(S^{-1}\mathfrak{i} = S^{-1}A.\) Segue que \(s/s\in S^{-1}\mathfrak{i},\) o que implica que \(s\in \mathfrak{i}\) para algum
      \(s\in S\) ou \(1\in \mathfrak{i}.\) Portanto, \(1\in S\), e o resultado segue.

      4.) Se \(\mathfrak{p}\in \mathrm{Spec}(A)\) e \(\mathfrak{p}\cap S = \emptyset,\) então \(S^{-1}\mathfrak{p}\vartriangleleft S^{-1}A.\)
      Considere elementos \(a/s, b/t\in S^{-1}A\) tais que \(a/s \cdot b/t\in S^{-1}\mathfrak{p}.\) Então, \(ab/st\in S^{-1}\mathfrak{p}\) e \(ab\in \mathfrak{p}.\)
      Utilizando que \(\mathfrak{p}\) é primo, \(a\in \mathfrak{p}\) ou \(b\in \mathfrak{p}\) implica que \(a/s\in S^{-1}\mathfrak{p}\) ou \(b/t\in S^{-1}\mathfrak{p}.\)
      Logo, \(S^{-1}\mathfrak{p}\in \mathrm{Spec}(S^{-1}A).\)

      Por outro lado, suponha que \(\mathfrak{q}\in \mathrm{Spec}(S^{-1}A).\) tome \(\mathfrak{p}\coloneqq \rho^{-1}(\mathfrak{q})\) e
      sabemos que, como a imagem inversa de um morfismo de um ideal primo é um ideal primo, o item 2 fornece \(\mathfrak{q} = S^{-1}\mathfrak{p}.\) 
      Quanto à última parte da afirmação, considere \(\rho :A\hookrightarrow S^{-1}A\) e então \(S^{-1}\mathfrak{p}\cap A = \rho^{-1}(S^{-1}\mathfrak{p}) = \mathfrak{p} \trianglelefteq{A}.\) \qedsymbol
    \end{proof*}
    \begin{prop*}[Exercício]
      \begin{itemize}
        \item[1)] Se \(f:A\rightarrow B\) é morfismo de domínios, \(S\subseteq A\) tal que \(S\cap\ker{(f)} = \emptyset\), então
          \(f(S)\subseteq B\) é multiplicativo sem o 0, no modelo em que precisamos para realizar a localização. Além disso, \(f|_S:S\rightarrow S\) 
          é injetivo.
        \item[2)] Se \(\mathfrak{p}\in \mathrm{Spec}(A)\) e \(\mathfrak{p}\cap S = \emptyset,\) então \(\overline{S} = \pi (S)\) é multiplicativo sem o 0, com
          morfismo quociente. Além disso, os morfismos
          \begin{align*}
  &\varphi :\frac{S^{-1}A}{S^{-1}\mathfrak{p}}\rightarrow \overline{S}^{-1}\biggl(\frac{A}{\mathfrak{p}}\biggr)\\
  &\frac{a}{s}+S^{-1}\mathfrak{p}\mapsto \frac{\pi (a)}{\pi (s)}.
          \end{align*}
          e 
          \begin{align*}
  &\psi:\overline{S}^{-1}\biggl(\frac{A}{\mathfrak{p}}\biggr)\rightarrow \frac{S^{-1}A}{S^{-1}\mathfrak{p}}\\
  &\frac{\pi (a)}{\pi (s)}\mapsto \frac{a}{s} + S^{-1}\mathfrak{p}
          \end{align*}
          estão bem definidos e são um o inverso do outro, ou seja, obtivemos um isomorfismo 
          \[
            \frac{S^{-1}A}{S^{-1}\mathfrak{p}}\cong{\overline{S}^{-1}\biggl(\frac{A}{\mathfrak{p}}\biggr)}.
          \]
      \end{itemize}  
    \end{prop*}
    \begin{example}
      \begin{itemize}
        \item[1)] Se \(t\in A, t\neq0,\) temos o conjunto multiplicativo \(S=\{1, t, t^{2}, \dotsc\}\).
          Denotaremos \(S^{-1}A\coloneqq A[1/t] = \{a/t^{i}: a\in A \text{ e } i\geq 0\}\). Além disso, 
          \(\mathrm{Spec}(A[1/t]) = \{S^{-1}\mathfrak{p}:\mathfrak{p}\in \mathrm{Spec}(A) \text{ e }t\not\in \mathfrak{p}\}.\)
        \item[2)] Se \(t_{1}, \dotsc, t_{n}\in A\setminus{\{0\}},\) podemos considerar o conjunto multiplicativo \(S = \{t_{1}^{r_{1}}\cdot \dotsc \cdot t_{n}^{r_{n}}: r_{i}\geq 0\}.\)
          Como \textbf{exercício}, encontre todos os ideais primos de \(S^{-1}A.\)
        \item[3)] Se \(\mathfrak{p}\in \mathrm{Spec}(A),\) considere \(S = A\setminus{\mathfrak{p}}.\) Neste caso, denotaremos 
          \(S^{-1}A\coloneqq A_{\mathfrak{p}} = \{a/s: a\in A, s\in A\setminus{\mathfrak{p}}.\}\) Em particular, se \(\mathfrak{p}=(0), A_{(0)}\)
          é um corpo. De fato, se \(a/s\in A_{(0)},\) com \(a/s\neq0/1,\) temos \(a\neq0\) e \(s\neq0.\) Com isso, o elemento
          \(s/a\in A_{(0)}\) e é o inverso multiplicativo de \(a/s\). Chamamos \(A_{(0)}\) de \textbf{corpo de frações} de A,
          denotado por \(\mathrm{Frac}(A).\) Em geral, se S é multiplicativo, \(S^{-1}A\hookrightarrow \mathrm{Frac}(A).\)
      \end{itemize}
    \end{example}
    \begin{prop*}[Exercício]
      \begin{itemize}
        \item[1)] Se \(S\subseteq A^{*},\) mostre que o morfismo localização \(\rho :A\hookrightarrow S^{-1}A\) é um isomorfismo;
        \item[2)] Se A é corpo, mostre que \(\mathrm{Frac}(A)\cong{A;}\)
        \item[3)] Observe que \(\mathrm{Frac}(A)\) é o menor corpo que contém A;
      \end{itemize}  
    \end{prop*}
    \begin{example}[Exercício]
      \begin{itemize}
        \item[1)] Mostre que \(\mathbb{Z}\) não possui subanéis próprios;
        \item[2)] Seja K um corpo. Se K não possui subcorpo próprio, mostre que \(K\cong{\mathbb{Q}}\) ou
          \(K\cong{\mathbb{Z}/p \mathbb{Z}}\) para algum primo p.
      \end{itemize}  
    \end{example}
   \begin{example}
    \begin{itemize}
      \item[1)] Seja \(S = \{2^{n}: n\in \mathbb{N}\cup\{0\}\}\subseteq \mathbb{Z}\).
    Neste caso, o anel de localização com o conjunto multiplicativo S é 
    \[
      S^{-1}\mathbb{Z} = \mathbb{Z}\biggl[\frac{1}{2}\biggr] = \biggl\{\frac{a}{2^{n}}: a\in \mathbb{Z}, n \geq  0\biggr\}.
    \]
    Considere um ideal \(\mathfrak{j}\trianglelefteq \mathbb{Z}\biggl[\frac{1}{2}\biggr] = S^{-1}\mathfrak{i}, \) em que \(\mathfrak{i} = n \mathbb{Z}\).
Assim, 
    \[
      \mathfrak{j} = \biggl\{\frac{nr}{2^{r}}:r\geq 0, x\in \mathbb{Z}\biggr\} = \langle n \rangle = \langle \frac{n}{2} \rangle = n \mathbb{Z}\biggl[\frac{1}{2}\biggr].
    \]
    Disto, conclui-se que \(\mathbb{Z}\biggl[\frac{1}{2}\biggr]\) é um domínio de ideais principais. Podemos ir mais a fundo e analisar os ideais primos e maximais 
de \(\mathbb{Z}\biggl[\frac{1}{2}\biggr].\) Para isso, observe que 
    \[
      \mathrm{Spec}(\mathbb{Z}) = \{(0), p \mathbb{Z}: p \text{ é primo}\} \quad\&\quad \mathrm{Specm}(\mathbb{Z}) = \{p \mathbb{Z}: p \text{ é primo}\}.
    \]
  Desta forma,
  \begin{align*}
    \mathrm{Spec}\biggl(\mathbb{Z}\biggl[\frac{1}{2}\biggr]\biggr) &= \{S^{-1}p: p\in \mathrm{Spec}(\mathbb{Z}), p\cap S = \emptyset \} \\
                                                                   &= \{(0), \langle \frac{p}{2} \rangle: p\text{ primo}, p \neq 2\}\\
                                                                   &= \{(0), \langle p \rangle: p \text{ primo}, p\neq 2\}.
  \end{align*}
  Note que, se \(p \mathbb{Z}\cap S = \emptyset \), então \(p\neq 2\). Da mesma forma que vimos no exemplo anterior desses conjuntos,
o raciocínio feito para \(\mathrm{Spec}\biggl(\mathbb{Z}\biggl[\frac{1}{6}\biggr]\biggr),\) de modo que 
  \[
    \mathrm{Spec}\biggl(\mathbb{Z}\biggl[\frac{1}{6}\biggr]\biggr) = \{(0), \langle p \rangle: p \text{ é primo}, p\neq 2, 3\}.
  \]
  É possível estender esse raciocínio de modo que 
    \[
      \mathbb{Z}\biggl[\frac{1}{3}\biggr]\subseteq \mathbb{Z}\biggl[\frac{1}{3}, \frac{1}{5}\biggr]\subseteq \mathbb{Z}\biggl[\frac{1}{3}, \frac{1}{5}, \frac{1}{7}\biggr]\subseteq \dotsc \subseteq \mathbb{Q}.
    \]
  Com isso, 
    \[
      A = \bigcup_{p \text{ primo}, p\neq 2}^{}\mathbb{Z}\biggl[\frac{1}{3}, \dotsc , \frac{1}{p}\biggr] = \mathbb{Z}_{(2)} = \biggl\{\frac{a}{b}\in \mathbb{Q}: a, b\in \mathbb{Z}, \text{2 não divide b}\biggr\}.
    \]
    Da mesma forma que estudamos os primos de \(\mathbb{Z}\biggl[\frac{1}{2}\biggr],\) faremos isso com os primos de A. Considere o conjunto mulitplicativo
 \(S = A\setminus{\mathfrak{p}}\) e o anel de localização \(A_{\mathfrak{p}}\coloneqq S^{-1}A.\) Se \(A = \mathbb{Z}\) e \(p = (2),\) então 
    \[
      A_{p} = \mathbb{Z}_{(2)} = \biggl\{\frac{a}{s}: a\in \mathbb{Z}, s\in \mathbb{Z}\setminus{2 \mathbb{Z}}\biggr\} \Rightarrow \mathbb{Z}_{(2)}\text{ é D.I.P.}.
    \]
  Para todo primo diferente de 2, \(p\in S = \mathbb{Z}\setminus{2 \mathbb{Z}}.\) Logo, p tem inverso em \(\mathbb{Z}_{(2)}\) dado por \(\frac{1}{p}\). Consequentemente, 
    \[
      \mathrm{Spec}(\mathbb{Z}_{(2)}) = \biggl\{S^{-1}\mathfrak{p}: \mathfrak{p}\in \mathrm{Spec}(\mathbb{Z}), p\cap S = \emptyset \biggr\} = \{(0), \langle \frac{2}{1} \rangle\}.
    \]
  Analogamente, \(\mathrm{Spec}(\mathbb{Z}_{(p)}) = \{(0), \langle \frac{p}{1} \rangle\}, \mathrm{Specm}(\mathbb{Z}_{(p)}) = \{\langle p \rangle\}\).
  \end{itemize}
  \end{example}
  \begin{example}
    Seja F um corpo e \(A = F[x]\). Dado um ideal \(\mathfrak{i}\trianglelefteq F[x]\), sabe-se que F[x] é D.I.P. e que 
    \[
      \mathfrak{i} = \langle f(x) \rangle = \{f(x)\cdot h(x):h(x)\in F[x]\}.
    \]
    Segue que \(\mathrm{Spec}(F[x]) = \{(0), \langle f(x) \rangle: f \text{ é irredutível}\}\) e que \(\mathrm{Specm}(F[x]) = \{\langle f(x) \rangle: f \text{ é irredutível}\}\).
Fazendo \(\mathfrak{p} = \langle f(x) \rangle\in \mathrm{Spec}(F[x]),\) temos 
  \[
    A_{\mathfrak{p}} = F[x]_{p} = \biggl\{\frac{g(x)}{h(x)}: g(x)\in F[x], h(x)\in F[x]\setminus{\mathfrak{p}}\biggr\} = \biggl\{\frac{g(x)}{h(x)}: g(x)\in F[x], f(x)\text{ não divide }h(x)\biggr\}.
  \]
  Com isto, \(\mathrm{Spec}(F[x]_{\mathfrak{p}}) = \{(0), \langle \frac{f(x)}{1} \rangle\}\), ou seja, \(F[x]_{\mathfrak{p}}\) é um ideal local com ideal maximal \(\langle \frac{f(x)}{1} \rangle.\)
  \end{example}
    \begin{prop*}
      Seja \(\mathfrak{p}\in \mathrm{Spec}(A)\). Então, o domínio \(A_{\mathfrak{p}}\) é um anel local
      com ideal maximal \(\mathfrak{p}A_{\mathfrak{p}}\coloneqq \{a/s: a\in \mathfrak{p}\text{ e }s\not\in \mathfrak{p}\}.\) Além disso,
      \[
        \frac{A_{\mathfrak{p}}}{\mathfrak{p}A_{\mathfrak{p}}}\cong{\mathrm{Frac}\biggl(\frac{A}{\mathfrak{p}}\biggr).}
      \]
      Em particular, se \(\mathfrak{m}\in \mathrm{Specm}(A),\) então:
      \[
        \frac{A_{\mathfrak{m}}}{\mathfrak{m}A_{\mathfrak{m}}}\cong{\mathrm{Frac}\biggl(\frac{A}{\mathfrak{m}}\biggr)} \cong{\frac{A}{\mathfrak{m}}}
      \]
    \end{prop*}
    \begin{proof*}
      Observe que \(A_{\mathfrak{p}}=S^{-1}A,\) com \(S=A\setminus{\mathfrak{p}}.\) É claro que \(\mathfrak{p}A_{\mathfrak{p}} = S^{-1}\mathfrak{p}\), e então
      \(\mathfrak{p}A_{\mathfrak{p}}\in \mathrm{Spec}(A_{\mathfrak{p}}).\) Provemos que este ideal é maximal e é o único.

      Seja \(\mathfrak{p}a_{\mathfrak{p}}\subsetneq \mathfrak{j}\trianglelefteq{A_{\mathfrak{p}}}. \) Então, \(\mathfrak{j} = S^{-1}\mathfrak{i}\) para algum \(\mathfrak{i}\trianglelefteq{A}\) e 
      \(\mathfrak{p}A_{\mathfrak{p}} = S^{-1}\mathfrak{p}\subsetneq S^{-1}\mathfrak{i}\subseteq S^{-1}A \), tal que \(\mathfrak{p}\subsetneq \mathfrak{i}\subseteq A. \)
      Como \(\mathfrak{p}\neq \mathfrak{i}, \mathfrak{i}\cap A\setminus{\mathfrak{p}} = \mathfrak{i}\cap S \neq\emptyset\). Logo,
      \(S^{-1}\mathfrak{i} = \mathfrak{j} = A_{\mathfrak{p}}.\) Assim, provamos que \(\mathfrak{p}A_{\mathfrak{p}}\) é maximal em \(S^{-1}A.\) 

      Agora, considere \(\mathfrak{k}\in \mathrm{Specm}(S^{-1}A)\). Então, \(\mathfrak{k} = S^{-1}\mathfrak{i}\) para algum ideal 
      \(\mathfrak{i}\trianglelefteq{A}.\) Suponha que \(\mathfrak{k}\neq \mathfrak{p}A_{\mathfrak{p}}.\) Então, \(\mathfrak{i}\neq \mathfrak{p}\) e
      \(\mathfrak{i}\cap A\setminus{\mathfrak{p}}\neq\emptyset.\) Com isso, \(S^{-1}\mathfrak{i} = S^{-1}A = A_{\mathfrak{p}},\) o que é um absurdo com a
      suposição de que \(\mathfrak{k}\) é maximal e, então, próprio.

      Como \(\mathfrak{p}A_{\mathfrak{p}}\in \mathrm{Spec}(A_{\mathfrak{p}}), A_{\mathfrak{p}}/\mathfrak{p}A_{\mathfrak{p}}\) é um corpo. Além disso,
      note que os mapas 
      \begin{align*}
  &\rho :\frac{A_{\mathfrak{p}}}{\mathfrak{p}A_{\mathfrak{p}}}\rightarrow \mathrm{Frac}\biggl(\frac{A}{\mathfrak{p}}\biggr)\\
  &\frac{a}{s} + \mathfrak{p}A_{\mathfrak{p}}\mapsto \frac{\pi (a)}{\pi (s)}
      \end{align*}
      e
      \begin{align*}
  &\psi:\mathrm{Frac}\biggl(\frac{A}{\mathfrak{p}}\biggr)\rightarrow \frac{A_{\mathfrak{p}}}{\mathfrak{p}A_{\mathfrak{p}}}\\
  &\frac{\pi (a)}{\pi (s)}\mapsto \frac{a}{s} + \mathfrak{p}A_{\mathfrak{p}}
      \end{align*}
      estão bem definido e são um o inverso do outro, provando o teorema (fica de exercício mostrar isso). \qedsymbol
    \end{proof*}
  \begin{example}
    Sejam F um corpo, \(A = F[x, y], \mathfrak{p} = \langle x \rangle\). Tome dois elementos \(f = f(x, y), g = g(x, y)\in A\) tais que 
 \(f, g\in \langle x \rangle.\) Equivalentemente, \(x\mid fg\), ou seja, \(x\mid f\) ou \(x\mid g\). Quando fazemos o quociente, obtemos 
  \[
    \frac{A}{\langle x \rangle} = \frac{F[x, y]}{\langle x \rangle}\cong F[y], \quad \frac{F[x, y]}{\langle x, y \rangle}\cong F.
  \]
  Considerando, então, \(A_{\mathfrak{p}},\) chegamos em 
    \[
      A_{\langle x \rangle}= \biggl\{\frac{f}{g}: f\in A, g\in A\setminus{\langle x \rangle}\biggr\} = \biggl\{\frac{f}{g}: f\in A, \text{x não divide }g(x)\biggr\}.
    \]
  Em particular, para todo f(y) em F[x] tal que x não divide f(y), segue que \(\frac{1}{f(y)}\in A_{\langle x \rangle}.\) Com isso,
    \[
      F(y)[x]_{\langle x \rangle}\subseteq A_{\langle x \rangle}
    \]
  e, portanto, \(A_{\langle x \rangle} = F(y)[x]_{\langle x \rangle}\). Ademais, 
    \[
      \frac{A_{\langle x \rangle}}{\langle x \rangle A_{\langle x \rangle}} \cong \mathrm{Frac}\biggl(\frac{F[x, y]}{\langle x \rangle}\biggr) \cong \mathrm{Frac}(F[y]) = F[y].
    \]
  Podemos concluir, então, que \(\mathrm{Spec}(A_{\langle x \rangle}) = \{S^{-1}\mathfrak{p}:\mathfrak{p}\in \mathrm{Spec}(A), \mathfrak{p}\subseteq \langle x \rangle\} = \{(0), \langle x \rangle A_{\langle x \rangle}\}\).
  \end{example}
  \begin{example}
    Para \(A = F[x, y]\) novamente, considere \(\mathfrak{p} = \langle x, y \rangle\). O Anel local \(A_{\mathfrak{p}} = F[x, y]_{\langle x, y \rangle} = \biggl\{\frac{f}{g}: g, f\in F[x, y], g\not\in \langle x, y \rangle\biggr\}.\)
Neste caso, \(\mathrm{Spec}(A_{\mathfrak{p}}) = \{S^{-1}Q: Q\subseteq \langle x, y \rangle\} = \{(0), \langle \frac{x}{1} \rangle, \langle \frac{y}{1} \rangle, \underbrace{\langle \frac{x}{1}, \frac{y}{1} \rangle}_{\text{maximal}}\}\).
  \end{example}
    \begin{prop*}[Exercício]
      Sejam \(\mathfrak{p}_{1}, \dotsc, \mathfrak{p}_{n}\in \mathrm{Spec}(A)\) tais que \(\mathfrak{p}_{i}\subsetneq \mathfrak{p}_{j} \) para todo
      \(i\neq j\). Se \(S = A\setminus{\bigcup_{}^{}\mathfrak{p}_{i}}\), então \(S^{-1}A\) é um anel semilocal, com \(\mathrm{Specm}(S^{-1}A) = \{S^{-1}\mathfrak{p}_{1}, \dotsc, S^{-1}\mathfrak{p}_{n}\}\).
    \end{prop*}
    \newpage

    \section{Aula 10 - 18/09/2023}
    \subsection{Motivações}
    \begin{itemize}
      \item Polinômios em Anéis;
      \item Algoritmo de Divisão;
      \item Exemplos.
    \end{itemize}
    \subsection{Anéis de Polinômios}
    \begin{def*}
      Seja A um anel. O \textbf{anel de polinômios com coeficientes em A}, \(A[x]\), é definido como 
      \[
        A[x]\coloneqq \{a_{n}x^{n}+\cdots+a_{1}x + a_{0}: a_{i}\in A, i = 1, \cdots, n\},
      \]
      em que x é chamado de variável. \(\square\)
    \end{def*}
    Nestes estudos, assumiremos sempre algumas coisas:
    \begin{itemize}
      \item[i)] \(ax = xa\) para todo a em A;
      \item[ii)] \(a_{n}x^{n} + \cdots + a_{0} = 0\) se, e somente se, \(a_{i} = 0\) para todo \(i=1, \cdots, n.\)
      \item[iii)] \(a_{n}x^{n} + \cdots + a_{1}x + a_{0} = b_{m}x^{m} + \cdots + b_{1}x + b_{0}\) se, e somente se, \(n=m\) e \(a_{i} = b_{i}\) para todo \(i=1, \cdots, n\).
    \end{itemize}

    \begin{def*}
      Seja \(f(x) = a_{n}x^{n} + \cdots + a_{0}\in A[x]\). Definimos
      \begin{itemize}
        \item[1)] Se \(a_{n}\neq0, a_{n}\) é chamado o \textbf{coeficiente líder} de f(x). Denotamos o coeficiente líder
          de f(x) por \(\ell_{f}\) ou \(\ell_{f(x)}\)
        \item[2)] Se \(a_{n}\neq0\), n é chamado o \textbf{grau} de f(x), denotado por \(\deg(f(x))\)
        \item[3)] \(a_{0}\) é chamado o \textbf{coeficiente constante} de f(x)
        \item[4)] Diz-se que \(\alpha\in A\) é uma raiz de \(f(x)\in A[x]\) se \(f(\alpha ) = 0.\quad\square\)
      \end{itemize}
    \end{def*}
    Algumas propriedades simples que temos são:
    \begin{itemize}
      \item[i)] \(\deg{(f(x) + g(x))}\leq \max\{\deg{f(x)}, \deg{g(x)}\}\). A igualdade ocorre se, e somente se,
        \(\deg{f(x)}\neq \deg{g(x)}\) ou \(\deg{f(x)}=\deg{g(x)},\) mas \(\ell_{f} + \ell_{g}\neq0.\)
      \item[2)] \(\deg{(f(x)\cdot g(x))}\leq \deg{f(x)} + \deg{g(x)}.\) A igualdade ocorre se, e somente se,
        \(\ell_{f}\cdot \ell_{g} \neq0\)
    \end{itemize}
    Fica como um bom exercício prová-las.
    \begin{example}
      Se A é um domínio, então 
      \[
        \deg{(f \cdot g)} = \deg{f} + \deg{g},
      \]
      para todo \(f, g\neq0.\)
    \end{example}
    \begin{example}
      Em \(\biggl(\mathbb{Z}/4\biggr)[x],\) temos 
      \[
        2x \cdot 2x = 0,
      \]
      tal que \(\deg{2x \cdot 2x} = 0 < 2 = \deg{2x} + \deg{2x}.\)
    \end{example}
    \begin{def*}
      Um polinômio \(f(x)\in A[x]\) é dito \textbf{mônico} se \(\ell_{f} = 1.\) Em outras palavras, 
      \[
        f(x) = x^{n} + a_{n-1}x^{n-1} + \cdots + a_{0}.
      \]
      Um polinômio mônico \(f(x)\) é dito \textbf{irredutível} se não existem polinômios mônicos
      \(g(x), h(x)\) tais que 
      \[
        f(x) = g(x)\cdot h(x).\square
      \]
    \end{def*}
    \begin{example}
      Em \(\mathbb{R}[x]\), o polinômio mônico \(x^{2} + 1\) é irredutível. De fato, suponha que existam 
      \(h(x), g(x)\) tais que 
      \[
        x^{2} + 1 = h(x)g(x).
      \]
      Então,
      \[
        x^{2} + 1 = h(x)g(x) = (x-a)(x-b) = x^{2} - (a+b)x +ab.
      \]
      Podemos reescrever isso em forma de sistema de equações 
      \[
        \left\{\begin{array}{ll}
            a + b = 0\\
            ab = 1
        \end{array}\right. 
        \Rightarrow  
        \left\{\begin{array}{ll}
            a = -b\\
            a(-a) = 1
        \end{array}\right.
        \Rightarrow 
        a^{2} = -1,
      \]
      mas isso não é possível para coeficientes reais. Portanto, não existe tal redução de \(x^{2} + 1.\)
    \end{example}
    Se \(f(x)\in \mathbb{R}[x]\) é irredutível, então 
    \[
      1\leq \deg{f(x)}\leq 2.
    \]
    Suponha, por outro lado, que \(f(x)\in \mathbb{R}[x]\) tem grau n, isto é, \(\deg{f(x)} = n.\) Então, \(f(x)\) tem n raízes complexas,
    digamos \(\alpha_{1}, \cdots, \alpha_{n}\). Assim, 
    \[
      f(x) = (x-\alpha_{1})\cdots(x-\alpha_{n}).
    \]
    Vamos separar essas raízes nas puramente reais e nas puramente complexas. Sejam \(\alpha_{1}, \cdots, \alpha_{h}\in \mathbb{R}\) e
    \(\alpha_{h+1}, \cdots, \alpha_{n}\in \mathbb{C}\setminus{\mathbb{R}}.\) Com isso, 
    \[
      f(x) = g(x)h(x),
    \]
    em que \(g(x) = (x-\alpha_{1})\cdots(x-\alpha_{h})\) e \(h(x) = (x-\alpha_{h+1})\cdots(x-\alpha_{n}).\)
    Se \(z = a + ib\in \mathbb{C}\) é uma raiz de \(h(x),\) então \(\overline{z} = a - ib\) também é uma raiz de h, ou seja, 
    \[
      (x-z)(x-\overline{z})\mid h(x).
    \]
    Se \(h(x) = (x-z)(x-\overline{z})h_{1}(x) = (x^{2}-2a + (a^{2}+b^{2}))h_{1}(x)\), o termo 
    \(x^{2}-2a + (a^{2}+b^{2})\) é irredutível em \(\mathbb{R}[x]\). Continuando assim, podemos ver que 
    \[
      f(x) = (x-a)\cdots(x-a_{n})(x^{2}+a_{h+1}x + b_{h+1})(x^{2}+a_{h+2}x+b_{h+2})\cdots(x^{2}+a_{n}x + b_{n}),\quad n = h = 2m.
    \]
    Uma consequência deste raciocínio é que um polinômio complexo \(f(x)\in \mathbb{C}[x]\) é irredutível se, e somente se, \(\deg{f(x)} = 1.\)
    \begin{example}
      Em \(\mathbb{Q}[x]\), os polinômios irredutíveis podem ter qualquer grau. Polinômios irredutíveis em \(\mathbb{Q}[x]\)
      são, por exemplo, 
      \begin{align*}
  &x-a,\quad a\in \mathbb{Q}\\
  &x^{2}-2ax + (a^{2}+b^{2}),\quad a, b\in \mathbb{Q}\\
  &x^{2} - p,\quad p\in \mathbb{Z}\text{ livre de quadrados}\\
  &x^{n} - 2.
      \end{align*}
    \end{example}
    \begin{def*}
      Dados \(f, g\in A[x]\), dizemos que \textbf{f divide g,} denotado \(f\mid g\), se existe \(h\in A[x]\) tal que 
      \[
        g(x) = f(x)h(x).\quad\square
      \]
    \end{def*}
    \begin{lemma*}
      Se \(f(x)\in A[x]\) é mônico, então existem polinômios mônicos irredutíveis \(f_{1}, \cdots, f_{m}\in A[x]\) tais que 
      \[
        f(x) = f_{1}(x)^{r_{1}}\cdot \dotsc \cdot f_{m}(x)^{r_{m}}.
      \]
    \end{lemma*}
    \begin{proof*}
      Exercício.
    \end{proof*}
    Agora, vejamos o \textbf{Algoritmo da Divisão:}
    \begin{theorem*}
      Sejam \(f(x), g(x)\in A[x], \deg{g(x)}\geq 1\) e \(\ell_{g}\in A^{*}.\) Então,
      existem polinômios únicos \(q(x)\) e \(r(x)\) em \(A[x]\) satisfazendo 
      \[
        f(x) = q(x)g(x) + r(x),
      \]
      com \(\deg{r(x)} < \deg{g(x)}.\)
    \end{theorem*}
    \begin{proof*}
      A prova é por indução no grau de f.

      Se \(\deg{f}=0,\) f é constante. Como \(\deg{g}\geq 1, 0 = \deg{f} < \deg{g},\) g não é constante.
      Tome \(q(x) = 0\) e \(r(x) = f(x).\) Então, \(f(x) = 0g(x) + f(x)\) e \(\deg{r} < \deg{g}.\)

      Suponha agora que o resultado seja válido para polinômios de grau menor que n e sejam \(f(x) = \sum\limits_{i=0}^{n+1}a_{i}x^{i}, g(x) = \sum\limits_{i=0}^{m}b_{i}x^{i}, b_{m}\in A^{*}\).
      Se \(\deg{f(x)} < \deg{g(x)},\) basta copiarmos o caso em que f era constante, colocando \(q(x) = 0\) e \(r(x) = f(x)\).
      Assim, podemos assumir que \(\deg{f(x)}\geq \deg{g(x)}.\) Neste caso, o polinômio 
      \[
        f_{1}(x)\coloneqq f(x) - a_{n+1}b_{m}^{-1}x^{n+1-m}g(x) = c_{n}x^{n} + \cdots + c_{0}
      \] 
      é de grau menor ou igual que n. Aplicando a hipótese de indução, existem \(q_{1}(x), r_{1}(x)\) 
      com \(\deg{r_{1}(x)} < \deg{g(x)},\) de forma que \(f_{1} = gq_{1} + r_{1}.\) Temos:
      \begin{align*}
        f(x) &= f(x)-a_{n}b_{m}^{1}x^{n+1-m}g(x) = q_{1}(x)g(x) + r_{1}(x)\\
             &= (a_{n}b_{m}^{-1}x^{n+1-m}+q_{1}(x))g(x) + r_{1}(x).
      \end{align*}
      Definindo \(q(x) = a_{n}b_{m}^{-1}x^{n+1-m}+q_{1}(x), r = r_{1},\) temos o resultado. Pela hipótese
      de indução, ele vale para todo \(n\in \mathbb{N},\) restando apenas provar que os polinômios são únicos.

      Destarte, suponha que existem \(q_{1}, r_{1}, q_{2}, r_{2}\in A[x]\) tais que 
      \[
        f(x) = q_{1}g+r_{1} = q_{2}g+r_{2},
      \]
      com \(\deg{r_{1}}, \deg{r_{2}} < \deg{g}.\) Temos \(r_{1}-r_{2} = (q_{2}-q_{1})g.\) Suponha que \(q_{2}\neq q_{1},\)]
      de modo que \(q_{2}-q_{1}\neq0\) e 
      \[
        \deg{(r_{1}-r_{2})} = \deg{(q_{2}-q_{1})}g = \deg{(q_{2}-q_{1})} + \deg{g},
      \]
      em que a igualdade do produto dos graus ocorre pois, se \(f, g\in A[x]\) e o coeficiente líder de g é \(b\in A^{*}\),
      então \(\deg{fg} = \deg{f} + \deg{g}.\) De fato, \(f(x)g(x) = a_{n}b_{m}x^{n+m} + \cdots + a_{0}b_{0}\) com
      \(a_{n}b_{m}\neq0\), já que, se \(a_{n}b_{m} = 0, a_{n} = 0.\) Assim, o grau do produto é \(n+m\).

      Logo, da igualdade, segue que \(\deg{(r_{1}-r_{2})} = \deg{(q_{2}-q_{1})g}\geq \deg{g}.\) Como
      \(\deg{(r_{1}-r_{2})}\leq \max\{\deg{r_{1}}, \deg{r_{2}}\}\leq \deg{g},\) temos um absurdo. Portanto,
      \(q_{2} = q_{1}.\) \qedsymbol
    \end{proof*}
    \begin{example}
      Em \(\mathbb{Z}[x],\) considere \(f(x) = 2x^{3} + x^{2} + x + 1\)  e \(g(x) = -x + 1.\)
      Podemos aplicar o algoritmo da divisão da seguinte forma:
      \begin{align*}
  &2x^{3} + x^{2} + x + 1\quad \text{(Polinômio original)}\\
  &2x^{3} + x^{2} + x + 1 \underbrace{- 2x^{3} + 2x^{2}}_{-2x^{2}g(x)} \quad (-2x^{2}g(x)\text{ pois é o que falta para g ``alcançar'' o grau de f)}\\
  &3x^{2} + x + 1 - 3x^{2} \underbrace{- 3x^{2} + 3x}_{-3xg(x)} \quad (\text{Novamente, faltava multiplicar g(x) por } -3x \text{ a fim de cancelar).}\\
  &4x + 1 \underbrace{- 4x + 4}_{-4g(x)} \quad\text{(Mesma coisa, multiplicamos g(x) para cancelar o } 4x)\\
  &5\quad\text{(Finaliza-se com um polinômio de grau 0 (constante)).}
      \end{align*}
      Coletamos os termos que não foram usados para cancelar os maiores graus, ou seja, \(2x^{2}, 3x\) e \(4\), multiplicamos eles por -1 e, assim, o algoritmo da divisão nos dá 
      \[
        f(x) = g(x)(-2x^{2} -3x -4) + 5 = (-x+1)(-2x^{2}-3x-4)+5.
      \]
    \end{example}
    \newpage

    \section{Aula 11 - 20/09/2023}
    \subsection{Motivações} 
    \begin{itemize}
      \item Corolários do Algoritmo da Divisão;
      \item Decomposição de polinômios.
    \end{itemize}
    \subsection{Corolários do Algoritmo da Divisão}
    \begin{prop*}
      Se F é um corpo, então \(F[x]\) é um Domínio de Ideais Principais.
    \end{prop*}
    \begin{proof*}
      Seja \(\mathfrak{i} \trianglelefteq{F[x]}\) e seja 
      \[
        n = \min\{\deg{f(x)}: f(x)\in \mathfrak{i}\setminus{\{0\}}\}.
      \]
      Se \(n=0,\) então existe um polinômio constante \(f(x) = a\neq 0, a\in F.\) Logo, 
      \[
        1 = \frac{1}{a}\cdot a = \frac{1}{a}f(x)\in \mathfrak{I}.
      \]
      Assim, 
      \[
        I = F[x] = \langle 1 \rangle.
      \]
      Então, podemos assumir que \(n > 0\). Seja \(g(x)\in \mathfrak{i}\) com \(\deg{g} = n\) e observe que
      \(\deg{g}\geq 1\) e \(\ell_{g}\in F^{*}\). Se \(f(x)\in \mathfrak{i},\) pelo algoritmo da divisão, existem
      \(g(x), r(x)\in F[x]\) tais que \(f(x) = q(x)g(x) + r(x),\) com \(\deg{r(x)} < \deg{g(x)}.\) Se \(r(x)\neq0, r(x) = f(x) 
      - q(x)g(x)\in \mathfrak{i}\) e, junto com \(\deg{r} < \deg{g},\) obtemos um absurdo, pois g deveria ser o polinômio de
      grau mínimo.

      Logo, \(r(x) = 0\) e, assim, \(f(x) = q(x)g(x)\) para algum polinômio \(q(x)\in F[x]\), tal que \(f(x)\in \langle g(x) \rangle.\) 
      Portanto, \(\mathfrak{i} = \langle g(x) \rangle.\) \qedsymbol
    \end{proof*}
    \begin{example}[Exercício]
      \begin{itemize}
        \item[1)] \(\mathbb{Z}[x]\) não é D.I.P: O ideal gerado \(\mathfrak{i} = \langle 2, x \rangle\) não é principal, mas é primo.
        \item[2)] Mostre que \(\langle 3, x^{2}-2 \rangle \trianglelefteq{\mathbb{Z}[x]}\) é ideal maximal e não é principal. (Dica:
          mostre que 
          \[
            \frac{\mathbb{Z}[x]}{\langle 3, x^{2}-2 \rangle}\cong{\frac{\mathbb{Z}_{3}[x]}{\langle x^{2}-2 \rangle}}.
          \]
      \end{itemize} 
    \end{example}
    \begin{crl*}
      Seja F um corpo e \(f, g, h\in F[x]\) tais que f(x) é mônico e irredutível.
      Se \(f(x)\mid g(x)h(x),\) então \(f(x)\mid g(x)\) ou \(f(x)\mid h(x).\)
    \end{crl*}
    \begin{proof*}
      Faremos indução no grau de f.

      Se \(\deg{f(x)} = 1,\) como f é mônico, \(f(x) = x-a\) para algum a em F. Por hipótese, existe \(p(x)\) tal que
      \(g(x)h(x) = (x-a)p(x).\) Aplicando em \(x=a,\) segue que \(g(a)h(a) = 0.\) Como f é domínio, é preciso que
      \(g(a) = 0\) ou \(h(a) = 0\). Em outras palavras, \((x-a)\mid g(x)\) ou \((x-a)\mid h(x),\) como desejado.

      Suponha agora que o resultado vale grau menor que n, \(\deg{f} = n\) e \(f\mid g\) e \(f\mid h.\) Assim, existem 
      \(q, r, q', r'\in A[x]\) tais que \(g(x) = f(x)q(x) + r(x)\) e \(h(x) = q'(x)f(x) + r'(x),\) satisfazendo
      \(\deg{r'} < \deg{g'}, \deg{r} < \deg{g}\) e \(r, r'\neq 0.\) Assim,
      \begin{align*}
        g(x)h(x) &= (qf + r)(q'f + r')\\
                 &= qq'f^{2} + (qr' + q'r)f + rr'\\
        \Rightarrow & rr' = gh - qq'f^{2} - (qr' + rq')f.
      \end{align*}
      Sabemos que f divide \(gh, qq'f^{2}\) e \((qr' + q'r)f.\) Com isso, \(f\mid rr'\) e, então,
      existe \(p(x)\in F[x]\) tal que \(fp = rr'.\) Note também que \(\deg{p} < \deg{r} \) e \(\deg{r'}.\) De fato,
      caso contrário, teríamos \(\deg{p}\geq \deg{r}\) e como \(\deg{f} > \deg{r'},\) temos 
      \[
        \deg{p}\geq \deg{r} \Rightarrow \deg{f} + \deg{p} > \deg{r'} + \deg{r} \Rightarrow \deg{fp} > \deg{rr'},
      \]
      o que é um absurdo. O mesmo vale para \(\deg{p}\geq r'.\) Escrevamos \(p = \alpha p_{1}^{\alpha_{1}}\cdot \dotsc \cdot p_{r}^{\alpha_{r}},\) 
      em que \(p_{1}, \cdots, p_{r}\) são irredutíveis (isso é possível pois F é corpo). Assim, \(p_{i}\mid fp = rr'.\) 
      Como \(\deg{p} < \deg{r} < \deg{f}\) e \(\deg{p_{i}}\leq \deg{p},\) podemos aplicar a hipótese de indução em \(p_{i}\) e então
      \(p_{i}\mid r\) ou \(p_{i}\mid r'\) para cada \(i=1, \cdots, r\). Suponha, sem perda de generalidade, que \(p_{1}\mid r,\) tal que 
      \(r=p_{1}r_{1}\). Consequentemente, 
      \[
        fp = rr' \Rightarrow fp_{1}^{\alpha_{1}}\cdot \dotsc \cdot p_{r}^{\alpha_{r}} = p_{1}rr' \Rightarrow fp_{1}^{\alpha_{1} -1}\cdot \cdots \cdot p_{r}^{\alpha_{r}} = r_{1}r.
      \]
      Pode-se continuar esse processo até sobrarem os termos \(p_{i}\)'s, chegando em
      \(f = r_{t}r_{t}'\) com \(\deg{r_{t}}\geq 1\) e \(\deg{r_{t}'}\geq 1\) (pois \(\deg{p} < \deg{r}, \deg{r'}\)). Isso é uma contradição
      com o fato de r ser irredutível.
    \end{proof*}
    \begin{example}[Exercício]
      Seja A um domínio, \(a\in A\) e \((x-a)\mid f(x)g(x).\) Mostre que \((x-a)\mid f(x)\) e \((x-a)\mid g(x).\) 
    \end{example}
    \begin{prop*}
      Seja F um corpo e \(f(x)\in F[x]\). Então, temos uma decomposição única 
      \[
        f(x) = af_{1}^{\alpha_{1}}(x)\cdot \dotsc \cdot f_{r}^{\alpha_{r}}(x),
      \]
      com \(a\in F\) e \(f_{1}, \cdots, f_{r}\in F[x]\) polinômios irredutíveis e mônicos.
    \end{prop*}
    \begin{proof*}
      Já provamos a existência da decomposição para polinômios mônicos. Como F é corpo, existe \(b\in F\) tal que 
      \(bf(x)\) é mônico (basta tomarmos \(b=a^{-1},\) em que a é o coeficiente líder de f). Então, \(bf(x) = f_{1}^{\alpha_{1}}\cdot \dotsc \cdot f_{r}^{\alpha_{r}}\)
      e assim \(f(x) =af_{1}^{\alpha_{1}}\cdot \dotsc \cdot f_{r}^{\alpha_{r}}.\)

      Provemos agora a unicidade. Seja 
      \[
        f(x)=af_{1}^{\alpha_{1}} \cdot \dotsc \cdot f_{r}^{\alpha_{r}} = bg_{1}^{\beta_{1}} \cdot \cdots \cdot g_{s}^{\beta_{s}}.
      \]
      Mostraremos que \(\{f_{1}, \cdots, f_{r}\} = \{g_{1}, \cdots, g_{s}\}\) e \(\{\alpha_{1}, \cdots, \alpha_{r}\} = \{\beta_{1}, \cdots, \beta_{s}\}\).
      Se \(\deg{f} = 1, f = a(x-\alpha ) = b(x-\beta )\) e \(ax -a\alpha =bx - b\beta,\) tal que \(a=b\) e \(\alpha =\beta .\)

      Agora, suponha que o resultado vale para grau menor que n e \(\deg{f} = n.\) Então: 
      \[
        f_{1}\mid f(x) \Rightarrow f_{1}\mid bg_{1}^{\beta_{1}}\cdot \dotsc g_{s}^{\beta_{s}} \Rightarrow \exists j: f_{1}\mid g_{j}.
      \]
      Como \(g_{j}\) é irredutível e mônico, se \(g_{j}=f_{1}\), existe \(h_{j}\) tal que \(g_{j} = h_{j}f_{1}.\) Se \(\deg{(h_{j})}\geq 1, g_{j}\) não seria
      irredutível e, então, \(h_{j} = a\in F.\) Mas, se \(a\neq1,\) como \(f_{1}\) é mônico, \(g_{j}\) não seria mônico. Logo, \(h_{j} = 1\) e \(g_{j} = f_{1}.\) 
      Podemos reordenar tal que \(f_{1} = g_{1}. \) Assim, 
      \[
        af_{2}^{\alpha_{2}}\cdot \dotsc \cdot f_{r}^{\alpha_{r}} = bg_{2}^{\beta_{2}} \cdot \dotsc \cdot g_{s}^{\beta_{s}}.
      \]
      Portanto, pela hipótese indutiva, o resultado segue. \qedsymbol
    \end{proof*}
    \newpage

    \section{Aula 12 - 25/09/2023}
    \subsection{Motivações}
    \begin{itemize}
      \item Corpos de Elementos Finitos;
      \item Característica de um Corpo.
    \end{itemize}
    \subsection{Corpos Finitos}
    Seja F um corpo. Definimos o homomorfismo de anéis
    \begin{align*}
  &\varphi :\mathbb{Z}\rightarrow F\\
  &n\mapsto n \cdot 1_{F}\coloneqq \underbrace{1_{F} + \cdots + 1_{F}}_{\text{n-vezes}}.
    \end{align*}
    Já vimos que \(\overline{\varphi }:\frac{\mathbb{Z}}{\ker{(\varphi )}}\rightarrow F, \overline{n}\mapsto \varphi (n)=n \cdot 1_{F}\) é um monomorfismo. Como F é um
    domínio, \(\mathbb{Z}/\ker{(\varphi )}\) é um domínio. Logo, \(\ker{(\varphi )}\in \mathrm{Spec}(\mathbb{Z}).\) Então, 
    \[
      \ker{(\varphi )} = (0)\quad\text{ou}\quad \ker{(\varphi )}=\langle p \rangle,\text{ p primo.}
    \]
    Olhando mais atentamente a estes casos, se \(\ker{(\varphi )} = (0),\) então \(\varphi \) é um monomorfismo. Agora, 
    o mapa 
    \begin{align*}
   &\overline{\varphi }:\mathbb{Q}\rightarrow F\\
   &\frac{m}{n}\mapsto m \cdot 1_{F} \cdot (n \cdot 1_{F})^{-1} = m \cdot n^{-1}
    \end{align*}
    é um homomorfismo de corpos que é injetivo.

    Por outro lado, se \(\ker{(\varphi )} = \langle p \rangle,\) p um primo, então 
    \[
      \overline{\varphi }:\mathbb{F}_{p}=\mathbb{Z}_{p}\hookrightarrow F,
    \]
    logo \(\overline{\varphi }:\mathbb{F}_{p}\rightarrow F\) é um homomorfismo de corpos que é injetivo.
    \begin{def*}
      Dizemos que um corpo F \textbf{é de característica} 0 se pudermos mergulhar \(\mathbb{Q}\) em F e diremos que \textbf{é de característica positiva} p, p
      um primo, se pudemos mergulhar \(\mathbb{F}_{p} em \mathbb{F}.\) Neste caso, 
      \[
        \mathrm{char}{(F)} = \left\{\begin{array}{ll}
            0,\quad \mathbb{Q}\hookrightarrow F\\
            p,\quad \mathbb{F}_{p}\hookrightarrow \mathbb{F}_{p}.
        \end{array}\right.\quad\square
      \]
    \end{def*}
    Algumas observações devem ser feitas.
    \begin{itemize}
      \item[1)] \(\mathbb{Q}\) e \(\mathbb{F}_{p}\) não têm subcorpos próprios. Por isso, são chamados corpos primos.
      \item[2)] Se F é um corpo finito, então \(\mathrm{char}(F) = p\) por um primo p.
      \item[3)] Se F é um subcorpo do corpo E, então \(\mathrm{char}(E) = \mathrm{char}(F).\)
    \end{itemize}
    \begin{example}
      \begin{itemize}
        \item[1)] \(\mathrm{char}(\mathbb{Q}) = 0,\quad \mathrm{char}(\mathbb{R})=0 = \mathrm{char}(\mathbb{C})\)
        \item[2)] Seja \(\mathbb{Q}(\sqrt[]{2})\coloneqq \{a + b\sqrt[]{2}: a, b\in \mathbb{Q}\}\supseteq{\mathbb{Q}}.\) Segue que 
          \[
            \mathbb{Q}(\sqrt[]{2})\cong{\frac{\mathbb{Q}[x]}{\langle x^{2}-2 \rangle}},
          \]
          em que \(x^{2}-2\) é irracional. Assim, o mapa \(\mathbb{Q}\hookrightarrow \mathbb{Q}[x]/\langle x^{2}-2 \rangle,\quad a\mapsto \overline{a}\) é vazio.
          Com isso, tomando \(f(x)\in \mathbb{Q}[x]\) irracional e colocando \(F = \frac{\mathbb{Q}[x]}{\langle f(x) \rangle},\)
          segue que \(\mathrm{char}(F) = 0.\)
        \item[3)] Coloque \(A = \mathbb{Q}[x]\) e
          \[
            F = Q(A)\coloneqq \biggl\{\frac{f(x)}{g(x)}: f, g\in \mathbb{Q}[x], g(x)\neq0\biggr\}.
          \]
          Segue que \(\mathbb{Q}\subseteq{F}\), tal que \(\mathrm{char}(F) = 0\). Em particular, 
          \[
            E = Q(\mathbb{Z}[x]) = \biggl\{\frac{f(x)}{g(x)}:f, g\in \mathbb{Z}[x], g(x)\neq0\biggr\} = F.
          \]
        \item[4)] Colocando \(K = Q(\mathbb{R}[x]),\) segue que \(\mathbb{Q} \subseteq{\mathbb{R}}\subseteq{K}\) e, logo, \(\mathrm{char}(K) = 0.\)
        \item[5)] Temos \(\mathrm{char}(\mathbb{F}_{p}) = p > 0.\) Além disso, colocando 
          \[
            \mathbb{F}_{p}[\sqrt[]{2}]\coloneqq \frac{\mathbb{F}_{p}[x]}{\langle x^{2}-2 \rangle},\quad p\neq2,
          \]
          se chamarmos de \(F = \mathbb{F}_{3}[\sqrt[]{2}]\) e \(K = \mathbb{F}_{5}[\sqrt[]{5}],\) então
          F é um corpo com \(\mathrm{char}(F) = 3\) e K é um corpo com \(\mathrm{char}(F) = 5.\) No entanto, para p = 7,
          \(x^{2}-2\) tem uma raiz em \(\mathbb{F}_{7}\) e, logo, não é irracional. Portanto, \(\mathbb{F}_{7}[\sqrt[]{2}]\) não é um
          corpo.
        \item[6)] Defina 
          \[
            \mathbb{Q}(\mathbb{F}_{p}[x])\coloneqq \biggl\{\frac{f}{g}: f, g\in \mathbb{F}_{p}[x], g\neq 0\biggr\}.
          \]
          Então, \(\mathbb{F}_{p}[x]\) é infinito, logo F é infinito. Apesar disso, \(\mathrm{char}(F) = p > 0.\)
      \end{itemize}
    \end{example}
    Seja K um corpo finito tal que \(\mathrm{char}(K) = p > 0.\) Então, \(\mathbb{F}_{p}\hookrightarrow K\) e 
    \[
      n = [K:\mathbb{F}_{p}] = \dim_{\mathbb{F}_{p}}K < \infty.
    \]
    Isto implica que \(|K| = p^{n}\). Como \(\mathbb{F}_{p}\)-espaço vetorial, temos 
    \[
      K\cong{\mathbb{F}_{p}^{n}} = \underbrace{\mathbb{F}_{p}\times \cdots\times \mathbb{F}_{p}}_{\text{n-vezes}}
    \]
    Como um fato geral, se F é um corpo e olhamos para V como um F-espaço vetorial com dimensão \(n = \dim_{F}V < \infty\), então 
    \[
      V\cong{F^{n}}.
    \]
    \begin{theorem*}
      \begin{itemize}
        \item[1)] Para todo primo p e todo número natural \(n > 0,\) temos um corpo com \(p^{n}\) elementos;
        \item[2)] Quaisquer dois corpos com \(p^{n}\) elementos são isomorfos.
      \end{itemize}
    \end{theorem*}
    \begin{def*}
      Denotamos um corpo finito com \(p^{n}\) elementos por \(\mathbb{F}_{p^{n}}.\square\)
    \end{def*}
    \begin{theorem*}
      Seja p um primo e \(m, n\in \mathbb{N}\) com \(m\leq n.\) Então, podemos mergulhar \(\mathbb{F}_{p^{m}}\) em \(\mathbb{F}_{p^{n}}\) se, e somente se, \(m\mid n.\)
    \end{theorem*}
    \newpage

    \section{Aula 13 - 27/09/2023}
    \subsection{Motivações}
    \begin{itemize}
      \item Espectro e Maximais de um Corpo de Polinômios;
      \item Extensões e Quocientes.
    \end{itemize}
    \subsection{Espectro de Polinômios}
    \begin{prop*}
      Se F é um corpo, então o espectro de \(F[x]\) é o conjunto \(\mathrm{Spec}(F[x]) = \{0\}\cup\{\langle f(x) \rangle: f\text{ é irredutível}\}\) e o espectro maximal de 
      \(F[x]\) é \(\mathrm{Specm}(F[x]) = \{\langle f(x) \rangle: f\text{ é irredutível}\}\).
    \end{prop*}
    \begin{proof*}
      Como \(F[x]\) é um domínio, \((0)\in \mathrm{Spec(F[x])}.\) Seja \( \mathfrak{p} = \langle f(x) \rangle\in \mathrm{Spec}(F[x])\) e seja
      f(x) não irredutível. Então, \(f(x) = g(x)h(x),\) em que \(\deg{g}, \deg{h} < \deg{f}.\) Temos 
      \[
        gh = f\in \langle f(x) \rangle = \mathfrak{p} \Rightarrow g\in \mathfrak{p}\text{ ou } h\in \mathfrak{p}.
      \]
      Contradição. Logo, f(x) é irredutível. Agora, suponhamos que \(f(x)\) é irredutível, de forma que nosso objetivo será mostrar que 
      \(\langle f(x) \rangle\) é primo. Sejam \(g, h \in \langle f(x) \rangle,\) tal que \(f\mid gh.\) Como f é irredutível, vale que 
      ou \(f\mid g\), ou \(f\mid h\), ou seja, \(g\in \langle f(x) \rangle\) ou \(h\in \langle f \rangle\), o que significa, exatamente,
      que \(\langle f \rangle\) é primo, isto é, \(\langle f \rangle\in \mathrm{Spec}(F[x]).\)

      Agora, tome \(\mathfrak{m}\in \mathrm{Specm}(F[x]).\) Duas possibilidades surgem - ou \(\mathfrak{m}\) é \(\langle f(x) \rangle\) com f irredutível,
      ou \((0)\), ou seja, \(\mathrm{Specm}(F[x]) \subseteq{\{\langle f(x) \rangle: f\text{ é irredutível}\}}\). Por outro lado,
      seja f(x) irredutível. Pela primeira parte, \(\langle f(x) \rangle\) é primo. Considere 
      \[
        \langle f(x) \rangle \subseteq{\mathfrak{j}}\subsetneq{F[x]}
      \]
      e coloque \(\mathfrak{j} = \langle g(x) \rangle.\) Temos 
      \[
        \langle f(x) \rangle \subseteq{\langle g(x) \rangle} \Rightarrow f(x)\in \langle g(x) \rangle,
      \]
      ou seja, \(g\mid f,\) mas f é irredutível, donde segue que \(g=1\) ou \(g = f.\) Se \(g=1,\) então \(\mathfrak{j} = F[x],\) 
      uma contradição. Logo, \(g=f\) e \(\langle f \rangle = \mathfrak{j}.\) Portanto, \(\langle f \rangle\) é maximal, provando que 
      \(\mathrm{Specm}(F[x]) = \{\langle f(x) \rangle: f\text{ é irredutível}\}\). \qedsymbol
    \end{proof*}
    \subsection{Corpos}
    \begin{def*}
      Sejam F e K dois corpos tais que \(F\subseteq{K}.\) Neste caso, dizemos que K é uma \textbf{extensão} de F e denotaremos por
      \(K/F.\) Se \(K/F\) é uma extensão de corpos, K pode ser visto como espaço vetorial sobre F, e a dimensão de K
      como F-espaço vetorial (\(\dim_{F}K\) é denotada por \([K:F]\) e chamada o \textbf{grau} da extensão. \(\square\)
    \end{def*}
    \begin{example}
      \begin{itemize}
        \item[1)] \([\mathbb{C}:\mathbb{R}] = 2,\) com base \(\{1, i\}\);
        \item[2)] \([\mathbb{R}:\mathbb{Q}] = \infty\) (não enumerável);
        \item[3)] Se F é um corpo, \([F:f]=1\);
        \item[4)] \([\mathbb{Q}[\sqrt[]{2}]:\mathbb{Q}] = 2,\) com base \(\{1, \sqrt[]{p}\}\).
      \end{itemize}
    \end{example}
    \begin{crl*}
      Se \(f(x)\in F[x]\) é irredutível e mônico, então \(K\coloneqq \frac{F[x]}{\langle f(x) \rangle}\) é um corpo.
    \end{crl*}
    \begin{proof*}
      Segue automaticamente de \(\langle f(x) \rangle \) ser maximal e o quociente de um corpo por um maximal ser outro corpo. \qedsymbol
    \end{proof*}
    Podemos definir o mapa 
    \begin{align*}
   &\varphi :F\rightarrow K=\frac{F[x]}{\langle f(x) \rangle}\\
   &a\mapsto \overline{a} = a + \langle f(x) \rangle.
    \end{align*}
    Observe que \(\varphi \) é um morfismo de corpos injetor, já que \(\ker{\varphi } \trianglelefteq{F}\) e,
    sendo F um corpo, os únicos ideais são \((0)\) ou \(F\). Sabendo que \(\varphi(1) = \overline{1}\neq\overline{0}, \ker{\varphi }\neq F\) e,
    assim, \(\ker{\varphi } = (0).\)

    Com isso, podemos pensar em F como subcorpo de K usando a injeção dada por \(\varphi \), e denotaremos \(\overline{a}\) simplesmente por a.
    \begin{prop*}
      Seja \(f(x)\in F[x]\) irredutível e mônico e \(K = \frac{F[x]}{\langle f(x) \rangle}.\) Então, \([K:F] = \deg{f(x)}.\)
    \end{prop*}
    \begin{proof*}
      Seja \(\alpha  = \overline{x}\), tal que \(K = F(\alpha ) = \{g(\alpha ): g(x)\in F[x]\}.\) Tome \(\beta \in K.\) Então,
      \(\beta = \overline{g(x)} = g(\overline{x}) = g(\alpha ).\) Mostraremos que \(B = \{1, \alpha , \cdots, \alpha ^{n-1},\}\) em que
      \(n=\deg{f}\), é base de \(K/F.\)

      Começamos provando que B gera todo o espaço. Pelo algoritmo da divisão, \(g(x) = q(x)f(x) + r(x),\) com \(\deg{r} < \deg{g}.\)
      Assim, \(\beta  = g(\alpha ) = \underbrace{\overline{q(x)f(x)}}_{=\overline{0}} + \overline{r(x)} = \overline{r(x)}\).
      Então, se \(r(x) = b_{m}x^{m} + \cdots + b_{0}, m < n,\) segue que 
      \[
        \beta = 0\alpha ^{n-1} + \cdots + 0\alpha ^{m+1} + b_{m}\alpha ^{m} + \cdots + b_{0}.
      \] 

      Agora, vamos mostrar que B é linearmente independente. Seja \(c_{0}1 + c_{1}\alpha + \cdots + c_{n-1}\alpha ^{n-1} = 0\) com coeficientes
      \(c_{i}\in F.\) Se \(h(x) = c_{n-1}x^{n-1} + \cdots + c_{1}x + c_{0},\) então \(\overline{h(x)} = \overline{0}\in K\) e, assim, 
      \(h(x)\in \langle f(x) \rangle,\) ou seja, \(f(x)\mid h(x).\) Como \(\deg{h} < \deg{f}, \) isso é possível se, e somente se, \(h(x) = 0,\) o
      que implica que \(c_{i} = 0\) para todo \(i=0, \cdots, n-1\).

      Portanto, B é um conjunto gerador linearmente independente de K, ou seja, uma base. \qedsymbol
    \end{proof*}
    \begin{example}
      \begin{itemize}
        \item[1)] Como \([\mathbb{C}:\mathbb{R}]=2, \mathbb{C}\cong{\frac{\mathbb{R}}{\langle x^{2} + 1 \rangle}}\)
        \item[2)] Seja F um corpo qualquer. Como \([F:F] = 1, F\cong{\frac{F[x]}{\langle x-a \rangle}}\)
        \item[3)] Para \(\mathbb{Q}[\sqrt[]{p}],\) tal que \([\mathbb{Q}[\sqrt[]{p}]:\mathbb{Q}]=2, \mathbb{Q}[\sqrt[]{p}]\cong{\mathbb{Q}/\langle x^{2}-p \rangle}\)
        \item[4)] Seja \(f(x) = x^{2} + x + 1\in \mathbb{F}_{2}[x].\) Note que \(f(0) = 1\neq0\) e \(f(1) = 1^{2} + 1 + 1 = 1\neq0\), ou seja,
          f não tem raízes em \(\mathbb{F}_{2}\), o que significa que ele é irredutível. Assim, \(K = \mathbb{F}_{2}[x]/\langle x^{2}+x+1 \rangle = \mathbb{F}_{2}(\alpha )\) 
          e, pelo resultado acima, \([K:\mathbb{F}_{2}] = \deg{f} = 2,\) tal que \(K = \{\overline{0}, \overline{1}, \overline{x}, \overline{x+1}: \overline{x}^{2} = \overline{x} + \overline{1}\}\)
          e \(|K| = 4.\) Construímos, assim, um corpo com 4 elementos.
        \item[5)] Se \(K = \frac{\mathbb{F}_{2}[x]}{\langle x^{3} + x + 1 \rangle}\) e \(L = \mathbb{F}_{2}[x]/\langle x^{3}+x^{2}+1 \rangle,\) então \([K:\mathbb{F}_{2}] = 3, [L:\mathbb{F}_{2}]=3\).
          Assim, \(|K| = 2^{3} = 8\) e \(|L| = 2^{3} = 8\).
      \end{itemize}
    \end{example}
    \begin{example}[Exercícios]
      \begin{itemize}
        \item[1)] Se E é um corpo com 4 elementos, mostre que \(E\cong{\frac{\mathbb{F}_{2}[x]}{\langle x^{2}+x+1 \rangle}}\)
        \item[2)] Se E é um corpo com 8 elementos, então \(E\cong{K}\cong{L}.\) Denotaremos um corpo com 8 elementos por \(\mathbb{F}_{8}.\)
        \item[3)] Mostre que não existe monomorfismo \(\mathbb{F}_{4}\hookrightarrow \mathbb{F}_{8}.\)
        \item[4)] Seja F corpo e \(f(x)\in F[x]\) tal que \(2\leq \deg{f}\leq 3\). Mostre que f é irredutível se, e somente se, f noa possui raízes em F.
        \item[5)] Mostre que \(\mathbb{F}_{4}^{*}\cong{\mathbb{Z}/3}, \mathbb{F}_{8}^{*}\)
      \end{itemize}
    \end{example}
    \newpage

    \section{Aula 14 - 09/10/2023}
    \subsection{Motivações}
    \begin{itemize}
      \item Números Algébricos e Transcendentes;
      \item Extensões Algébricas;
    \end{itemize}
    \subsection{Números Algébricos e Transcendentes}
    \begin{def*}
      Seja \(K/F\) uma extensão de corpos e seja \(a\in K\). Definimos o \textbf{mapa de evaluação em} \(\alpha \) por:
      \begin{align*}
   &e_{\alpha }:F[x]\rightarrow K\\
   &f(x)\mapsto f(\alpha ).\quad\square
      \end{align*}
      Note que \(e_{\alpha }\) é um morfismo de anéis. Dizemos que um elemento \(\alpha \in K\) é
      \begin{itemize}
        \item[1)] \textbf{algébrico} sobre F se \(\ker{(e_{\alpha })}\neq(0)\)
        \item[2)] \textbf{transcendente} sobre F se \(\ker{(e_{\alpha })}=(0).\)
      \end{itemize}
    \end{def*}
    \begin{example}[Exercícios]
      \begin{itemize}
        \item[1)] Mostre que \(\alpha \in K\) é algébrico sobre F se, e somente se, existe \(f(x)\neq 0\) em \(F[x]\) tal que
          \(f(\alpha ) = 0.\)
        \item[2)] Mostre que \(\alpha \in K\) é transcendente sobre F se, e somente se, para todo
          \(f(x)\in F[x]\) não nulo, \(f(\alpha)\neq0.\)
      \end{itemize}
    \end{example}
    Podemos observar algumas coisas. A primeira delas é que \(\ker{e_{\alpha }}\neq(0).\) Como
    F[x] é D.I.P., existe f(x) tal que \(\ker{e_{\alpha }} = \langle f(x) \rangle.\) Já que \((0)\in \mathrm{Spec(K)}, \ker{e_{\alpha }}=e_{\alpha }^{-1}((0))\in \mathrm{Spec(F[x])}.\)
    Assim, \(f(x)\) é irredutível e \(\ker{e_{a}}\in \mathrm{Specm(F[x])},\) donde obtemos a injeção de corpos:
    \[
      \overline{e}_{\alpha }:\frac{F[x]}{\langle f(x) \rangle}\rightarrow K.
    \]
    Denotaremos \(Im(e_{\alpha }) = F(\alpha ).\) Desta forma, \(F(\alpha )\subseteq{K}\) e \([F(\alpha ): F]=\deg{f}\) pela última aula.

    Além disso, note que se \(e_{\alpha }\) é injetor, então \(F[x]\cong{F[\alpha ]}\) (isomorfismo de anéis).
    \begin{def*}
      Dizemos que uma extensão \(K/F\) é \textbf{algébrica} se todo elemento de K é algébrico. \(\square\)
    \end{def*}
    \begin{example}
      \begin{itemize}
        \item[1)] Se \([K:F]\) é finita, então \(K/F\) é algébrica. De fato, seja \(n=[K:F]\) e \(\alpha \in K\) um elemento qualquer.
          Considere o conjunto \(S = \{1, \cdots, \alpha ^{n}\}\subseteq{K};\) Como \(\dim_{F}K = n,\) S é necessariamente um conjunto linearmente
          dependente. Então, existe \(a_{0}, \cdots, a_{n}\in F\) não todos nulos tais que \(a_{0} + a_{1}\alpha + \cdots + a_{n}\alpha^{n} = 0.\)

          Considere o polinômio \(f(x) = a_{0}+a_{1}x + \cdots + a_{n}x^{n}.\) Ele é não nulo, pois nem todos os \(a_{i}\) são nulos,
          e \(f(\alpha ) = 0\). Logo, \(\alpha \) é algébrico sobre F.

        \item[2)] Seja \(K/\mathbb{C}\) uma extensão algébrica sobre \(\mathbb{C}.\) Então, \(K = \mathbb{C}.\) Com efeito, seja \(\alpha \in K\)
          e \(f(x)\in \mathbb{C}[x]\) não nulo e mônico tal que \(f(\alpha ) = 0\). Pelo Teorema Fundamental da Álgebra, podemos escrever 
          \[
            f(x) = a(x-\alpha_{1})\cdot \dotsc \cdot (x-\alpha_{n}),
          \]
          para algum \(a\in \mathbb{C}\setminus{\{0\}},\alpha_{i}\in \mathbb{C}\) e \(n = \deg{f}.\) Assim, \(0 = f(\alpha ) = a(\alpha -\alpha_{1})\cdot \dotsc(\alpha -\alpha_{n})\)
          e, como K é domínio, existe i tal que \((\alpha -\alpha_{i})=0,\) o que significa que \(\alpha=\alpha_{i}\in \mathbb{C}.\) Logo, \(\ker{e_{\alpha }}=\langle x-a \rangle\) e \([K:\mathbb{C}] = 1.\)
      \end{itemize}
    \end{example}
    \begin{prop*}
      Seja \(K/F\) uma extensão de corpos. Suponha que F é não enumerável e \([K:F]=\dim_{F}K\) é enumerável. Então,
      \(K/F\) é algébrica.
    \end{prop*}
    \begin{proof*}
      Suponha que \(K/F\) não seja algébrica. Escolha \(\alpha \in K/F\) que não é
      algébrico sobre F. Mostraremos que 
      \[
        S = \biggl\{\frac{1}{\alpha - a}: a\in F\biggr\}\subseteq{K}
      \]
      é linearmente independente, e assim teremos um absurdo, pois \(|S| = |F|\) e teríamos 
      \(\dim_{F}K\geq |S|.\)

      Com efeito, seja uma combinação linear finita:
      \begin{align*}
  &\frac{c_{1}}{\alpha - a_{1}} + \cdots + \frac{c_{n}}{\alpha - a_{n}} = 0\\
        \underbrace{\Rightarrow}_{\times (\alpha -a_{1})}&c_{1} + \frac{c_{2}(\alpha - a_{1})}{\alpha - a_{2}} + \cdots + \frac{c_{n}(\alpha - a_{1})}{\alpha - a_{n}} = 0.
      \end{align*}
      Como \(\alpha \) é transcendente, temos \(F[x]\cong{F(\alpha )}\) e então \(\mathrm{Frac}(F[x])\cong{\mathrm{Frac}(F(\alpha ))}\).
      Assim, aplicando o isomorfismo, substituindo \(\alpha \) por x, obtemos a mesma igualdade com 0: 
      \[
        c_{1} + \frac{c_{2}(x-a_{1})}{x-a_{2}} + \cdots + \frac{c_{n}(x-a_{1})}{x-a_{n}} = 0.
      \]
      Como a igualdade vale para todo x, substituindo \(x=a_{1}, c_{1} = 0.\) Podemos fazer esse processo para
      todo \(i=1, \cdots, n\) e conseguindo, assim, \(c_{i} = 0\) para \(i=1, \cdots, n.\) Então, o conjunto é linearmente independente
      e temos o que queríamos demonstrar. \qedsymbol
    \end{proof*}
    \begin{theorem*}
      Seja \(A = F[x_{1}, \cdots, x_{n}],\) F corpo. Então, \(\mathfrak{i} = \langle x_{1}-a_{1}, \cdots, x_{n}-a_{n}\rangle \trianglelefteq{A}\) 
      é maximal, para \(a_{1}, \cdots, a_{n}\in F\).
    \end{theorem*}
    \begin{proof*}

      A prova segue da indução sobre o número de variáveis

      \textbf{\underline{Caso Base}:} Sejam \(\mathfrak{i} = \langle x - a_{1} \rangle, 0\neq \overline{f(x)}\in A/\mathfrak{i}.\)
      Existem \(q(x), r(x)\) tais que \(f(x) = q(x)(x-a_{1})+r(x),\) com \(\deg{r}\leq \deg{(x-a)}=1.\) Com isso,
      r(x) é constante, digamos \(r(x) = a\). Então, \(\overline{f(x)}=\overline{a}\) é constante e
      \(g(x) = \overline{a^{-1}}\) é o inverso de f(x) em \(A/\mathfrak{i}.\) Assim, \(A/\mathfrak{i}\) é um corpo isomorfo a F por meio
      de \(f(x)\mapsto a.\)

      \textbf{\underline{Hipótese Indutiva}:} Suponha que o resultado vale para \(k < n\). Agora, \(\mathfrak{i} = \langle x-a_{1}, x-a_{2}, \cdots, x-a_{n} \rangle\) e tomemos
      \(0\neq \overline{f(x_{1}, \cdots, x_{n})}\in A/\mathfrak{i}\). Podemos escrever
      \(f(x_{1}, \cdots, x_{n})\) como um elemento de \(F[x_{2}, \cdots, x_{n}][x_{1}]\) da seguinte forma:
      \[
        f(x_{1}, \cdots, f_{n}) = f_{m}x_{1}^{m} + f_{m-1}x_{1}^{m-1} + \cdots + f_{0},
      \]
      em que \(f_{i}\in F[x_{2}, \cdots, x_{n}].\) Pelo algoritmo de Euclides, existem \(q, r\in A\) com 
      \[
        f(x_{1}, \cdots, x_{n}) = q(x_{1} - a_{1} ) + r,
      \]
      com \(r\neq0\) e \(\deg_{x_{1}}{r} = 0, r \in F[x_{2}, \cdots, x_{n}].\) Por indução, \(\frac{F[x_{2}, \cdots, x_{n}]}{\langle x_{2}-a_{2}, \cdots, x_{n}-a_{n} \rangle}\cong{F}\) e
      então \(\overline{r} = \overline{b}\in F, r = b+g, g\in \langle x_{2}-a_{2}, \cdots, x_{n}-a_{n} \rangle\). Logo, 
      \[
        f(x_{1}, \cdots, x_{n}) = q(x_{1}-a_{1}) + g + b \Rightarrow \overline{f(x_{1}, \cdots, x_{n})} = \overline{b}.
      \]
      Assim, \(\overline{f(x_{1}, \cdots, x_{n})}\) corresponde a um polinômio constante \(\overline{b}\in F.\)
      Portanto, \(\mathfrak{i}\) é maximal. \qedsymbol
    \end{proof*}
    \newpage

    \section{Aula 15 - 18/10/2023}
    \subsection{Motivações}
    \begin{itemize}
      \item Nullstelensatz de Hilbert;
      \item Fecho Algébrico.
    \end{itemize}
    \subsection{Nullstelensatzs}
    \begin{theorem*}[Nullstelensatz Algébrico]
      Seja K um corpo algebricamente fechado. Então, 
      \[
        \mathrm{Specm}(K[x_{1}, \cdots, x_{n}]) = \{\langle x-a_{1}, \cdots, x-a_{n} \rangle: a_{1}, \cdots, a_{n}\in K\}
      \] 
    \end{theorem*}
    \begin{proof*}
      Vamos provar para o caso \(K = \mathbb{C}.\) Tome \(M\in \mathrm{Specm}(K[x_{1}, \cdots, x_{n}])\) e considere o quociente \(F \coloneq K[x_{1}, \cdots, x_{n}]/M\). Como M é maximal
      este quociente é um corpo. É claro que \(F\hookrightarrow K\) e que \(S=\{\overline{x_{1}}^{i_{1}}\cdot \dotsc \cdot \overline{x_{n}}^{i_{n}}: i_{j}\in \mathbb{N}\} \subseteq{F}\)
      gera F como K-espaço vetorial e é enumerável. Em outras palavras, \(\dim_{K}F\) é enumerável e, como K é não enumerável, a extensão
      \(K/F\) é algébrico, do que segue que \(K\cong{F}.\)

      Seja \(\overline{x_{i}} = \overline{a_{i}}, a_{i}\in \mathbb{C}\). Então, \(\overline{x_{i} - a_{i}} = \overline{0}\) e, assim, \(x_{i}-a_{i}\in M\) para cada \(i=1, \cdots, n.\)
      Sabemos que \(\langle x_{1}-a_{1}, \cdots, x_{n}-a_{n} \rangle\) é maximal e \(\langle x_{1}-a_{1}, \cdots, x_{n}-a_{n} \rangle \subseteq{M}\) implica, portanto, que
      \(M = \langle x_{1}-a_{1}, \cdots, x_{n}-a_{n} \rangle\). \qedsymbol
    \end{proof*}
    \begin{theorem*}[Nullstelensatz Geométrico]
      Seja \(S \subseteq{K[x_{1}, \cdots, x_{n}]}\) tal que \(\langle S \rangle\neq K[x_{1}, \cdots, x_{n}].\) Então, existe \(a=(a_{1}, \cdots, a_{n})\in K^{n}\)
      tal que, para todo \(f(x_{1}, \cdots, x_{n})\in K[x_{1}, \cdots, x_{n}], f(a_{1}, \cdots, a_{n}) = 0.\)
    \end{theorem*}
    \begin{proof*}
      Tome \(M\in \mathrm{Specm}(K[x_{1}, \cdots, x_{n}])\) tal que \(\langle S \rangle \subseteq{M}.\) Pelo Teorema
      Nullstelensatz Algébrico, sabemos que \(M = \langle x_{1}-a_{1}, \cdots, x_{n}-a_{n} \rangle\), para \(a_{i}\in \mathbb{C}.\)

      Pelo isomorfismo \(\varphi :K[x_{1}, \cdots, x_{n}]/M\rightarrow \mathbb{C},\) definido por \(\varphi (\overline{x_{i}}) = a_{i}\) temos
      um polinômio genérico \(\varphi (\overline{g(x_{1}, \cdots, x_{n})}) = g(a_{1}, \cdots, a_{n}).\) Como \(\langle S \rangle \subseteq{M},\)
      dado \(f(x_{1}, \cdots, x_{n})\in S, \varphi (\overline{f(x_{1}, \cdots, x_{n})})=\varphi (\overline{0}) = 0\) e, por outro, \(\varphi (\overline{f(x_{1}, \cdots, x_{n})}) =
      f(a_{1}, \cdots, a_{n}).\) Assim, o ponto \((a_{1}, \cdots, a_{n})\) é raiz de todo polinômio \(f\in M\). Portanto,
      é raiz de \(f\in S.\) \qedsymbol
    \end{proof*}
    \subsection{Fecho Algébrico}
    \begin{def*}
      Um corpo K é dito \textbf{algebricamente fechado} se todo polinômio não constante \(f(x)\in K[x]\) tem
      uma raiz em K. \(\square\)
    \end{def*}
    \begin{lemma*}
      Se K é algebricamente fechado, então todo \(f(x)\in K[x]\) fatora-se na forma 
      \[
        f(x) = (x-a_{1})^{r_{1}}\cdot \dotsc \cdot (x-a_{n})^{r_{n}},
      \]
      em que \(a_{1}, \cdots, a_{n}\in K\).
    \end{lemma*}
    \begin{example}
      \begin{itemize}
        \item[1)] \(\mathbb{C}\) é algebricamente fechado pelo Teorema Fundamental da Álgebra;
        \item[2)] \(\overline{\mathbb{Q}} = \{\alpha \in \mathbb{C}: f(\alpha ) = 0, f(x)\in \mathbb{Q}[x]\}\) é algebricamente fechado
      \end{itemize}
    \end{example}
    Todo corpo pode ser mergulhado em um outro algebricamente fechado.
    \begin{def*}
      Seja F um corpo e e K um corpo algebricamente fechado que é uma extensão de F. Chamamos de
      \textbf{fecho algébrico de F} o conjunto 
      \[
        \overline{F} = \{\alpha \in K: \exists f(x)\in F[x], f(\alpha ) = 0\}.\quad\square
      \]
    \end{def*}
    \begin{prop*}
      Seja F um corpo e K um corpo algebricamente fechado tal que \(F \subseteq{K}.\) Então, 
      \[
        \overline{F} = \{\alpha \in K: \exists f(x)\in F[x], f(\alpha ) = 0\}
      \]
      é um corpo algebricamente fechado.
    \end{prop*}
    Vale notar que o fecho algébrico de um corpo é único a menos de isomorfismo.
    \subsection{Parte Extra - Construindo a Topologia de Zariski sobre \(\mathbb{C}\)}
    Para \(\mathfrak{i}\trianglelefteq{\mathbb{C}[x_{1}, \cdots, x_{n}]},\) defina \(V(I) = \{a\in \mathbb{C}^{n}: f(a) = 0 \forall f\in \mathfrak{i}\}\). 
    Mostre que:
    \begin{itemize}
      \item[1)] \(\bigcap_{k\in K}^{}{V(\mathfrak{i}_{k})} = V(\sum\limits_{k\in K}^{}\mathfrak{i}_{k})\);
      \item[2)] \(\bigcup_{i=1}^{n}{V(\mathfrak{i_{k}})} = V(\bigcap_{i=1}^{n}{\mathfrak{i}_{k}}) = V(\mathfrak{i}_{1}\cdot \dotsc \mathfrak{i}_{k})\);
      \item[3)] \(V(\mathbb{C}[x_{1}, \cdots, x_{n}]) = \emptyset\) e \(V((0)) = \mathbb{C}^{n}\);
      \item[4)] \(\mathfrak{i}\subseteq{\mathfrak{j}} \Rightarrow V(\mathfrak{j}) \subseteq{} V(\mathfrak{i})\);
      \item[5)] \(V(\mathfrak{i}) = V(\sqrt[]{\mathfrak{i}}),\) em que \(\sqrt[]{\mathfrak{i}}\coloneqq \{g\in A: \exists n, g^{n}\in \mathfrak{i}\}\)
    \end{itemize}
    As propriedade 1, 2 e 3 fornece-nos uma topologia sobre \(\mathbb{C}^{n}\) gerada pelos fechados \(V(\mathfrak{i}),\) chamada \textbf{topologia de Zariski.} Se
    \(f\in \mathbb{C}[x_{1}, \cdots, x_{n}],\) definimos \(V(f) = V(\langle f \rangle)\).
    \begin{itemize}
      \item[1)] Mostre que \(S^{1}\coloneqq \{z=(a, b)\in \mathbb{C}^{2}: |a|^{2} + |b|^{2}=1\}\) é fechado na topologia usual em
        \(\mathbb{C}^{2},\) mas não é fechado na topologia de Zariski.
      \item[2)] Mostre que \(\{V(f): f\in \mathbb{C}[x_{1}, \cdots, x_{n}]\}\) é uma base para a topologia de Zariski em \(\mathbb{C}^{n}.\)
      \item[3)] Mostre que \(\mathbb{C}^{n}\) é compacto com a topologia de Zariski.
    \end{itemize}
    \newpage

    \section{Aula 16 - 23/10/2023}
    \subsection{Motivações} 
    \begin{itemize}
      \item Domínios Euclideanos;
      \item Normas.
    \end{itemize}
    \subsection{Domínios Euclidianos}
    \begin{def*}
      Um domínio A é dito \textbf{euclidiano} se existe uma função \(N:A\setminus{\{0\}}\rightarrow \mathbb{N}\cup\{0\}\)
      satisfazendo:
      \begin{itemize}
        \item[1)] \(N(a)\leq N(ab)\) para todos \(a, b\in A\setminus{\{0\}}\)
        \item[2)] (Algoritmo da Divisão) Para todos \(a, b\in A, b\neq0,\) existem \(q,r\in A\) tais que
          \(a = bq + r,\) em que \(N(r) < N(b)\) ou \(r = 0\).
      \end{itemize}
      Neste caso, a função N é  chamada \textbf{norma} de A. \(\square\)
    \end{def*}
    \begin{example}
      \begin{itemize}
        \item[1)] Para \(A = \mathbb{Z},\) uma possível norma é 
          \begin{align*}
      &N=|\cdot |:\mathbb{Z}\rightarrow \mathbb{N}\\
      &n\mapsto |n|
          \end{align*}
        \item[2)] Se F é um corpo e \(A = F[x]\), temos a norma 
          \begin{align*}
    &N:F[x]\rightarrow \mathbb{N}\\
    &p(x)\mapsto \deg{p(x)}.
          \end{align*}
      \end{itemize}
    \end{example}
    \begin{theorem*}
      Domínios Euclidianos são D.I.P.'s
    \end{theorem*}
    \begin{proof*}
      Seja \(\mathfrak{i} \trianglelefteq{A}\) não nulo e seja \(0\neq a\in \mathfrak{i}\) elemento com
      menor norma, i.e., \(N(A) = \min\{N(x): x\in \mathfrak{i}\}\). Mostraremos que \(\mathfrak{i} = \langle a \rangle\).

      Por um lado, é claro que \(\langle a \rangle \subseteq{\mathfrak{i}}.\) Por outro, tome \(x\in \mathfrak{i}\).
      Pela segunda parte da definição, existem \(q, r\in A\) tais que \( x = qa + r\) e \(N(r) < N(a).\)
      Então, \(r=x-qa\in \mathfrak{i}.\) Se \(r\neq0, N(r) < N(a)\), o que é um absurdo pela definição de a como 
      elemento com menor norma. Assim, \(r=0\), ou seja, \(x = qa\). Portanto, \(x\in \mathfrak{i}\) e \(\mathfrak{i} = \langle a \rangle\). \qedsymbol
    \end{proof*}
    \begin{example}
      \begin{itemize}
        \item[1)] Corpos são Domínios Euclidianos:
          \begin{align*}
    &N:F^{*}\rightarrow \mathbb{N}\\
    &a\mapsto 1
          \end{align*}
        \item[2)] Em \(A = K[x, y],\) o ideal \(\mathfrak{i} = \langle x, y \rangle\) não é principal. Logo, A não é Domínio Euclidiano.
      \end{itemize}
    \end{example}
    Vale observar que a norma num Domínio Euclidiano não é única, no sentido de que pode ter mais que uma norma. Por exemplo,
    \[
      \left\{\begin{array}{ll}
          N_{1}:&\mathbb{Z}\setminus{\{0\}}\rightarrow \mathbb{N}\\
                &a\mapsto |a|\\
          N_{2}:&\mathbb{Z}\setminus{0}\rightarrow \mathbb{N}\\
                &a\mapsto |2a|.
      \end{array}\right.
    \]
    \begin{prop*}
      Seja \((A, N)\) um Domínio Euclidiano. Então, \(A^{*}=\{a\in A\setminus{\{0\}}: N(a) = N(1)\}.\)
    \end{prop*}
    \begin{proof*}
      Sejam \(a\in A^{*}\) e \(b\in A^{*}\) tais que \(ab = 1.\) Pela propriedade 1,
      \(N(a)\leq N(ab) = N(1).\) Além disso, vale também que \(N(1)\leq N(a \cdot 1) = N(a).\) Logo,
      \(N(a) = N(1).\) Com isso, provamos que \(A^{*}\subseteq{\{a\in A\setminus{\{0\}}: N(a) = N(1)\}}.\)

      Por outro lado, tome um elemento \(a\in A\setminus{\{0\}}\) tal que \(N(a) = N(1).\) Aplicando o algoritmo da divisão,
      existem \(q, r\in A\) tais que \(1 = aq + r, r = 0\) ou \(N(r) < N(a) = N(1).\)

      Se \(r\neq0, N(1)\leq N(1 \cdot r) = N(r)\) implica que \(N(1)\leq N(r).\) No entanto, por hipótese,
      \(N(r) < N(1)\), ou seja, temos um absurdo. Logo, \(r = 0\) e \(aq = 1\) implica que \(a\in A^{*},\) assim provando
      a outra relação de contenção.

      Portanto, \(A^{*}=\{a\in A\setminus{\{0\}}: N(a) = N(1)\}.\) \qedsymbol
    \end{proof*}
    \begin{theorem*}
      O anel \(\mathbb{Z}[i] = \{a + ib: a, b\in \mathbb{Z}\}\subseteq{\mathbb{C}}\) é um Domínio Euclidiano com a norma
      \begin{align*}
   &N:\mathbb{Z}[i]\setminus{\{0\}}\rightarrow \mathbb{N}\\
   &z = a + ib\mapsto |z|^{2} = a^{2} + b^{2}.
      \end{align*}
    \end{theorem*}
    \begin{proof*}
      Note que \(N(a+ib) = (a+ib)(a-ib)\), donde fica fácil ver que \(N((a+ib)(c+id)) = N(a+ib)N(c+id)\).
      Observe também que \(N(z) = 0\) se, e somente se, \(z = 0,\) tal que dados \(x, y\in \mathbb{Z}[i], N(xy) = N(x)N(y)
      \geq N(x)\) para todo \(y\neq0\).

      O corpo de frações \(\mathrm{Frac}(\mathbb{Z}[i]) = \mathbb{Q}[i] \subseteq{\mathbb{C}}\) pode receber a mesma norma
      induzida que a de \(\mathbb{Z}[i]\). Assim, sejam \(x, y\in \mathbb{Z}[i], y\neq0.\) Buscamos \(q, r\in \mathbb{Z}[i]\) tais que
      \(x = qy + r,\) em que \(r=0\) ou \(N(r) < N(y)\). Isso é equivalente a
      \begin{align*}
        N(x-qy) < N(y) &\Longleftrightarrow N \biggl(y \biggl(\frac{x}{y}\biggr) - qy\biggr) < N(y)\\
                       &\Longleftrightarrow N(y) N \biggl(\frac{x}{y}-q\biggr) < N(y)\\
                       &\Longleftrightarrow N \biggl(\frac{x}{y} - q\biggr) < 1.
      \end{align*}
      Com isso, basta acharmos \(q\in \mathbb{Z}[i]\) tal que \(N(x/y - q) < 1.\)

      Dado \(x/y\in \mathbb{Q}[i],\) seja \(x/y = e + if.\) Existem \(g, h\in \mathbb{Z}\)
      tais que \(|e-g|\leq 1/2\) e \(|e-h|\leq 1/2\) através do seguinte processo:

      Se \(e\in \mathbb{Q},\) existe \(n\in \mathbb{Z}\) tal que \(n\leq e < n+1.\) Se \(0 < e-n\leq 1/2\),
      então \(|e-n|\leq 1/2\). Caso contrário, \(1/2 < e-n < 1,\) tal que \(-1/2 < e - n - 1 < 0\) e, assim,
      \(|e-(n+1)| < 1/2\).

      Com isso, se \(q\coloneqq g + ih,\) temos 
      \begin{align*}
        N \biggl(\frac{x}{y} - q\biggr) &= N(e + if - g - ih)\\
                                        &= N(e - q + (f-h)i)\\
                                        &= (e-q)^{2} + (f-h)^{2}\\
                                        &=\frac{1}{2} < 1.
      \end{align*}
      Portanto, provamos o resultado. \qedsymbol
    \end{proof*}
    \begin{crl*}
      \(\mathbb{Z}[i]\) é um D.I.P.
    \end{crl*}
    \begin{example}
      Mostremos que o ideal \(\langle (3, 2 + \sqrt[]{-5}) \rangle \trianglelefteq{\mathbb{Z}[\sqrt[]{-5}]}\) não é principal. Logo, que
      \(\mathbb{Z}[\sqrt[]{-5}]\) não é um domínio euclidiano.

      Com efeito, seja \(\langle 3, 2 + \sqrt[]{-5} \rangle = \langle a + b\sqrt[]{-5} \rangle.\) Observe que 
      \[
        N':\mathbb{Z}\sqrt[]{-5}\rightarrow \mathbb{N},\quad x = a+b\sqrt[]{-5}\mapsto a^{2}+5b^{2} = N(x)
      \]
      satisfaz a condição \(N(xy) = N(x)N(y)\geq N(x).\) Dado o elemento \(3\in \langle a + b\sqrt[]{-5}, \rangle\)
      vale que \(3 = (a+b\sqrt[]{-5})(c+d\sqrt[]{-5}),\) do que obtemos que 
      \[
        9 = N(3) = N(a+b\sqrt[]{-5})N(c+d\sqrt[]{-5}) = (a^{2}+5b^{2})(c^{5}+5d^{5}).
      \]
      Com isso, \(a^{2}+5b^{2}\mid 9\) dá-nos três possibilidade - \(a^{2} + 5b^{2} = 1, 3,\text{ ou }9.\)
      Se \(a^{2} + 5b^{2} = 1\) ou \(a^{2} + 5b^{2} = 3,\) obtemos \(b = 0\), duas contradições. Fazendo mais contas,
      chegamos em ainda mais contradições, do que segue que a condição \(N(xy)\geq N(x)\) não pode ser satisfeito e, portanto,
      \(\mathbb{Z}[\sqrt[]{-5}]\) não é Domínio Euclidiano e nem \(\langle (3, 2+\sqrt[]{-5}) \rangle\) é um ideal principal.
    \end{example}
    \begin{example}[Exercício]
      Mostre que \(\mathbb{Z}[\sqrt[]{-2}] = \{a + b\sqrt[]{-2}: a, b\in \mathbb{Z}\}\) é um Domínio
      Euclidiano com a norma \(N(a+b\sqrt[]{-2}) = a^{2} + 2b^{2}.\)
    \end{example}
    \newpage

    \section{Aula 17 - 25/10/2023}
    \subsection{Motivações} 
    \begin{itemize}
      \item Corpos quadráticos;
      \item O anel dos inteiros de \(\mathbb{Q}[\sqrt[]{d}]\);
      \item Elementos irredutíveis;
      \item Anéis Noetherianos e Teorema da Base de Hilbert.
    \end{itemize}
    \subsection{Corpos Quadráticos}
    \begin{def*}
      Seja \(\alpha \) raiz de um polinômio irredutível e mônico de grau 2, ou seja, \(x^{2}+ax+b\in \mathbb{Z}[x].\) Vamos denotar
      por \(\mathbb{Z}[\alpha ]\) o menor subanel de \(\mathbb{C}\) que contém \(\mathbb{Z}\) e \(\alpha \): 
      \[
        \mathbb{Z}[\alpha ]\coloneqq \bigcap_{A\in \mathcal{A}}^{}{A},
      \]
      em que \(\mathcal{A} = \{A\text{ subanel de }\mathbb{C}: \mathbb{Z}\subseteq{A}\text{ e }\alpha \in A\}. \square\)
    \end{def*}
    \begin{prop*}
      Temos \(\mathbb{Z}[\alpha ] = \{a + b\alpha : a, b\in \mathbb{Z}\}.\)
    \end{prop*}
    \begin{proof*}
      De fato, coloque \(B\coloneqq \{c+d\alpha : c, d\in \mathbb{Z}\}.\) Note que B é subanel de \(\mathbb{C}\).
      Com isso, se \(a + b\alpha , c + d\alpha \in B,\) então \(a + b\alpha + c +d\alpha = a + c + (b+d)\alpha \in B\)
      e, a multiplicação: 
      \[
        (a+b\alpha )(c+d\alpha ) = ac + bd\alpha^{2} + ad\alpha + bc\alpha 
      \]
      e note que, já que \(\alpha \) é raiz de um polinômio, \(x^{2} + a'x + b'\), então \(\alpha ^{2} + a'\alpha  + b' = 0\) implica 
      que \(\alpha^{2} = -a'\alpha  - b'\), isto é, \(\alpha^{2} = c'\alpha + d'\) para \(c', d'\in \mathbb{Z},\) donde segue que 
      \[
        ac + bd(c'+d'\alpha ) + ad\alpha + bc\alpha = ac + bdc' + \alpha(bdd'+ad+bc)\in B.
      \]
      Logo, B é subanel de \(\mathbb{C}.\) Além disso, temos \(\mathbb{Z}\subseteq{B}, \alpha \in B\), tal que
      \(\mathbb{Z}[\alpha ]\subseteq{B}.\)

      Por outro lado, seja \(A\) um subanel de \(\mathbb{C}\) que contém \(\mathbb{Z}\) e \(\alpha .\) Caso \(c+d\alpha \in B, c, d,\in A\)
      e a combinação \(c+d\alpha \in A\). Logo, \(B \subseteq{A}.\) Portanto, \(B\subseteq{\mathbb{Z}[\alpha ]}\). \qedsymbol
    \end{proof*}
    \begin{example}[Exercícios]
      \begin{itemize}
        \item[1)] Seja \(\alpha \) raiz de \(x^{2} + ax + b\in \mathbb{Z}[x]\) irredutível. Então, 
          \[
            \mathbb{Z}[\alpha ] \cong{\frac{\mathbb{Z}[x]}{\langle x^{2} + ax + b \rangle}}
          \]
        \item[2)] Se \(\beta \) é outra raiz do mesmo polinômio, então \(\mathbb{Z}[\alpha ] = \mathbb{Z}[\beta ]\)
      \end{itemize}
    \end{example}
    Definimos 
    \begin{align*}
  &N:\mathbb{Z}[\alpha ]\setminus{\{0\}}\rightarrow \mathbb{N}\\
  &x+y\alpha \mapsto |x^{2} - axy + y^{2}b|.
    \end{align*}
    \begin{def*}
      Se \(\alpha \) é a raiz do polinômio mônico de grau 2, o \textbf{conjugado} de \(\theta = c+d\alpha \) é definido como \(\overline{\theta } = c + d\beta \), em que \(\beta \) é a outra raiz. \(\square\)
    \end{def*}
    Note que um elemento e seu conjugado satisfazem a seguinte relação
    \begin{align*}
      \theta \overline{\theta } &= (x+y\alpha )(x+y\beta)\\
                                &= x^{2} + xy\beta  + xy\alpha  + y^{2}\alpha \beta \\
                                &= x^{2} + (\alpha +\beta )xy + \alpha \beta y^{2}\\
                                &= x^{2} + (-a)xy + (b)y^{2}\\
                                &= N(\theta).
    \end{align*}
    Disto segue que 
    \begin{itemize}
      \item[1)] \(N(\theta \varphi )=N(\theta )N(\varphi )\);
      \item[2)] \(\mathbb{Zj[\alpha ]^{*} = \{\theta \in \mathbb{Z}[\alpha ]: N(\theta ) = 1\}}\).
    \end{itemize}
    \begin{def*}
      Seja \(d\in \mathbb{Z}\) livre de quadrados. O corpo \(\mathbb{Q}[\sqrt[]{d}] = \{a + b\sqrt[]{d}: a, b\in \mathbb{Q}\}\) é 
      chamado \textbf{corpo quadrático.} \(\square\)
    \end{def*}
    \begin{example}
      \begin{itemize}
        \item[1)] Os inteiros de Gauss, \(\mathbb{Z}[i]\), com o polinômio correspondente \(x^{2} + 1 = 0\).
        \item[2)] Os inteiros de Eisenstein, \(\mathbb{Z}\biggl[\frac{1+\sqrt[]{-3}}{2}\biggr]\), com o polinômio \(x^{2} - x + 1 = 0\)
        \item[3)] Para o polinômio \(x^{2} - x - 1 = 0, \mathbb{Z}\biggl[\frac{1 + \sqrt[]{5}}{2}\biggr]\)
        \item[4)] Se \(d\in \mathbb{Z}\) é \textbf{livre de quadrados}, i.e., não existe \(l\in \mathbb{Z}\)
          tal que \(l^{2}\mid d, \mathbb{Z}[\sqrt[]{d}]\) é um anel quadrático do polinômio \(x^{2} - d = 0.\)
      \end{itemize}
    \end{example}
    \begin{def*}
      Seja d livre de quadrados, denotamos por \(\mathcal{O}_{d}\) o conjunto
      \(\{\alpha \in \mathbb{Q}[\sqrt[]{d}]:\text{ existe }f(x)\in \mathbb{Z}[x], \deg{f(x)} = 2, f \text{ mônico irracionais, tais que } f(\alpha)=0\}.\quad\square\)
    \end{def*}
    \begin{theorem*}
      Se \(d\in \mathbb{Z}\) é livre de quadrados, então 
      \[
        \mathcal{O}_{d} = \left\{\begin{array}{ll}
            \mathbb{Z}[\sqrt[]{d}],\quad\text{se } d\equiv 2, 3 \mod{4}\\
            \mathbb{Z}\biggl[\frac{1+\sqrt[]{d}}{2}\biggr],\quad\text{se } d\equiv 1\mod{4}.
        \end{array}\right.
      \]
    \end{theorem*}
    \begin{def*}
      O anel \(\mathcal{O}_{d}\) é chamado o \textbf{anel dos inteiros} do corpo \(\mathbb{Q}[\sqrt[]{d}].\) Nele,
      podemos definir a norma:
      \begin{align*}
  &N_{d}:\mathcal{O}_{d}\rightarrow \mathbb{N}\\
  &x+y\theta \mapsto \left\{\begin{array}{ll}
      x^{2} + dy^{2},\quad \text{se } d\equiv2,3\mod{4}\\
      x^{2} - xy - \frac{d-1}{4}y^{2},\quad\text{se } d\equiv1\mod{4},
  \end{array}\right.
      \end{align*}
      em que \(\theta = \sqrt[]{d}\) se \(d\equiv 2, 3\mod{4}\) e \(\theta = \frac{1+\sqrt[]{d}}{2}\), se \(d\equiv 1\mod{4}.\square\)
    \end{def*}
    \begin{def*}
      Dizemos que \(\mathcal{O}_{d}\) é \textbf{norma euclidiano} se é um domínio euclidiano com a norma definida \(N_{d}\) acima. \(\square\)
    \end{def*}
    \begin{theorem*}
      \(\mathcal{O}_{d}\) é norma euclidiano se, e somente se, 
      \[
        d\in\{-11, -7, -3, -2, -1, 2, 3, 5, 6, 7, 11, 13, 17, 19, 21, 29, 33, 37, 41, 57, 73\}.
      \]
    \end{theorem*}
    \begin{theorem*}
      \begin{itemize}
        \item[1)] \(\mathcal{O}_{69}\) é euclidiano (mas não é norma euclidiano)
        \item[2)] \(\mathcal{O}_{14}\) é euclidiano (mas não é norma euclidiano)
        \item[3)] Se \(d < 0\) e \(d\not\in\{-11, -7, -3, -2, 1\},\) então \(\mathcal{O}_{d}\) não é euclidiano
        \item[4)] Se \(d < 0\), então \(\mathcal{O}_{d}\) é D.I.P., se e somente se 
          \[
            d\in\{-1, -2, -3, -7, -11, -19, -43, -67, -163\}.
          \]
      \end{itemize}
    \end{theorem*}
    Relacionada a esses anéis dos inteiros, há a seguinte conjectura:

    \textbf{Conjectura:} Classificar para quais \(d\in \mathbb{Z}\) o anel \(\mathcal{O}_{d}\) é D.I.P. e para quais
    \(\mathcal{O}_{d}\) é euclidiano.

    \begin{def*}
      Seja A um domínio. Dizemos que \(a\in A\) é \textbf{irredutível} se \(a\not\in A^{*}\) e se \(a = bc,\) então \(b\in A^{*}\) ou \(c\in A^{*}.\square\)
    \end{def*}
    \begin{example}
      \begin{itemize}
        \item[1)] Em \(\mathbb{Z}\), apenas primos positivos ou negativos são irredutíveis;
        \item[2)] Se F é um corpo, em \(F[x]\) apenas os polinômios irredutíveis são irredutíveis.
      \end{itemize}
    \end{example}
    \begin{theorem*}
      Se A é um D.I.P, então todo elemento não irredutível de A pode ser escrito como produto de elementos irredutíveis.
    \end{theorem*}
    \begin{proof*}
      Suponha que existe um \(a\in A\) que não é irredutível, porém não pode ser escrito como produto de
      elementos irredutíveis. Seja \(a=bc\) tais que \(b, c\not\in A^{*}.\)

      Note que um dos dois, b ou c, não pode ser produto de elementos irredutíveis. Suponha que este seja b.
      Disto, segue que \(\langle a \rangle\subsetneq{\langle b \rangle}\), já que, como \(a=bc, a\in \langle b \rangle\) e, assim,
      \(\langle a \rangle\subset{\langle b \rangle}\). Suponha, agora, que \(\langle a \rangle = \langle b \rangle.\) Então, \(b\in \langle a \rangle\) implica
      que \(b = ar\) para algum \(r\in A\), de forma que \(b = (bc)r.\) Como A é domínio, \(cr = 1\), ou seja, \(c\in A^{*},\)
      contrariando a hipótese antes feita.

      Seja \(a_{1}\coloneqq b.\) Segue que \(\langle a \rangle\subsetneq{\langle a_{1} \rangle}\) e \(a_{1}\) não é produto de elementos irredutíveis.
      Continuando este processo indutivamente, obteremos uma cadeia de ideais \(\langle a \rangle \subsetneq{\langle a_{1} \rangle}\subsetneq{\langle a_{2} \rangle}\subsetneq{\cdots}\)
      tal que cada \(a_{i}\) não é produto de elementos irredutíveis. Seja \(\mathfrak{i} = \bigcup_{i=1}^{\infty}{}\). Já que A é D.I.P., seja \(x\in A\) tal que \(\mathfrak{i} = \langle x \rangle\).

      Caso \(x\in \mathfrak{i},\) existe \(n\in \mathbb{N}\) tal que \(x\in \langle a_{n} \rangle\) e, então, \(\mathfrak{i} = \langle x \rangle\subseteq{\langle a_{n} \rangle}.\) Portanto,
      \(\langle a_{n} \rangle = \langle a_{n+1} \rangle = \cdots = \mathfrak{i}\), o que é um absurdo. Portanto, todo elemento irredutível de A
      pode ser escrito como produto de elementos irredutíveis. \qedsymbol
    \end{proof*}
    \newpage

    \section{Aula 18 - 30/10/2023}
    \subsection{Motivações}
    \begin{itemize}
      \item Aneis Noetherianos;
      \item Anel de Inteiros Algébricos.
    \end{itemize}
    \subsection{Aneis Noetherianos}
    \begin{prop*}
      Todo elemento de um anel quadrático é produto de elementos irredutíveis.
    \end{prop*}
    \begin{proof*}
      Seja \(\mathbb{Z}[\alpha ] = \{c + d\alpha : c, d\in \mathbb{Z}\}, \alpha \) raiz de \(x^{2} + ax + b\in \mathbb{Z}[x]\) irredutível.
      Tome \(x\in \mathbb{Z}[\alpha ]\) tal que x não é o produto dos elementos irredutível com menor norma \(N(c+d\alpha ) = [c^{2}-acd + bd^{2}]\).
      Como x não é irredutível, \(x=yz\), com \(y,z\not\in \mathbb{Z}[\alpha ]^{*},\) em que um deles não é o produto de elementos
      irredutíveis.

      Agora, note que \(N(y)\neq1\) e \(N(z)\neq1\). De fato, se \(N(y)=1, y\in \mathbb{Z}[\alpha ]^{*}\) e então \(x=yz\) seria irredutível.
      Pela propriedade \(N(x) = N(y)N(z)\) e \(N(y) > 1, N(z) > 1\), temos \(N(y) < N(x)\) e \(N(z) < N(x)\). 
      Pelo menos um deles não é produto de irredutíveis, contrariando a minimalidade de x. Portanto, x deve ser o produto
      de elementos irredutíveis. \qedsymbol
    \end{proof*}
    \begin{example}
      2 é elemento irredutível em \(\mathbb{Z}[\sqrt[]{-5}].\)

      De fato, suponha que 2 possa ser escrito como \(2 = (a + b\sqrt[]{-5})(c+d\sqrt[]{-5}).\) Então,
      \begin{align*}
   &N(2) = N(a+b\sqrt[]{-5})N(c+d\sqrt[]{-5})\\
   &4 = (a^{2} + 5b^{2})(c^{2}+5d^{2}).
      \end{align*}
      Então, \(c^{2} + 5d^{2}\mid 4\), tal que \(c^{2} + 5d^{2} = 1, 2, 4.\). Caso \(a^{2}+5b^{2} = 1, b = 0 e a = \pm 1\in \mathbb{Z}[\sqrt[]{-5}]^{*}.\) 
      Com isso, \(a^{2} + 5b^{2} = 4\) e \(a = \pm 2\) e \(d=0.\) Assim, \(2=(\pm1)(\pm2).\)
      O caso \(c^{2} + 5d^{2} = 2\) não tem soluções e o caso \(c^{2} + 5d^{2} = 4\) é simétrico ao primeiro caso.
    \end{example}
    \begin{def*}
      Dizemos que um anel A é \textbf{noetheriano} se todo ideal de A é finitamente gerado. \(\square\)
    \end{def*}
    \begin{prop*}[Exercício]
      As propriedades a seguir são equivalentes:
      \begin{itemize}
        \item[1)] A é noetheriano;
        \item[2)] Toda cadeia de ideais de A: \(\mathfrak{i}_{1}\subseteq{}\mathfrak{i}_{2}\subseteq{}\mathfrak{i}_{3}\subseteq{\cdots}\) estabiliza, i.e.,
          existe \(n\in \mathbb{N}\) tal que \(\mathfrak{i}_{n} = \mathfrak{i}_{n+1} = \cdots.\)
        \item[3)] Todo conjunto S de ideias de A tem um elemento maximal.
      \end{itemize}
    \end{prop*}
    \begin{theorem*}[Base de Hilbert]
      Se A é noetheriano, então \(A[x]\) também é noetheriano.
    \end{theorem*}
    \begin{prop*}[Exercícios]
      \begin{itemize}
        \item[1)] Seja A noetheriano e \(\mathfrak{i}\trianglelefteq{A}.\) Então, \(A/\mathfrak{i}\) é noetheriano.
        \item[2)] Seja A um domínio noetheriano. Então, todo elemento de A é produto de irredutíveis
      \end{itemize}
    \end{prop*}
    \begin{example}
      \begin{itemize}
        \item[1)] Mostre que \(\mathbb{Z}[\sqrt[]{-3}]\) não é D.I.P. (Mostre que \(\langle 2, 1 + \sqrt[]{-3} \rangle\) não é principal);
        \item[2)] Seja \(F\subseteq{\mathbb{C}}\) subcorpo. Então, \(\mathbb{Q}\subseteq{F}.\);
        \item[3)] Seja F um subcorpo de \(\mathbb{C}\) tal que \([F:\mathbb{Q}]<\infty\). (Chamamos estes corpo de \textbf{corpos globais});
        \item[4)] Seja \(\mathcal{O}_{F}\coloneqq \{\alpha \in F: \exists f(x)\in \mathbb{Z}[x]\} \text{ mônico irredutível e } f(\alpha ) = 0\). Mostre que
          \(\mathcal{O}_{F}\) é um anel. (Chamamos este anel de \textbf{anel de inteiros algébricos} de F). Além disso, \(\mathrm{Frac}(\mathcal{O}_{F}) = F.\);
        \item[5)] Se \(F = \mathbb{Q}[\sqrt[]{d}],\) d livre de quadrados, então \(\mathcal{O}_{F} = \mathcal{O}_{d}.\)
      \end{itemize}
    \end{example}
    \begin{example}
      \begin{itemize}
        \item[1)] \(\mathcal{O}_{-3} = \{\pm1, \pm \frac{1\pm \sqrt[]{-3}}{2}\}\);
        \item[2)] \(\mathcal{O}_{-1} = \{\pm1, \pm i\}\);
        \item[3)] \(2, 3\in \mathcal{O}_{-19}\) são irredutíveis;
        \item[4)] Seja p primo. Se \(p\equiv 1 \mod 4,\) então \(p\in \mathbb{Z}[i]\) não é irredutível.
      \end{itemize}
    \end{example}
    \newpage

    \section{Aula 19 - 01/11/2023}
    \subsection{Motivações} 
    \begin{itemize}
      \item Teorema de Dedekind;
      \item Anéis quasi-euclidianos;
      \item Domínios de Fatoração Única.
    \end{itemize}
    \subsection{O Teorema de Dedekind}
    \begin{theorem*}
      \(\mathcal{O}_{-19}\) não é euclidiano.
    \end{theorem*}
    \begin{proof*}
      Suponha que exista \(N':\mathcal{O}_{-19}\setminus{\{0\}}\rightarrow \mathbb{N}\) tal que \(\mathcal{O}_{-19}\) é um domínio euclidiano com norma \(N'.\) Seja \(m\in \mathcal{O}_{-19}, m\neq \pm1, 0,\)
      tal que 
      \[
        N'(m) = \min\{N'(a): a\in \mathcal{O}_{-19}, a\neq \pm1, 0\}.
      \]
      Pelo algoritmo da divisão, existem \(q, r\) tais que \(2 = mq + r\), com \(N'(r) < N'(m)\)
      ou \(r=0\).

      Como m tem menor norma dentre os elementos diferentes de \(0\) e \(\pm 1\), então \(r=0, \pm1\). Se \(r=1\), temos
      \(mq = 1\) implica que \(m\in \mathcal{O}_{-19}^{*},\) um absurdo. Assim, \(r=0, -1\) e,
      logo, \(mq = 2, 3\) e, já que \(2, 3\) são irredutíveis, \(q\in \mathcal{O}_{-19}^{*} = \{\pm1\}.\) Com isso,
      \(m=\pm2, \pm 3\in \mathbb{Z}.\)

      Agora, seja \(\theta = \frac{1 + \sqrt[]{-19}}{2}.\) Então, existem \(q', r'\) tais que \(\theta  = mq'+r'.\) Usando o mesmo
      argumento que o feito acima, temos \(r'=0, \pm1.\) Destarte, \(mq'=\theta - r', r'=0, \pm1.\) Seja \(q'= a + b\theta ,\)
      em que \(a, b\in \mathbb{Z}.\) Logo, 
      \[
        m(a + b\theta ) = -r + \theta  \Rightarrow ma + mb\theta = -r + \theta  \Rightarrow mb = 1,
      \]
      mas, note que, como \(m = \pm2, \pm3\), temos \(b\not\in \mathbb{Z}\), ou seja, não existe \(q'\) que satisfaça \(\theta  = mq' + r'\) e, portanto,
      \(\mathcal{O}_{-19}\) não é euclidiano. \qedsymbol
    \end{proof*}
    \begin{def*}
      Um domínio A é dito \textbf{quasi-euclidiano} se existe uma aplicação \(d:A\setminus{\{0\}}\rightarrow \mathbb{N}\) tal que para todos \(a, b\in A, b\neq0\),
      existem \(p, q, r\in A\) com \(p\neq0\) tais que: 
      \[
        ap = bq + r\quad\&\quad d(r) < d(b) \text{ ou } (p=1 e r=0).\quad\square
      \]
    \end{def*}
    \begin{lemma*}
      Todo domínio quasi-euclidiano é D.I.P.
    \end{lemma*}
    \begin{proof*}
      Seja \((0)\neq \mathfrak{i}\trianglelefteq{A}\) e \(b\in \mathfrak{i}\) tal que \(d(b) = \min\{d(x):x\in \mathfrak{i}\setminus{\{\}}\}\).
      Tome \(a\in \mathfrak{i}.\) Então, existem \(p, q, r\in A\) tais que \(p\neq 0\) e \(pa = bq + r\), com \(d(r) < d(b)\) ou \(p=1, r=0.\)
      Caso \(r\neq 0,\) então \(r = pa-bq\in \mathfrak{i}\) e, assim, \(d(r) < d(b),\) um absurdo com a minimalidade de b em \(\mathfrak{i}.\)

      Com isso, \(r=0\) e, então, \(a = ap = bq\), tal que \(a\in \langle b \rangle.\) Além disso, segue que \(\mathfrak{i}\subseteq{\langle b \rangle}.\)
      Por outro lado, como temos \(b\in \mathfrak{i}, \mathfrak{I} = \langle b \rangle\). Portanto, todo ideal de A é principal, ou seja, A é D.I.P. \qedsymbol
    \end{proof*}
    \begin{theorem*}
      \(\mathcal{O}_{-19}\) é quasi-euclidiano. (Em particular, é D.I.P).
    \end{theorem*}
    \begin{proof*}
      Considere a norma \(N:\mathcal{O}_{-19}\rightarrow \mathbb{N}\) definida por \(N(a+b\theta ) = a^{2} + ab + 5b^{2}.\) Essa norma é a
      restrição dos complexos em \(\mathcal{O}_{-19}.\) De fato,
      \begin{align*}
        a + b\theta &= a +b \biggl(\frac{1+\sqrt[]{-19}}{2}\biggr) = \biggl(a + \frac{b}{2}\biggr) + \frac{\sqrt[]{19}}{2}ib.\\
                    &= a' + b'i\in \mathbb{C}.
      \end{align*}
      Representando \(a + b\theta \) desta forma, conseguimos calcular a norma como 
      \begin{align*}
        |a+b\theta |^{2} &= \biggl(a + \frac{b}{2}\biggr)^{2} + \biggl(\frac{\sqrt[]{19}}{b}\biggr)^{2}\\
                         &= a^{2} + ab + \frac{b^{2}}{4} + \frac{19}{4}b^{2}\\
                         &= a^{2} + ab + 5b^{2} = N(a + b\theta ).
      \end{align*}
      Sejam \(a, b\in a, b\neq0.\) Se \(a\in \langle b \rangle, a = bq\) para algum \(q\in A\) (ou seja, \(p=1, r = 0\)).
      Caso \(a\not\in \langle b \rangle\), a existência de \(p, q, r\in A\) com \(p\neq 1\) e \(r\neq 0\) depende da condição 
      a seguir:
      \begin{align*}
        N(r) &= N(ap-bq) < N(p)\\
        \Longleftrightarrow & N(b)N \biggl(\frac{a}{b}p - q\biggr) < N(b)\\
        \Longleftrightarrow & N \biggl(\frac{a}{b}p - q\biggr) < 1\\
        \Longleftrightarrow & \biggl|\frac{a}{b}p - q\biggr| < 1. \quad\text{\qedsymbol}
      \end{align*} 
    \end{proof*}
    \begin{lemma*}
      Para todo \(a, b\in A = \mathcal{O}_{-19},\) com \(b\neq0, a\not\in \langle b \rangle,\) existem \(p, q\in A\) tais que 
      \(\biggl|\frac{a}{b}p - q\biggr| < 1.\)
    \end{lemma*}
    \subsection{Domínios de Fatoração Única}
    \begin{def*}
      Seja A um domínio e sejam \(a, b\in A\).
      \begin{itemize}
        \item[1)] Dizemos que a \textbf{divide} b (ou b é \textbf{múltiplo} de x), e escrevemos \(a\mid b,\) se existe \(c\in A\) tal que \(b = ac.\)
        \item[2)] Dizemos que a, b são \textbf{associados} se existe \(c\in A^{*}\) tal que a = bc. \(\square\)
      \end{itemize}
    \end{def*}
    \begin{prop*}[Exercício]
      \begin{itemize}
        \item[1)] \(a\mid b \Longleftrightarrow \langle b \rangle \subseteq{\langle a \rangle}\);
        \item[2)] a, b são associados se, e somente se, \(\langle b \rangle = \langle a \rangle\);
        \item[3)] a é irredutível se, e somente se, todo divisor de a é associado, ou a 1, ou a a.
      \end{itemize}
    \end{prop*}
    \begin{def*}
      Dizemos que \(0\neq a\in A\setminus{A^{*}}\) é \textbf{primo} se \(a\mid xy \Rightarrow a\mid x\) ou \(a\mid y.\quad\square\)
    \end{def*}
    \begin{prop*}[Exercício]
      Um elemento \(a\in A\) é primo se, e somente se, \(\langle a \rangle\in \mathrm{Spec}(A).\)
    \end{prop*}
    \begin{lemma*}
      Elementos primos são irredutíveis.
    \end{lemma*}
    \begin{proof*}
      Seja \(a\in A\) e \(a = xy, x, y\in A.\) Como \(a\mid a, a\mid xy\) e \(a\mid x\) ou \(a\mid y\).

      Suponha, sem perda de generalidade, que \(a\mid x\). Então, existe \(c\in A\) tal que \(x=ac\) e, assim, \(a=xy=acy\), de forma que \(cy = 1\), ou seja,
      \(y\in A^{*}.\) Portanto, a é irredutível. \qedsymbol
    \end{proof*}
    \begin{example}[Exercício]
      Seja \(A = \mathbb{Z}\) ou \(A = F[x]\), em que F é um corpo. Mostre que os elementos irredutíveis de A são primos.
    \end{example}
    \begin{def*}
      Um domínio A é chamado \textbf{Domínio de Fatoração Única (D.F.U.)} ou \textbf{Domínio Fatorial} se 
      todo elemento não nulo \(a\in A\) pode ser escrito como produto de elementos irredutíveis de maneira única a menos 
      da ordem dos fatores e de elementos associados. Mais explicitamente, dado \(a\in A\setminus{\{0\}},\)
      \begin{itemize}
        \item[1)] Existem \(\pi_{1}, \cdots, \pi_{r}\) irredutíveis em A tais que \(a = \pi_{1}\cdot \dotsc \cdot \pi_{r}.\)
        \item[2)] Se \(a = \rho_{1}\cdot\dotsc \cdot \rho_{m} = \pi_{1}\cdot \dotsc \cdot \pi_{r}\), com \(\pi_{i}, \rho_{j}\) irredutíves em A, então \(r = m\)
          e existe \(b\in S_{r}\) permutação tal que \(\pi_{i}\) é associado com \(\rho_{b(i)}\) para todo \(i=1, \cdots, r. \quad\square\)
      \end{itemize}
    \end{def*}
    \begin{prop*}
      Num D.F.U., elementos irredutíveis são primos.
    \end{prop*}
    \begin{proof*}
      Seja a irredutível e \(a\mid xy\). Então, existe \(c\in A\) tal que \(ac = xy.\) Coloque 
      \[
        x = \pi_{1} \cdot \dotsc \cdot \pi_{n},\quad y = \pi_{n+1}\cdot \dotsc \cdot \pi_{n+m}\quad\&\quad c=\rho_{1}\cdot \dotsc \rho_{r}.
      \]
      Temos \(a\mid ac = xy = \pi_{1}\cdot \dotsc \cdot \pi_{n+m},\) tal qual existe \(d\in A\) satisfazendo \(a = d\pi_{1}\cdot \dotsc \cdot \pi_{n+m}.\) Como a é 
      irredutível e cada \(\pi_{i}\) é irredutível, \(d\in A^{*}.\) Assim, a é associado a \(\pi_{i}\) para algum \(i.\) Se \(i\leq n, a\mid x\) e, se \(n < i\leq m + n, a \mid y.\)
      Portanto, a é primo. \qedsymbol
    \end{proof*}
    \begin{prop*}
      Se A é um domínio Noetheriano e todo elemento irredutível é primo, então A é D.F.U.
    \end{prop*}
    \begin{proof*}
      Sabemos que, se A é noetheriano, todo elemento de A é produto de elementos irredutíveis. Precisamos, então, provar apenas que essa representação é única.

      Seja \(a\in A\setminus{\{0\}}\) e \(a = \pi_{1}\cdot \dotsc \cdot \pi_{n} = \rho_{1}\cdot \dotsc \cdot \rho_{m},\) em que \(\rho_{j}, \pi_{i}\) são elementos
      irredutíveis em A. Faremos essa prova por indução. Se n = 1, então \(a = \pi_{1}\) e, como \(a\mid a\), \(\pi_{1}\mid\rho_{1}\cdot \dotsc \cdot \rho_{m}\). Já que 
      \(\pi_{1}\) é primo, temos \(\pi_{1}\mid\rho_{j}\) para algum \(1\leq j\leq m.\) Sem perda de generalidade, podemos supor \(j=1\) (basta reordenar via permutação). Então, existe \(c\in A\) tal que \(\rho _{1} = c\pi_{1}.\) Como \(\rho_{j}\) é irredutível,
      \(c\in A^{*},\) pois \(\pi_{1}\) é irredutível, ou seja, \(\pi_{1}\not\in A^{*}\).

      Suponha que \(m > 1,\) então 
      \[
        \pi_{1} = \rho_{1}\cdot \dotsc \cdot \rho_{m} \Rightarrow c\pi_{1}\cdot \rho_{2}\cdot \dotsc \cdot \rho_{m} = \pi_{1} \Rightarrow c\rho_{2} \cdot \dotsc \cdot \rho_{m} = 1.
      \]
      Assim, temos \(\rho_{2}\in A^{*},\) um absurdo. Logo, \(m=1\) e \(\pi_{1} = \rho_{1}\). 

      Agora, seja \(n > 1\). Segue que \(\pi_{1}\mid \rho_{1}\cdot \dotsc \cdot \rho_{m}\) e existe j tal que \(\pi_{1}\mid\rho_{j}.\) Podemos supor novamente
      que j=1, tal que, pelo mesmo argumento anterior, existe \(c\in A^{*}\) tal que \(\pi_{1} = c\rho_{1}\) e, assim, 
      \[
        \pi_{1}\cdot \dotsc \cdot \pi_{n} = c\rho_{1}\cdot \rho_{2} \cdot \dotsc \cdot \rho_{m} \Rightarrow \pi_{2}\cdot \dotsc \pi_{n} = \rho_{2}\cdot \dotsc \cdot \rho_{m}
      \]
      Renumerando os termos da última igualdade tal que eles vão de 1 até n-1 do lado esquerdo e de 1 até m-1 do lado direito, pela hipótese de indução, existe \(b\in S_{m-1}\)
      tal que \(\pi_{i} = c_{i}\rho_{b(i)}\) e \(n-1 = m-1.\) Portanto, n = m e os termos são associados. \qedsymbol
    \end{proof*}
    \newpage

    \section{Aula 20 - 06/11/2023}
    \subsection{Motivações}
    \begin{itemize}
      \item Domínios de Fatoração e primos;
      \item Primos de \(\mathcal{O}_{-3}\);
      \item M.D.C. e M.M.C.
    \end{itemize}
    \subsection{D.F.U.'s e Exemplos}
    \begin{theorem*}
      Se A é um D.I.P., então A é D.F.U.
    \end{theorem*}
    \begin{proof*}
      Como um D.I.P. é noetheriano, basta mostrar que elementos irredutíveis são primos.
      Seja \(\pi \in A\) irredutível e seja \(\pi \mid ab.\) Suponha que \(\pi\) não divide a e considere
      o ideal \(\langle \pi , a \rangle.\) Seja 
      \[
        \langle \pi , a \rangle = \langle d \rangle.
      \]
      Temos 
      \[
        \pi \in \langle d \rangle \Rightarrow \exists d'\in A: \pi  = dd'.
      \]
      Assim, \(d\in A^{*},\) ou \(d'\in A^{*}\), tal que \(d'^{-1}\pi  = d\), do que segue que \(\pi \mid d\) e,
      como \(d\mid a,\) temos \(\pi \mid a.\) Contradição. Logo, a única possibilidade é \(d\in A^{*}\), donde segue que 
      \[
        \langle \pi , a \rangle = \langle d \rangle = \langle 1 \rangle = A.
      \]
      Já que \(1\in A, 1\in \langle \pi , a \rangle\), ou seja, existem r e s tais que 
      \[
        1 = r\pi +sa.
      \]
      Desta forma,
      \begin{align*}
        b = b \cdot 1 &= b(r\pi + sa)\\
        \Rightarrow & b = br\pi + sab.
      \end{align*}
      Como \(\pi \mid ab\) e \(\pi \mid \pi ,\) portanto, 
      \[
        \pi \mid(rb\pi +sab) = b.\quad \text{\qedsymbol}
      \]
    \end{proof*}
    \begin{example}
      \begin{itemize}
        \item[1)] Sabe-se que \(\mathbb{Z}\) é D.F.U;
        \item[2)] Se F é um corpo, \(F[x]\) é um D.F.U.;
        \item[3)] Domínio Euclidianos são D.F.U.;
        \item[4)] Domínios quasi-euclidianos são D.F.U. (em particular, \(\mathcal{O}_{-19})\);
        \item[5)] Seja A um D.I.P. Se \(0\neq\mathfrak{p}\in \mathrm{Spec}(A)\), então \(A/\mathfrak{p}\) é D.I.P. e é um corpo;
        \item[6)] Se A é um D.I.P. e \(S\subseteq{A}\) é um conjunto multiplicativo, então \(S^{-1}A\) é um D.I.P. e, assim, um D.F.U. Seja \(\mathfrak{p}\in \mathrm{Spec}(A), \mathfrak{p} = \langle \pi  \rangle, \pi\) primo. Considere o conjunto 
          \[
            A_{\mathfrak{p}} = \biggl\{\frac{a}{b}: a, b \in A, \pi \text{ não divide } b\biggr\}.\quad \text{(\textbf{``Discrete Valuation Ring'' - Anel de Avaliação Discreta})}
          \]
          Então, se \(\mathfrak{p} = (0),\) segue que \(A_{\mathfrak{p}} = \mathrm{Frac}(A)\) e, se \(\mathfrak{p}\neq (0),\) então
          \(A_{\mathfrak{p}}\) é anel local com ideal maximal 
          \[
            \mathfrak{p}A_{\mathfrak{p}} = \biggl\{\frac{a}{b}: a,b\in A, \pi \mid a, \pi \text{ não divide } b\biggr\}
          \]
          Assim, \(\mathrm{Spec}(A_{\mathfrak{p}}) = \{(0), \mathfrak{p}A_{\mathfrak{p}}\}\) e os elementos irredutíveis primos de \(A_{\mathfrak{p}}\)
          são associados com \(\pi \). Além disso, dado \(\mathfrak{j}\trianglelefteq{A_{\mathfrak{p}}}, \mathfrak{j} = \langle x \rangle, x\in A,\) então
          \(x = \pi ^{n}\) para algum n, pois só temos um elemento irredutível. Com isso, \(\mathfrak{j} = \langle \pi ^{n} \rangle = \langle \pi  \rangle^{n} = (\mathfrak{p}A_{\mathfrak{p}})^{n} = \mathfrak{p}^{n}A_{\mathfrak{p}}\).
        \item[7)] Tome \(p\in A\) primo, A D.I.P. e coloque 
          \[
            S = \{p^{n}: n\in \mathbb{N}\cup \{0\}\} = \{1, p, p^{2}, p^{3}, \cdots\}.
          \]
          O conjunto \(A \biggl[\frac{1}{p}\biggr] = S^{-1}A\) é um D.I.P. e, além disso, 
          \[
            S^{-1}A = \biggl\{\frac{a}{p^{n}}:a\in A, n\in \mathbb{N}\cup\{0\}\biggr\}.
          \]
          Se \(\mathcal{P}\) é o conjunto dos primos de A, então o conjunto dos primos de \(A \biggl[\frac{1}{p}\biggr]\) é 
          \(\mathcal{P}\setminus{\{p\}}\).
      \end{itemize}
    \end{example}
    É importante notar que \textbf{nem todo D.F.U.} é D.I.P. Antes disso, porém, vamos ver como encontrar
    todos os ideais primos dos inteiros de Gauss:
    \begin{example}
      Considere \(\mathcal{O}_{-1} = \mathbb{Z}[i] = \{a + ib: a, b\in \mathbb{Z}\}\subseteq{\mathbb{C}}\) e considere o elemento \(13\in \mathbb{Z}\subseteq{\mathbb{Z}[i]}\).
      Em \(\mathbb{Z},\) 13 é primo, mas em \(\mathbb{Z}[i]\), temos 
      \[
        13 = 9 + 4 = (3+2i)(3-2i).
      \]
      Como \(\mathbb{Z}[i]^{*} = \{\pm1, \pm i,\}\) temos \(3-2i, 2+2i\not\in \mathbb{Z}[i]^{*}\) e, assim,
      13 é o produto de dois elementos não invertíveis, fazendo com que 13 não seja primo nos inteiros de Gauss.
    \end{example}
    \begin{example}
      Seja \(\mathcal{O}_{-3} = \{a + b\theta : a, b\in \mathbb{Z}\} \subseteq{\mathbb{C}}, \theta  = \frac{1 + \sqrt[]{-3}}{2}\). Vimos que 
      \[
        \mathcal{O}_{-3} = \biggl\{\pm 1, \frac{\pm1 \pm \sqrt[]{-3}}{2}\biggr\}.
      \]
      Então, 

      \begin{center}
        \begin{table}[h!]
          \caption{Casos dos elementos}
          \centering
          \begin{tabular}{| c | c |}
            \hline
            Produtos de \(\mathcal{O}_{-3}^{*}\) & Resultado\\
            \hline
            \(\theta ^{2} - \theta  + 1 = 0\) & \(\mathbb{Z} \subseteq{}\mathcal{O}_{-3}\)\\
            \hline
            \(\theta (\theta -1)=-1\) & \(-2\in \mathbb{Z}\) é primo\\
            \hline
            \(-\theta (\theta -1) = 1\) & \(-2\in \mathcal{O}_{-3}, 2 = (1-\sqrt[]{-3})\underbrace{\biggl(\frac{1+\sqrt[]{-3}}{2}\biggr)}_{\in \mathcal{O}_{-3}^{*}} = (1+\sqrt[]{-3})\underbrace{\biggl(\frac{1-\sqrt[]{-3}}{2}\biggr)}_{\in \mathcal{O}_{-3}^{*}}\)\\
            \hline
          \end{tabular}
        \end{table}
      \end{center}

      Assim, vemos que \(2\) e \((1\pm\sqrt[]{-3})\) são associados. Como \(2\in \mathcal{O}_{-3}\) é primo, \(1\pm\sqrt[]{-3}\) são primos também. Logo,
      \(\langle 2 \rangle = \langle 1 + \sqrt[]{-3} \rangle = \langle 1 - \sqrt[]{-3} \rangle.\) Para ver isso,
      utiliza-se a norma induzida dos complexos em \(\mathcal{O}_{-3}\) e escreve-se 2 e 4 como elementos de \(\mathcal{O}_{-3}\).
    \end{example}
    \subsection{Máximo Divisor Comum e Mínimo Múltiplo Comum}
    \begin{def*}
      Sejam A um domínio e \(a_{1}, \cdots, a_{n}\in A\setminus{\{0\}}\). O \textbf{maior divisor comum (m.d.c)} de \(a_{1}, \cdots, a_{n}\)
      é um elemento \(d = \mathrm{mdc}(a_{1}, \cdots, a_{n})\in A\setminus{\{0\}}\) tal que 
      \begin{itemize}
        \item[i)] \(d\mid a_{i}\) para todo \(i = 1, \cdots, n;\)
        \item[ii)] Se \(a\in A\setminus{\{0\}}\) e \(a\mid a_{i}\) para todo \(i=1, \cdots, n\), então \(a\mid d.\quad\square\)
      \end{itemize}
    \end{def*}
    \begin{def*}
      Sejam A um domínio e \(a_{1}, . ., a_{n}\in A\setminus{\{0\}}\). O \textbf{menor múltiplo comum} de \(a_{1}, \cdots, a_{n}\) é um elemento
      \(m = \mathrm{mmc}(a_{1}, \cdots, a_{n})\in A\setminus{\{0\}}\) tal que
      \begin{itemize}
        \item[i)] \(a_{i}\mid m\) para todo \(i=1, \cdots, n\);
        \item[ii)] Se \(m'\in A\setminus{\{0\}}\) e \(a_{i}\mid m'\) para todo \(i=1, \cdots, n\), então \(m\mid m'.\quad\square\)
      \end{itemize}
    \end{def*}
    \begin{prop*}
      \begin{itemize}
        \item[1)] Se \(d, d'\) são maiores divisores comuns de \(a_{1}, \cdots, a_{n}\in A\), então eles são associados (existe \(s\in A^{*}: d = d's\));
        \item[2)] Se \(m, m'\) são menores múltiplos comuns de \(a_{1}, \cdots, a_{n}\in A\), então eles são associados (existe \(s\in A^{*}: m = m's\)).
      \end{itemize}
    \end{prop*}
    \begin{example}
      Considere \(2, 3\in \mathbb{Z}.\) Então, \(\mathrm{mdc}(2, 3) = \pm1\)
    \end{example}
    \begin{theorem*}[Exercício]
      Sejam A um D.F.U. e \(a_{1}, \cdots, a_{n}\in A\setminus{\{0\}}.\) Suponha que 
      \[
        a_{j} = a_{j}\pi_{1}^{e_{j1}}\cdot\dotsc \cdot \pi_{n}^{e_{jn}} \quad (e_{ji} = 0 \text{ se } \pi_{j} \text{ não divide } a_{j}),
      \]
      em que \(\pi_{i}\) são irredutíveis e dois-a-dois distintos. Com isso, 
      \begin{align*}
    &i)\quad\mathrm{mdc}(a_{1},\cdots,a_{n}) = \prod\limits_{i=1}^{n}\pi_{i}^{\min\{e_{1i}, \cdots, e_{nj}\}};\\
    &ii)\quad \mathrm{mmc}(a_{1}, \cdots, a_{n}) = \prod\limits_{i=1}^{n}\pi_{i}^{\max\{e_{1i}, \cdots, e_{ni}\}};\\
    &iii)\quad a_{1}\cdot \dotsc \cdot  a_{n} = \mathrm{mdc}(a_{1}, \cdots, a_{n})\cdot \mathrm{mmc}(a_{1}, \cdots, a_{n});\\
    &iv)\quad \text{Se A é um D.I.P., então }  \langle a_{1}, \cdots, a_{n} \rangle = \langle \mathrm{mdc}(a_{1}, \cdots, a_{n}) \rangle.
      \end{align*}
    \end{theorem*}
    \begin{prop*}
      Os primos de \(\mathcal{O}_{-1} = \mathbb{Z}[i]\), a menos de associação, são:
      \begin{itemize}
        \item[1)] Primo \(p\in \mathbb{Z}\) tal que \(p\equiv 3 \mod 4\)
        \item[2)] \(a+ib\in \mathbb{Z}[i]\) tais que \(a^{2} + b^{2} = p, p\) primo com \(p = 2\) ou \(p\equiv 1 \mod 4\)
      \end{itemize}
    \end{prop*}
    \begin{proof*}
      Seja \(p\equiv 3\mod 4\) e seja 
      \[
        p = \alpha \cdot \beta , \alpha , \beta \in \mathbb{Z}[i].
      \]
      Temos \(N(p) = N(\alpha )\cdot N(\beta )\), tal que \(p^{2} = N(\alpha )\cdot N(\beta )\). Logo,
      \(N(\alpha ) = 1, p, ^{2},\) a partir donde analisamos casos.
      \begin{itemize}
        \item Se \(N(\alpha ) = 1\), então \(\alpha \in \mathbb{Z}[i]^{*}\);
        \item Se \(N(\alpha ) = p^{2},\) então \(N(\beta ) = 1\) e \(\beta \in \mathbb{Z}[i]^{*}\);
      \end{itemize}
      Resta o caso \(N(\alpha ) = N(\beta ) = p.\) Com este caso, segue que 
      \[
        \alpha = a + ib \Rightarrow N(\alpha ) = a^{2} + b^{2} = p,
      \]
      tal que \(\overline{a}^{2} + \overline{b}^{2} = 0\) em \(\mathbb{F}_{p} = \mathbb{Z}/\{p\}\). Consequentemente,
      \(p^{2}\mid a^{2}+b^{2} = p\), um absurdo. Logo, \(\overline{a}\neq \overline{0}\) e \(\overline{b}\neq \overline{0}\), o que implica que 
      \[
        \biggl(\frac{\overline{a}}{\overline{b}}\biggr)^{2} + 1 = 0 \text{ em } \mathbb{F}_{p}.
      \]
      Logo, \(x^{2} + 1 = 0\) tem uma raiz em \(\mathbb{F}_{p},\) o que ocorre apenas se \(p\equiv 1 \mod 4,\) que também
      é um absurdo. Conclui-se, assim, que o caso \(N(\alpha ) = p\) é impossível, restando apenas as outras opções.

      Agora, seja \(a+ib\in \mathbb{Z}[i]\) primo tal que \(a^{2} + b^{2} = p, p\) um primo de \(\mathbb{Z}.\)
      Vamos mostrar que \(p = 2\) ou \(p\equiv 1\mod 4\). Com efeito, considerando \(\pi  = \alpha \cdot \beta \), vale 
      \[
        N(\pi ) = N(\alpha )\cdot N(\beta ).
      \]
      Com isto, 
      \[
        p = a^{2} + b^{2} = N(\pi) = N(\alpha )\cdot N(\beta ) \Rightarrow N(\alpha ) = 1 \text{ ou } N(\beta ) = 1,
      \]
      do que segue o fato de \(\pi \) ser um primo. Analogamente ao argumento acima, se \(N(\alpha ) = p\), então \(x^{2} + 1\) tem raiz em \(\mathbb{F}_{p}\),
      tal que \(p\equiv 1 \mod 4\) ou \(p = 2\).

      Fica como exercício mostrar que não há outros primos além destes. \qedsymbol
    \end{proof*}
    \begin{example}
      Os elementos \(3, 7, 11, 19, 23, \dotsc\) são primos, assim como \(1\pm i, 1\pm 2i, 2\pm i, 3\pm 2i, 4\pm i, 5\pm 2i, \dotsc\) 
    \end{example}
    \begin{example}[Exercício]
      Encontre todos os primos de \(\mathcal{O}_{-3}\). (Dica: \(x^{2}+x+1\) tem taíz em
      \(\mathbb{F}_{p} \Longleftrightarrow p = 3\) ou \(p\equiv 1\mod 3.\) Além disso, \(N(a + b\theta ) 
      = a^{2} + ab + b^{2} = p \Rightarrow \overline{a}^{2} + \overline{a}\overline{b} + \overline{b}^{2} = \overline{0}\)
      em \(\mathbb{F}_{p},\) e então \((\overline{a}/\overline{b})^{2} + (\overline{a}/\overline{b}) + 1 = \overline{0}.\))
    \end{example}
    \newpage

    \section{Aula 21 - 08/11/2023}
    \subsection{Motivações} 
    \begin{itemize}
      \item Polinômios Primitivos;
      \item O Lema de Gauss;
      \item O Critério de Eisenstein;
    \end{itemize}
    \subsection{O Lema de Gauss}
    \begin{def*}
      Seja A um D.F.U. Um polinômio não nulo \(f(x)\in A[x]\) com coeficientes não nulos \(a_{i_{1}},\dotsc, a_{i_{k}}\in A\) é dito
      \textbf{primitivo} se \(\mathrm{mdc}(a_{1}, \dotsc, a_{k})\) é associado a 1. \(\square\)
    \end{def*}
    \begin{example}
      Todo polnômio mônico é primitivo.
    \end{example}
    \begin{lemma*}
      Seja A um D.F.U. e \(K = \mathrm{Frac}(A)\).
      \begin{itemize}
        \item[1)] Sejam \(a, b\in A\) e \(f(x), g(x)\in A[x].\) Se \(f(x)\) é primitivo e \(af(x) = bg(x),\)
          então \(b\mid a.\)
        \item[2)] Se \(f(x)\in K[x], f(x)\neq 0\), então \(f(x)\) é associado (em K[x]) a um polinômio primitivo em \(A[x]\).
      \end{itemize}
    \end{lemma*}
    \begin{proof*}
      \(1.):\) Seja \(b = u\pi_{1}^{r_{1}}\cdot \dotsc \cdot \pi_{n}^{r_{n}}\), com \(\pi_{i}\) irredutíveis e distintos, \(u\in A^{*},\)
      e seja \(f(x) = \sum\limits_{i=1}^{m}c_{i}x^{i}.\) Como \(f(x)\) é primitivo, dado \(\pi_{j}\), existe \(c_{ij}\) tal que \(\pi_{j}\) 
      não divide \(c_{ij}.\) Da igualdade \(af(x) = bg(x),\) segue que \(b\mid ac_{i}\), pois b divide os coeficientes
      de bg, que é iguais aos de af. Como \(\pi_{j}^{r_{j}}\mid b\) e \(\pi_{j}\) não divide \(c_{ij},\) resta que \(\pi_{j}^{r_{j}}\mid a.\) 
      Como isso vale para todo j, temos \(b\mid a.\)

      \(2.):\) Seja \(f(x)\in K[x],\) ou seja, 
      \[
        f(x) = \sum\limits_{i=0}^{n}\frac{a_{i}}{b_{i}}x_{i},
      \]
      em que \(a_{i}, b_{i}\in A.\) Coloque \(c = b_{n}\cdot \dotsc \cdot b_{0}.\) Então,
      \(f(x) = \frac{1}{c} \cdot cf(x),\) tal que 
      \[
        f(x) = \frac{1}{c}\biggl(d_{n}x^{n}+ \dotsc +d_{0},\biggr)
      \]
      em que \(d_{i} = a_{i}b_{n} \cdot \dotsc \cdot \hat{b_{i}}\cdot \dotsc \cdot b_{0}.\) Note que o polinômio 
      \(d_{n}x^{n} + \dotsc  + d_{0}\in A[x]\).
      Caso \(e = \mathrm{mdc}(d_{i1}, \dotsc , d_{ik},\) em que \(d_{ij}\) são os não nulos, então 
      \[
        f(x) = \frac{3}{c}\biggl(\frac{d_{n}}{e}x^{n} + \dotsc + \frac{d_{0}}{e}\biggr)\coloneqq \frac{e}{c}f_{0}(x).
      \]
      Por definição, \(f_{0}\in A[x]\) e, além disso, \(\mathrm{mdc}(d_{n}/e, \dotsc , d_{0}/e) = 1,\) ou seja, \(f_{0}\) é primitivo.
      Caso \(\alpha\coloneqq e/c \in K^{*},\) segue que \(f = \alpha f_{0}.\) Portanto, \(f\) e \(f_{0}\) estão associados, como queríamos demonstrar. \qedsymbol
    \end{proof*}
    \begin{lemma*}
      Seja A um D.F.U. Se \(f, g\in A[x]\) são primitivos, então \(f(x)g(x)\) também é primitivo.
    \end{lemma*}
    \begin{proof*}
      Seja \(\pi \in A\) um primo. Então, \(\mathfrak{p} = \langle \pi  \rangle\in \mathrm{Spec}(A).\) Como f, g
      são primitivos, seus coeficientes não são todos múltiplos de \(\pi \), i.e., \(\overline{f(x)}\neq 0\) e 
      \(\overline{g(x)}\neq 0\) em \((A/\mathfrak{p})[x].\)

      Sabemos que \((A/\mathfrak{p})[x]\) é domínio, já que A é D.F.U. Assim, \(\overline{f(x)g(x)} = \overline{f(x)}\overline{g(x)}\neq0.\) Fazendo este processo para 
      todo elemento \(\pi \in A\) primo, segue que \(f(x)g(x)\) é primitivo. \qedsymbol
    \end{proof*}
    \begin{lemma*}
      Seja A um D.F.U., \(f(x)\in A[x]\) primitivo e \(g(x)\in A[x].\) Caso \(K = \mathrm{Frac}(A)\) e se \(f(x)\mid g(x)\) em \(K[x]\), então
      \(f(x)\mid g(x)\) em A[x].
    \end{lemma*}
    \begin{proof*}
      Como \(f(x)\mid g(x)\) em \(K[x]\), existe \(h(x)\in K[x]\) tal que \(g(x) = f(x)h(x).\) Pelos últimos lemas,
      h é associado a um polinômio primitivo em A[x], digamos \(h_{0}\), por meio de \(h(x) = \frac{a}{b}h_{0}(x),\) em que \(a, b\in A\) e 
      \(h_{0}(x)\) é primitivo em A[x].

      Além disso, vimos também que o produto de polinômios primitivos é um polinômio primitivo, ou seja, \(f(x)h_{0}(x)\) é primitivo. Com isso,
      aplicando um deles em \(bg(x) = a(f(x)h_{0}(x)),\) temos \(b\mid a\) em A. Então, \(b/a\in A\) e, assim, \(h(x)\in A[x].\)
      Portanto, \(f(x)\mid g(x)\) em A[x]. \qedsymbol
    \end{proof*}
    Com estas ferramentas em mãos, vemos o (Teolema) Lema de Gauss:
    \begin{theorem*}[Lema de Gauss]
      Seja A um D.F.U. e \(K = \mathrm{Frac}(A).\) Se \(f(x)\in A[x]\) é primitivo, então:
      \[
        f(x) \text{ é irredutível em }A[x] \Longleftrightarrow f(x) \text{ é irredutível em }K[x].
      \]
    \end{theorem*}
    \begin{proof*}
      \(\Leftarrow ):\) Assuma que f(x) é irredutível em K[x], mas não é irredutível em A[x]. Então, \(f(x) = g(x)h(x),\) com 
      \(g(x), h(x)\in A[x]\setminus{A[x]^{*}}\). Se um dos dois tem grau nulo, por exemplo, \(\deg{g} = 0, \) então
      \(g(x) = a\in A\setminus{A^{*}},\) e, assim, \(f(x) = ah(x),\) o que é um absurdo, já que isso significaria que 
      a dividiria todos os coeficientes de f(x), contrariando a condição de f(x) ser primitivo.

      Conclui-se, assim, que \(\deg{g(x)} > 0\) e \(\deg{h(x)} > 0,\) ou seja, \(f(x) = g(x)h(x)\) é uma fatoração de 
      f(x) em K[x], um absurdo.

      \(\Rightarrow ):\) Suponha que \(f(x)\) é irredutível em \(A[x]\), mas não em \(K[x]\). Como \(K[x]^{*} = K^{*},\) temos uma
      fatoração não trivial dada por 
      \[
        f(x) = g_{1}(x)h_{1}(x),\quad g_{1}, h_{1}\in K[x],\quad \deg{g_{1}}, \deg{h_{1}} > 0.
      \]
      Agora, usamos os lemas anteriores. Podemos escrever \(g_{1}(x) = \alpha g(x)\) e \(h_{1}(x) = \beta h_{0}(x),\) em que \(\alpha , \beta \in K^{*}\)
      e \(g, h_{0}\) são polinômios primitivos de \(A[x].\) Logo, definindo \(h(x) = \alpha \beta h_{0}(x), f(x) = g(x)h(x).\) Como g(x) é primitivo e 
      \(g(x)\mid f(x)\) em \(K[x],\) e, pelo lema, \(g(x)\mid f(x)\) em A[x]. Portanto, \(h(x)\in A[x]\) e \(f(x) = g(x)h(x)\) é uma fatoração de f(x) em A[x],
      um absurdo. \qedsymbol
    \end{proof*}
    \begin{theorem*}
      Se A é um D.F.U., então \(A[x]\) também é D.F.U. Em particular, \(A[x_{1}, \dotsc , x_{n}], n\geq 1\) é D.F.U.
    \end{theorem*}
    \begin{proof*}
      A ideia desta prova é mostrar que todo irredutível é primo em A[x] e que todo elemento de A[x] possui uma fatoração, não necessariamente
      única, em elementos irredutíveis.

      Para ver que todo elemento irredutível é primo em A[x], seja \(f(x)\in A[x]\) irredutível. Caso \(\deg{f} = 0,\) segue que \(f(x) = a\in A.\) Já 
      que A é D.F.U., a é primo em A, ou seja, a é primo em A[x]. Assumma, agora, que \(\deg{f} = n > 0.\) Vamos mostrar que \(f(x)\) é primitivo. Suponha que
      \(a\coloneqq \mathrm{mdc}(a_{0}, \dotsc , a_{n}), a_{i}\neq 0\), ou seja, a é o mdc dos coeficientes não nulos de f(x) com \(a\in A\setminus{A^{*}} \) e
      \(f(x) = ag(x),\) em que \(g(x)\in A[x].\) Logo, f(x) não é irredutível, o que seria um absurdo. A única opção que resta é que \(a\in A^{*},\) o que equivale a dizer
      que a é associado com 1.

      Agora, pelo Lema de Gauss, \(f(x)\) é irredutível em \(K[x], K = \mathrm{Frac}(A).\) Como \(K[x]\) é D.I.P. e, em D.I.P., elementos irredutíveis são primos,
      conclui-se que \(f(x)\) é primo em K[x]. Caso \(f(x)\mid g(x)h(x)\) em A[x], então, \(f(x)\mid g(x)h(x)\) em K[x]. Por ser primo, isto quer dizer que ou
      \(f(x)\mid g(x)\), ou \(f(x)\mid h(x)\), ambos em K[x]. Consequentemente, ou \(f(x)\mid g(x)\) em A[x] ou \(f(x)\mid h(x)\) em A[x], isto é, f(x) é primo em A[x].

      A seguir, mostraremos a existência de fatoração em elementos irredutíveis em A[x]. Para isto, considere \(f(x)\in A[x]\) não nulo. Como \(f(x)\in K[x],\)
      podemos fatorar f(x) em K[x], pois \(K[x]\) é D.F.U. Digamos que essa fatoração é \(f(x) = a\tilde{f}_{1}\cdot \dotsc \cdot \tilde{f}_{n}\in K[x],\) sendo cada
      \(\tilde{f}_{i}\in K[x]\) irredutíveis. Sabemos, pelos lemas de antes, que existem \(\alpha_{i}\in K^{*}\) e \(f_{i}\in A[x]\) primitivos tais que \(\tilde{f}_{i} =
      \alpha_{i}f_{i}(x).\) Assim, \(f(x) = \alpha_{1}\cdot \dotsc \cdot \alpha_{n}f_{1}\cdot \dotsc \cdot f_{n}\). Além disso, pelo Lema de Gauss, \(f_{i}\) são irredutíveis
      em A[x]. Como o produto finito de primitivos é primitivo, segue que \(f_{1}\cdot \dotsc \cdot f_{n}\) é primitivo. Destarte, \(f_{1}\cdot \dotsc \cdot f_{r}\mid f(x)\)
      em K[x] implica que \(f_{1}\cdot \dotsc \cdot f_{r}\mid f(x)\) também em A[x]. Considere \(\alpha = \alpha _{1}\cdot \dotsc \cdot \alpha_{n}\in A.\) Visto que A é 
      D.F.U., \(\alpha = a_{1}\cdot \dotsc \cdot a_{s}, a_{i}\) irredutível em A. Logo, 
      \[
        f(x) = a_{1}\cdot \dotsc a_{s}\cdot f_{1}(x)\cdot \dotsc \cdot f_{n}(x)
      \]
      é uma fatoração de f(x) em \(A[x]\) em elementos irredutíveis. Portanto, \(A[x]\) é D.F.U. \qedsymbol
    \end{proof*}
    \begin{theorem*}[Critério de Eisenstein]
      Seja A um D.F.U. e \(f(x) = a_{n}x^{n} + \dotsc + a_{0}, f(x)\in A[x]\setminus{A}.\) Suponha que exista um primo \(\pi \in A\) tal que 
      \(\pi \) não divide \(a_{n}, \pi \mid a_{i}\) para todo \(i=0, \dotsc , n-1\) e \(\pi^{2}\) não divide \(a_{0}\). Então, f(x) é irredutível em 
      A[x].
    \end{theorem*}
    \begin{proof*}
      Podemos assumir que f(x) é primitivo. Caso contraŕio, se \(d = \mathrm{mdc}(a_{1}, \dotsc ,a_{n}),\) então
      \(\frac{1}{d}f(x) = \frac{a_{n}}{d}x^{n} + \dotsc + \frac{a_{0}}{d}\in A[x]\) satisfaz as hipóteses do teorema (\textbf{Exercício}).

      Seja, então, f(x) primitivo e suponha que f não seja irredutível em \(A[x]\setminus{\{0\}}\) e coloque
      \(f(x) = g(x)h(x), g, h\in A[x], g = b_{r}x^{r} + \dotsc  + b_{0}\) e \(h = c_{s}x^{s} + \dotsc +c_{0},\) em que
      \(g,h\not\in A^{*}\) e \(r+s = n.\) Note que nenhum dos dois pode ter grau 0, visto que, caso um deles tenha, digamos que seja g ( \(\deg{g} = 0\) ), então
      \(g(x)=a\in A\setminus{A^{*}}\) e \(f(x) = ah(x),\) o que é um absurdo com o fato de f ser primitivo, já que isso significaria que a dividiria todos os coeficientes
      de f. Com isso, \(\deg{g}, \deg{h} > 0.\)

      Agora, considere \(\mathfrak{p} = \langle \pi  \rangle.\) Como \(\mathfrak{p}\) é ideal primo, sabemos que 
      \((A/\mathfrak{p})[x]\) é domínio. Além disso, por \(\pi \) não dividir \(a_{n}\) e \(\pi \mid a_{i}\) para todo \(i=0, \dotsc , n-1,\) então:
      \[
        0\neq \overline{f(x)} = \overline{a_{n}x^{n}} = \overline{a_{n}}x^{n} = \overline{g(x)h(x)}
      \]
      em \(A/\mathfrak{p}[x].\) Suponha que \(s'\) seja o menor grau com coeficiente não nulo em \(\overline{h(x)}\) e r' o análogo em \(\overline{g}.\) 
      Com essas condições, \(\overline{g(x)} = \overline{b_{r}}x^{r} + \dotsc + \overline{b_{r'}x^{r'}} + \overline{0}x^{r'-1} + \dotsc + \overline{0}\) e
      \(\overline{h(x)} = \overline{c_{s}}x^{s} + \dotsc + \overline{c_{s'}x^{s'}}.\) Logo, 
      \[
        \overline{g(x)h(x)} = \overline{c_{s}b_{r}}x^{n} + \dotsc + \overline{b_{r'}c_{s'}}x^{r'+s'} = \overline{f(x)} = \overline{a_{n}}x^{n}.
      \]
      Em seguida, mostremos que \(\overline{b_{0}} = \overline{c_{0}} = \overline{0}.\) Assuma que \(\overline{c_{0}}\neq \overline{0},\) tal que \(\overline{b_{0}}=\overline{0}\), pois
      \(\overline{b_{0}c_{0}} = \overline{0}\) e \((A/\mathfrak{p})[x]\) é domínio. No entanto, observe que \(\overline{0} = \overline{a_{1}} = \overline{b_{0}c_{1}} + \overline{b_{1}c_{0}} = \overline{b_{1}c_{0}},\)
      donde segue que \(\overline{b_{1}} = \overline{0}.\) Analogamente, para \(\overline{0} = \overline{a_{2}} = \overline{b_{0}c_{2}}+\overline{b_{1}c_{1}}+\overline{b_{2}c_{0}} = \overline{b_{2}c_{0}}\),
      de modo que \(\overline{b_{2}} = \overline{0}.\) Assim, em r passos, chegarmoes que \(\overline{b_{0}} = \dotsc = \overline{b_{r}} = \overline{r} = \overline{0},\) uma contradição.
      Logo, \(\overline{c_{0}} = \overline{0}\) e segue que \(\pi \) divide \(b_{0}\) e \(c_{0}\) o que ipmlica que \(\pi^{2}\) divide \(b_{0}c_{0} = a_{0}.\) Novamente, 
      contradição. Portanto, f(x) é irredutível e o resultado segue. \qedsymbol
    \end{proof*}
    \newpage

    \end{document}
